% Amélioration de la conservation des fruits de saison — SVT Tle D — Cours signé OnBuch+
% Programme officiel MINESEC. Compiler avec : tectonic N17-conservation-fruits-de-saison.tex
\newcommand{\DOCMATIERE}{SVT}
\newcommand{\DOCNIVEAU}{Tle D}
\newcommand{\DOCMODULE}{Module IV : La Biotechnologie — Transformation et conservation des fruits de saison}
\newcommand{\DOCLECON}{N17}
\newcommand{\DOCTITRE}{Amélioration de la conservation des fruits de saison}
% OnBuch+ — préambule « cours signés » ENRICHI (compilé avec tectonic / XeLaTeX)
% Système visuel « Pop » d'OnBuch+ : fond crème (jamais blanc), bordures encre
% épaisses, ombres « dures » décalées, titres Archivo Black en capitales.
% Version MATHÉMATIQUES : graphes (pgfplots/tikz), boîtes pédagogiques
% (définition, propriété, méthode, attention, le savais-tu, exercice résolu,
% exercice, TP, bilan), page de garde et signature OnBuch+ ancrée en bas.
%
% Macros attendues dans main.tex AVANT \input{preamble} :
%   \DOCMATIERE  \DOCNIVEAU  \DOCMODULE  \DOCLECON (numéro)  \DOCTITRE
\documentclass[11pt]{article}

\usepackage{fontspec}
\usepackage[french]{babel}
\usepackage[a4paper,margin=2cm,top=3.4cm,bottom=2.8cm,headheight=1.4cm,footskip=1.3cm]{geometry}
\usepackage{amsmath,amssymb,amsfonts}
\usepackage{mathtools} % \xmapsto, \coloneqq... souvent utilisés par les modèles
\usepackage{cancel} % simplifications barrées (bilans de forces)
% \abs{z} (module, valeur absolue) et \norm{u} (norme), utilisés par les modèles
\providecommand{\abs}{}\let\abs\relax\DeclarePairedDelimiter{\abs}{\lvert}{\rvert}
\providecommand{\norm}{}\let\norm\relax\DeclarePairedDelimiter{\norm}{\lVert}{\rVert}
\usepackage[table]{xcolor} % [table] : fournit aussi \rowcolors (alternance de lignes)
\usepackage{pagecolor}
\usepackage{tikz}
\usetikzlibrary{arrows.meta,positioning,calc,shapes.geometric,shapes.misc,shapes.arrows,decorations.pathmorphing,decorations.pathreplacing,decorations.markings,patterns,shadows,mindmap,trees,fit,backgrounds,matrix,intersections,angles,quotes}
\usepackage{pgfplots}
\usepackage[version=4]{mhchem} % équations nucléaires \ce{^{235}_{92}U}
\usepackage[european,siunitx]{circuitikz} % schémas électriques
\pgfplotsset{compat=1.18}
\usepgfplotslibrary{fillbetween,groupplots} % aires entre courbes (calcul intégral)
\usepackage{enumitem}
\usepackage{fancyhdr}
% [most] charge toutes les bibliothèques tcolorbox (listings, minted, raster,
% documentation...) et gonfle inutilement la mémoire TeX sur un document long
% (~300 boîtes) : on ne charge que ce qui est réellement utilisé.
\usepackage[skins,breakable]{tcolorbox}
\usepackage{graphicx}
\usepackage{colortbl}
\usepackage{booktabs}
\usepackage{tabularx}
\usepackage{diagbox} % case d'en-tête diagonale des tableaux à double entrée
% Colonnes extensibles centrée / gauche / droite (C, L, R), souvent utilisées par les modèles
\newcolumntype{C}{>{\centering\arraybackslash}X}
\newcolumntype{L}{>{\raggedright\arraybackslash}X}
\newcolumntype{R}{>{\raggedleft\arraybackslash}X}
\usepackage{array}
\usepackage{multicol}
\usepackage{siunitx}
% parse-numbers=false : accepte aussi \SI{2,5 \times 10^{-3}}{...}
\sisetup{locale=FR,output-decimal-marker={,},group-minimum-digits=4,parse-numbers=false}
\DeclareSIUnit\ohm{\ensuremath{\symup{\Omega}}} % U+2126 absent de la police
\DeclareSIUnit\division{div} % réglages d'oscilloscope (ms/div, V/div)
\DeclareSIUnit\tr{tr} % tours (vitesse de rotation, ex. \SI{1500}{\tr\per\minute})
\DeclareSIUnit\molar{M} % concentration molaire (ex. \SI{3}{\molar}, \SI{1}{\micro\molar})
\DeclareSIUnit\an{an} % année (le modèle confond parfois avec \year, un compteur TeX) — ex. \SI{5}{\kilo\meter\per\an}
\DeclareSIUnit\FCFA{FCFA} % franc CFA (ex. \SI{500}{\FCFA\per\kilo\gram})
\DeclareSIUnit\degreeBrix{\textdegree Brix} % teneur en sucre d'un sirop/jus (réfractométrie)
\DeclareSIUnit\month{mois} % le modèle utilise parfois \month, un compteur TeX, comme unité de durée
\usepackage{qrcode}
% \degree : commande fréquemment utilisée par les modèles pour un angle ou une
% température en dehors de \SI{}{} (non fournie par les paquets chargés).
\providecommand{\degree}{^{\circ}}
% \qed (fin de démonstration), utilisé par les modèles sans amsthm
\providecommand{\qed}{\hfill$\square$}
% \barre : double barre des valeurs interdites dans un tableau de variations (tkz-tab)
\providecommand{\barre}{\ensuremath{\|}}
% environnement proof (amsthm non chargé) utilisé par les modèles
\ifdefined\proof\else\newenvironment{proof}[1][Démonstration]{\par\noindent\textit{#1.}\ }{\qed\par}\fi
\usepackage{lastpage}
\usepackage{hyperref}
\hypersetup{hidelinks}

% ============================================================
% Palette « Pop » OnBuch+ — mêmes hex que la classe P (plus_ui.dart)
% ============================================================
\definecolor{poporange}{HTML}{F59321}
\definecolor{popdark}{HTML}{E07A0C}
\definecolor{popcream}{HTML}{FFF6E8}
\definecolor{popink}{HTML}{1C1714}
\definecolor{poppurple}{HTML}{7A5AE0}
\definecolor{popgreen}{HTML}{2E9E6A}
\definecolor{poppink}{HTML}{E0457B}
\definecolor{popblue}{HTML}{2F7DE1}
\definecolor{popgold}{HTML}{C98A12}
\definecolor{popmuted}{HTML}{8A7F76}
\definecolor{popline}{HTML}{EADFD2}
\definecolor{popgray}{HTML}{8A7F76}
\colorlet{popgrey}{popgray}
% Alias tolérants (le modèle nomme parfois les couleurs différemment)
\colorlet{popwhite}{white}
\colorlet{popblack}{popink}
\colorlet{popred}{poppink}
\colorlet{popyellow}{popgold}
% Teintes claires (fonds de tableaux/encadrés) — le modèle nomme parfois
% les couleurs avec un suffixe L (ex. popblueL) sans les avoir définies.
\colorlet{poporangeL}{poporange!20}
\colorlet{popdarkL}{popdark!20}
\colorlet{poppurpleL}{poppurple!20}
\colorlet{popgreenL}{popgreen!20}
\colorlet{poppinkL}{poppink!20}
\colorlet{popblueL}{popblue!20}
\colorlet{popgoldL}{popgold!20}
\colorlet{popredL}{popred!20}
\colorlet{popgrayL}{popgray!20}
% Teintes claires pour fonds de tableaux / zones
\colorlet{popblueL}{popblue!10!white}
\colorlet{poporangeL}{poporange!14!white}
\colorlet{popgreenL}{popgreen!12!white}
\colorlet{poppurpleL}{poppurple!12!white}
\colorlet{poppinkL}{poppink!12!white}
\colorlet{popgoldL}{popgold!14!white}

\pagecolor{popcream}

% ============================================================
% Typographie
% ============================================================
\defaultfontfeatures{Ligatures=TeX}
\setmainfont{PlusJakartaSans-Medium.ttf}[Path=../../../fonts/,BoldFont=PlusJakartaSans-Bold.ttf]
% Maths en sans-serif (Fira Math) avec chiffres et lettres droites de Plus
% Jakarta, pour que les formules se fondent dans le texte.
\usepackage[math-style=ISO,bold-style=ISO]{unicode-math}
\setmathfont{FiraMath-Regular.otf}[Path=../../../fonts/]
\setmathfont{PlusJakartaSans-Medium.ttf}[Path=../../../fonts/,range={up/num,up/{Latin,latin},\mathup}]
\usepackage{icomma} % virgule décimale sans espace en mode math
% Caractères mathématiques Unicode absents des polices de texte (Plus Jakarta,
% Archivo Black) : rendus par leur équivalent mathématique, en texte comme en
% formule — sinon ils s'affichent en carré vide (ex. « Arithmétique dans ℤ »).
\usepackage{newunicodechar}
\newunicodechar{ℝ}{\ensuremath{\mathbb{R}}}
\newunicodechar{ℤ}{\ensuremath{\mathbb{Z}}}
\newunicodechar{ℂ}{\ensuremath{\mathbb{C}}}
\newunicodechar{ℕ}{\ensuremath{\mathbb{N}}}
\newunicodechar{ℚ}{\ensuremath{\mathbb{Q}}}
\newunicodechar{𝔻}{\ensuremath{\mathbb{D}}}
\newunicodechar{α}{\ensuremath{\alpha}}
\newunicodechar{β}{\ensuremath{\beta}}
\newunicodechar{γ}{\ensuremath{\gamma}}
\newunicodechar{δ}{\ensuremath{\delta}}
\newunicodechar{ε}{\ensuremath{\varepsilon}}
\newunicodechar{θ}{\ensuremath{\theta}}
\newunicodechar{λ}{\ensuremath{\lambda}}
\newunicodechar{μ}{\ensuremath{\mu}}
\newunicodechar{σ}{\ensuremath{\sigma}}
\newunicodechar{τ}{\ensuremath{\tau}}
\newunicodechar{φ}{\ensuremath{\varphi}}
\newunicodechar{ψ}{\ensuremath{\psi}}
\newunicodechar{ω}{\ensuremath{\omega}}
\newunicodechar{Σ}{\ensuremath{\Sigma}}
% majuscules produites par \MakeUppercase dans les titres (α-aminés -> Α)
\newunicodechar{Α}{\ensuremath{\alpha}}
\newunicodechar{Β}{\ensuremath{\beta}}
\newunicodechar{Γ}{\ensuremath{\Gamma}}
\newunicodechar{Δ}{\ensuremath{\Delta}}
\newunicodechar{Ε}{\ensuremath{\varepsilon}}
\newunicodechar{Θ}{\ensuremath{\Theta}}
\newunicodechar{Λ}{\ensuremath{\Lambda}}
\newunicodechar{Μ}{\ensuremath{\mu}}
\newunicodechar{Τ}{\ensuremath{\tau}}
\newunicodechar{Φ}{\ensuremath{\Phi}}
\newunicodechar{Ψ}{\ensuremath{\Psi}}
\newunicodechar{Ω}{\ensuremath{\Omega}}
\newunicodechar{↦}{\ensuremath{\mapsto}}
\newunicodechar{∈}{\ensuremath{\in}}
\newunicodechar{∉}{\ensuremath{\notin}}
\newunicodechar{∧}{\ensuremath{\wedge}}
\newunicodechar{⇔}{\ensuremath{\Leftrightarrow}}
\newunicodechar{⟺}{\ensuremath{\Longleftrightarrow}}
\newunicodechar{⇒}{\ensuremath{\Rightarrow}}
\newunicodechar{⟹}{\ensuremath{\Longrightarrow}}
\newunicodechar{∘}{\ensuremath{\circ}}
\newunicodechar{∪}{\ensuremath{\cup}}
\newunicodechar{∩}{\ensuremath{\cap}}
\newunicodechar{≡}{\ensuremath{\equiv}}
\newunicodechar{⊂}{\ensuremath{\subset}}
\newunicodechar{⊥}{\ensuremath{\perp}}
\newunicodechar{∃}{\ensuremath{\exists}}
\newunicodechar{∀}{\ensuremath{\forall}}
\newunicodechar{∛}{\ensuremath{\sqrt[3]{\,}}}
\newunicodechar{∜}{\ensuremath{\sqrt[4]{\,}}}
\newunicodechar{⃗}{\ensuremath{{}^{\to}}}
\newunicodechar{ⁿ}{\ifmmode^{n}\else\textsuperscript{\ensuremath{n}}\fi}
\newunicodechar{ˣ}{\ifmmode^{x}\else\textsuperscript{\ensuremath{x}}\fi}
\newunicodechar{ᵇ}{\ifmmode^{b}\else\textsuperscript{\ensuremath{b}}\fi}
\newunicodechar{ʳ}{\ifmmode^{r}\else\textsuperscript{\ensuremath{r}}\fi}
\newunicodechar{ᵘ}{\ifmmode^{u}\else\textsuperscript{\ensuremath{u}}\fi}
\newunicodechar{ʸ}{\ifmmode^{y}\else\textsuperscript{\ensuremath{y}}\fi}
\newunicodechar{ᵃ}{\ifmmode^{a}\else\textsuperscript{\ensuremath{a}}\fi}
\newunicodechar{ᵉ}{\ifmmode^{e}\else\textsuperscript{\ensuremath{e}}\fi}
\newunicodechar{ᵏ}{\ifmmode^{k}\else\textsuperscript{\ensuremath{k}}\fi}
\newunicodechar{ᵖ}{\ifmmode^{p}\else\textsuperscript{\ensuremath{p}}\fi}
\newunicodechar{ᴺ}{\ifmmode^{N}\else\textsuperscript{\ensuremath{N}}\fi}
\newunicodechar{ᶜ}{\ifmmode^{c}\else\textsuperscript{\ensuremath{c}}\fi}
\newunicodechar{ʰ}{\ifmmode^{h}\else\textsuperscript{\ensuremath{h}}\fi}
\newunicodechar{ᵠ}{\ifmmode^{q}\else\textsuperscript{\ensuremath{q}}\fi}
\newunicodechar{ₙ}{\ifmmode_{n}\else\textsubscript{\ensuremath{n}}\fi}
\newunicodechar{ₐ}{\ifmmode_{a}\else\textsubscript{\ensuremath{a}}\fi}
\newunicodechar{ᵢ}{\ifmmode_{i}\else\textsubscript{\ensuremath{i}}\fi}
\newunicodechar{ᵣ}{\ifmmode_{r}\else\textsubscript{\ensuremath{r}}\fi}
\newunicodechar{ₖ}{\ifmmode_{k}\else\textsubscript{\ensuremath{k}}\fi}
\newunicodechar{ᵦ}{\ifmmode_{\beta}\else\textsubscript{\ensuremath{\beta}}\fi}
\newunicodechar{ₓ}{\ifmmode_{x}\else\textsubscript{\ensuremath{x}}\fi}
\newfontfamily\popdisplay{ArchivoBlack-Regular.ttf}[Path=../../../fonts/]
\newfontfamily\poplogo{SpaceGrotesk-Bold.ttf}[Path=../../../fonts/]
\newfontfamily\popsemi{PlusJakartaSans-SemiBold.ttf}[Path=../../../fonts/]
\newfontfamily\popxbold{PlusJakartaSans-ExtraBold.ttf}[Path=../../../fonts/]
\newfontfamily\popxboldls{PlusJakartaSans-ExtraBold.ttf}[Path=../../../fonts/,LetterSpace=3.0]

\setlist[enumerate]{leftmargin=1.7em,itemsep=5pt,topsep=4pt}
\setlist[itemize]{leftmargin=1.7em,itemsep=4pt,topsep=4pt,label={\color{poporange}\textbullet}}
\setlength{\parindent}{0pt}
\setlength{\parskip}{5pt}
\renewcommand{\arraystretch}{1.25}

% pgfplots : style Pop par défaut
\pgfplotsset{
  every axis/.append style={axis line style={popink,line width=1pt},tick style={popink},
    grid style={popline,line width=0.5pt},label style={font=\small},tick label style={font=\footnotesize}},
  popcurve/.style={poporange,line width=1.8pt,no markers,smooth},
}
% Même style disponible dans un \draw TikZ simple (hors pgfplots)
\tikzset{popcurve/.style={poporange,line width=1.8pt}}
\tikzset{popgrid/.style={popline,very thin}}
\colorlet{popcurve}{poporange} % le nom du style est parfois utilisé comme couleur

% ============================================================
% Pill / squiggle
% ============================================================
\newcommand{\poppill}[3]{%
  \tikz[baseline=-0.55ex]\node[draw=popink,line width=1.2pt,rounded corners=2.6pt,%
    fill=#2,inner xsep=6.5pt,inner ysep=3pt]{%
    \popxboldls\fontsize{7}{7}\selectfont\color{#3}\MakeUppercase{#1}};%
}
\newcommand{\popsquiggle}[1]{%
  \tikz[baseline=-0.4ex]{
    \draw[#1,line width=2.6pt,line cap=round]
      plot [smooth] coordinates {(0,0.09) (0.34,0.01) (0.68,0.21) (1.02,0.01) (1.36,0.10)};
  }%
}
% Mot-clé surligné (marqueur orange sous le texte)
\newcommand{\cle}[1]{{\popxbold\color{popdark}#1}}

% ============================================================
% En-tête / pied
% ============================================================
\pagestyle{fancy}
\fancyhf{}
\renewcommand{\headrulewidth}{0pt}
\renewcommand{\footrulewidth}{0pt}
\lhead{\raisebox{-0.15ex}{\poplogo\fontsize{13}{13}\selectfont\color{popink}OnBuch\color{poporange}+}}
\rhead{\poppill{\DOCMATIERE\ \textbullet\ \DOCNIVEAU\ \textbullet\ Leçon \DOCLECON}{white}{popink}}
\lfoot{\popxbold\fontsize{7.6}{7.6}\selectfont\color{popmuted}\MakeUppercase{Cours signé OnBuch+}\ \textbullet\ {\popsemi\DOCTITRE}}
\rfoot{\tikz[baseline=-0.6ex]\node[rounded corners=8.5pt,fill=popink,minimum height=17pt,minimum width=17pt,inner xsep=5pt,inner sep=0]{\popxbold\fontsize{8}{8}\selectfont\color{popcream}\thepage\,/\,\pageref*{LastPage}};}

% ============================================================
% Titres
% ============================================================
\newcommand{\chaptitre}[2]{%
  \begin{tcolorbox}[enhanced,colback=poporange,colframe=popink,boxrule=2.6pt,arc=11pt,
      drop shadow=popink,left=15pt,right=15pt,top=12pt,bottom=11pt,before skip=3pt,after skip=18pt]
    \poppill{#2}{popink}{popcream}\par\vspace{9pt}
    {\raggedright\hyphenpenalty=10000\exhyphenpenalty=10000\popdisplay\fontsize{22}{24}\selectfont\color{popcream}\MakeUppercase{#1}\par}
  \end{tcolorbox}
}
% Section numérotée : pastille orange avec le numéro + titre Archivo
\newcounter{popsec}
\newcommand{\coursec}[1]{%
  \stepcounter{popsec}\setcounter{popsub}{0}%
  \par\vspace{16pt}\noindent
  \begin{minipage}{\linewidth}
  \tikz[baseline=-0.7ex]\node[draw=popink,line width=1.6pt,fill=poporange,rounded corners=4pt,minimum size=22pt,inner sep=0]{\popdisplay\fontsize{12}{12}\selectfont\color{popcream}\thepopsec};%
  \hspace{8pt}{\popdisplay\fontsize{15}{17}\selectfont\color{popink}\MakeUppercase{#1}}\par
  \vspace{4pt}\noindent{\color{poporange}\rule{36pt}{3.6pt}}
  \end{minipage}\par\nopagebreak\vspace{8pt}
}
\newcounter{popsub}
\newcommand{\courssub}[1]{%
  \stepcounter{popsub}%
  \par\vspace{9pt}\noindent{\popxbold\fontsize{11.5}{13}\selectfont\color{popink}{\color{poporange}\thepopsec.\thepopsub}\hspace{6pt}#1}\par\nopagebreak\vspace{4pt}
}

% ============================================================
% Boîtes pédagogiques — même recette : fond blanc, bordure épaisse colorée,
% ombre dure encre, badge plein. Titre optionnel [#1].
% ============================================================
\tcbset{popbase/.style={breakable,enhanced,colback=white,boxrule=2.3pt,arc=9pt,
  shadow={3pt}{-3pt}{0pt}{popink},left=14pt,right=14pt,top=11pt,bottom=11pt,before skip=11pt,after skip=15pt,
  fontupper=\popsemi\fontsize{10.6}{14.5}\selectfont\color{popink},
  fontlower=\popsemi\fontsize{10.6}{14.5}\selectfont\color{popink}}}
% Coche dessinée (le glyphe ✓ n'existe pas dans les polices du document)
\newcommand{\popcheck}[1]{\tikz[baseline=-0.5ex]\draw[#1,line width=1.6pt,line cap=round,line join=round](-0.13,0.0)--(-0.04,-0.09)--(0.13,0.1);}
\providecommand{\coche}{\popcheck{popgreen}}
\providecommand{\resultat}[1]{\par\smallskip\noindent\fcolorbox{popgreen}{popgreenL}{\parbox{\dimexpr\linewidth-2\fboxsep-2\fboxrule\relax}{\textbf{Résultat.} #1}}\par\smallskip}
\newcommand{\popboxhead}[3]{\poppill{#1}{#2}{white}\ifx\relax#3\relax\else\hspace{6pt}{\popxbold\fontsize{10.6}{12}\selectfont\color{popink}#3}\fi\par\vspace{7pt}}

\newtcolorbox{aretenir}[1][]{popbase,colframe=popgreen,before upper={\popboxhead{À retenir}{popgreen}{#1}}}
\newtcolorbox{exemplebox}[1][]{popbase,colframe=poporange,before upper={\popboxhead{Exemple}{poporange}{#1}}}
\newtcolorbox{definition}[1][]{popbase,colframe=popblue,before upper={\popboxhead{Définition}{popblue}{#1}}}
\newtcolorbox{propriete}[1][]{popbase,colframe=poppurple,before upper={\popboxhead{Propriété}{poppurple}{#1}}}
\newtcolorbox{methode}[1][]{popbase,colframe=popgold,before upper={\popboxhead{Méthode}{popgold}{#1}}}
\newtcolorbox{attention}[1][]{popbase,colframe=poppink,before upper={\popboxhead{Attention}{poppink}{#1}}}
\newtcolorbox{experience}[1][]{popbase,colframe=popblue,colback=popblueL,before upper={\popboxhead{Expérience}{popblue}{#1}}}
% « Le savais-tu ? » : carte violette pleine (contexte, culture, Cameroun)
\newtcolorbox{savaistu}[1][]{popbase,colback=poppurple,colframe=popink,
  fontupper=\popsemi\fontsize{10.4}{14.2}\selectfont\color{white},
  before upper={\poppill{Le savais-tu\,?}{white}{poppurple}\ifx\relax#1\relax\else\hspace{6pt}{\popxbold\color{white}#1}\fi\par\vspace{7pt}}}
% Exercice résolu : énoncé en haut, \tcblower puis la solution sur fond crème
\newcounter{popexo}\newcounter{popexr}
\newtcolorbox{exoresolu}[1][]{popbase,colframe=poporange,colbacklower=popcream,
  segmentation style={popink,line width=1.2pt,dashed},
  before upper={\stepcounter{popexr}\popboxhead{Exercice résolu \thepopexr}{poporange}{#1}},
  before lower={\poppill{Solution}{popink}{popcream}\par\vspace{6pt}}}
% Exercice (non résolu en place) — la correction est dans la partie Corrigés
\newtcolorbox{exercice}[1][]{popbase,colframe=popink,
  before upper={\stepcounter{popexo}\popboxhead{Exercice \thepopexo}{popink}{#1}}}
% Corrigé
\newtcolorbox{corrige}[1][]{popbase,colframe=popgreen,colback=popgreenL,
  before upper={\popboxhead{Corrigé}{popgreen}{#1}}}
% Objectifs / prérequis / bilan
\newtcolorbox{objectifs}{popbase,colframe=popink,colback=poporangeL,
  before upper={\popboxhead{Objectifs de la leçon}{popink}{}}}
\newtcolorbox{prerequis}{popbase,colframe=popink,colback=popblueL,
  before upper={\popboxhead{Prérequis}{popblue}{}}}
\newtcolorbox{bilan}{popbase,colframe=popink,colback=white,boxrule=2.8pt,
  before upper={\popboxhead{Fiche bilan}{poporange}{L'essentiel à connaître pour l'examen}}}
% Figure encadrée avec légende
\newtcolorbox{popfigure}[1][]{enhanced,breakable=false,colback=white,colframe=popline,boxrule=1.4pt,arc=8pt,
  left=10pt,right=10pt,top=10pt,bottom=8pt,before skip=10pt,after skip=12pt,halign=center,
  lower separated=false,halign lower=center,
  fontlower=\popsemi\fontsize{9}{11.5}\selectfont\color{popmuted},
  #1}
\newcommand{\legende}[1]{\par\vspace{6pt}{\popsemi\fontsize{9}{11.5}\selectfont\color{popmuted}#1}}

% ============================================================
% Signature OnBuch+
% ============================================================
\newcommand{\poplogoinline}[1]{{\poplogo\fontsize{#1}{#1}\selectfont\color{popink}OnBuch\color{poporange}+}}
\newcommand{\signatureonbuch}{%
  \begin{tcolorbox}[enhanced,colback=popink,colframe=popink,boxrule=2.2pt,arc=10pt,
      drop shadow=poporange,left=14pt,right=14pt,top=10pt,bottom=10pt,before skip=0pt,after skip=0pt]
    \tikz[baseline=-0.7ex]\node[circle,fill=popgreen,draw=popcream,line width=1.3pt,minimum size=22pt,inner sep=0]{\popcheck{white}};%
    \hspace{9pt}\begin{minipage}[c]{0.62\linewidth}
      {\popxbold\fontsize{12}{13}\selectfont\color{popcream}Signé\ \textbullet\ {\poplogo OnBuch\color{poporange}+}}\par\vspace{2pt}
      {\raggedright\popsemi\fontsize{8}{10.5}\selectfont\color{popline}\MakeUppercase{Équipe pédagogique OnBuch+}\\\MakeUppercase{Conforme au programme officiel MINESEC}\par}
    \end{minipage}\hfill
    {\poplogo\fontsize{22}{22}\selectfont\color{popcream}OnBuch\color{poporange}+}
  \end{tcolorbox}%
}

% ============================================================
% Page de garde
% #1 titre · #2 matière · #3 niveau · #4 module · #5 n° leçon · #6 durée ·
% #7 description / objectifs (texte ou itemize)
% ============================================================
\newcommand{\pagegarde}[7]{%
  \thispagestyle{empty}
  \newgeometry{a4paper,margin=1.7cm,top=1.6cm,bottom=1.4cm}%
  \begin{tikzpicture}[remember picture,overlay]
    \fill[poporange!18] ([xshift=-3.2cm,yshift=-3.6cm]current page.north east) circle (4.6cm);
    \fill[poppurple!14] ([xshift=2.4cm,yshift=7.6cm]current page.south west) circle (3.8cm);
  \end{tikzpicture}%
  {\poplogo\fontsize{30}{30}\selectfont\color{popink}OnBuch\color{poporange}+}\hfill
  \raisebox{0.6ex}{\poppill{Cours signé \textbullet\ Premium}{popink}{popcream}}\par
  \vspace{1pt}\popsquiggle{poppurple}\par
  \vspace{8pt}
  \begin{tcolorbox}[enhanced,colback=poporange,colframe=popink,boxrule=2.8pt,arc=13pt,
      drop shadow=popink,left=18pt,right=18pt,top=12pt,bottom=13pt,after skip=10pt]
    \poppill{#2 \textbullet\ #3}{popink}{popcream}\hspace{5pt}\poppill{#4}{white}{popink}\par\vspace{8pt}
    {\popdisplay\fontsize{14}{14}\selectfont\color{popink}LEÇON #5}\par\vspace{4pt}
    {\raggedright\hyphenpenalty=10000\exhyphenpenalty=10000\popdisplay\fontsize{25}{27}\selectfont\color{popcream}\MakeUppercase{#1}\par}
    \vspace{8pt}
    {\popxbold\fontsize{9.5}{9.5}\selectfont\color{popink}\MakeUppercase{Durée conseillée : #6}}
  \end{tcolorbox}
  \begin{tcolorbox}[enhanced,colback=white,colframe=popink,boxrule=2.2pt,arc=10pt,
      drop shadow=popink,left=16pt,right=16pt,top=10pt,bottom=10pt,after skip=8pt]
    \poppill{Dans cette leçon}{poppurple}{white}\par\vspace{5pt}
    \popsemi\fontsize{9.3}{12.6}\selectfont\color{popink}#7
  \end{tcolorbox}
  \noindent\begin{minipage}[t]{0.487\linewidth}
    \begin{tcolorbox}[enhanced,colback=popgreen,colframe=popink,boxrule=2.2pt,arc=10pt,
        drop shadow=popink,left=11pt,right=11pt,top=10pt,bottom=10pt,height=3.1cm,valign=center]
      \tcbox[colback=white,colframe=popink,boxrule=1.3pt,arc=5pt,boxsep=0pt,nobeforeafter,
        left=3pt,right=3pt,top=3pt,bottom=3pt,valign=center]{\qrcode[height=1.9cm]{https://play.google.com/store/apps/details?id=com.luvvix.onbuch&hl=fr}}%
      \hspace{8pt}\begin{minipage}[c]{0.5\linewidth}
        {\popdisplay\fontsize{11}{12.5}\selectfont\color{white}\MakeUppercase{Télécharge OnBuch}}\par\vspace{4pt}
        {\raggedright\popsemi\fontsize{8.4}{11}\selectfont\color{white}Gratuit \textbullet\ Google Play\\Tuteur IA, annales, concours, résultats\par}
      \end{minipage}
    \end{tcolorbox}
  \end{minipage}\hfill
  \begin{minipage}[t]{0.487\linewidth}
    \begin{tcolorbox}[enhanced,colback=poppurple,colframe=popink,boxrule=2.2pt,arc=10pt,
        drop shadow=popink,left=11pt,right=11pt,top=10pt,bottom=10pt,height=3.1cm,valign=center]
      \tikz[baseline=-0.7ex]\node[circle,fill=white,draw=popink,line width=1.3pt,minimum size=1.6cm,inner sep=0]{\popdisplay\fontsize{24}{24}\selectfont\color{poppurple}+};%
      \hspace{8pt}\begin{minipage}[c]{0.6\linewidth}
        {\popdisplay\fontsize{11}{12.5}\selectfont\color{white}\MakeUppercase{Passe à OnBuch+}}\par\vspace{4pt}
        {\raggedright\popsemi\fontsize{8.4}{11}\selectfont\color{white}Tous les cours signés, Tuteur IA illimité, fiches PDF — onglet OnBuch+ de l'app\par}
      \end{minipage}
    \end{tcolorbox}
  \end{minipage}
  \vspace{8pt}
  \popcontenu
  \vfill
  \signatureonbuch
  \restoregeometry
  \newpage
}

% Rangée « Ce cours contient » (page de garde)
\newcommand{\poptile}[3]{% #1 couleur #2 gros texte #3 légende
  \begin{tcolorbox}[enhanced,colback=#1,colframe=popink,boxrule=2pt,arc=9pt,shadow={2.5pt}{-2.5pt}{0pt}{popink},
      left=6pt,right=6pt,top=9pt,bottom=9pt,halign=center,valign=center,height=1.75cm,width=\linewidth,nobeforeafter]
    {\popdisplay\fontsize{12.5}{13}\selectfont\color{white}\MakeUppercase{#2}}\par\vspace{3pt}
    {\centering\popsemi\fontsize{7.8}{9.5}\selectfont\color{white}#3\par}
  \end{tcolorbox}}
\newcommand{\popcontenu}{%
  \par\noindent{\popxboldls\fontsize{7.5}{7.5}\selectfont\color{popmuted}CE COURS SIGNÉ CONTIENT}\par\vspace{7pt}
  \noindent\begin{minipage}[t]{0.235\linewidth}\poptile{popblue}{Cours}{complet, illustré, pas à pas}\end{minipage}\hfill
  \begin{minipage}[t]{0.235\linewidth}\poptile{poporange}{Méthodes}{et exercices résolus}\end{minipage}\hfill
  \begin{minipage}[t]{0.235\linewidth}\poptile{poppink}{Exercices}{type examen + corrigés}\end{minipage}\hfill
  \begin{minipage}[t]{0.235\linewidth}\poptile{popgreen}{Fiche bilan}{l'essentiel à retenir}\end{minipage}\par}

% Sommaire
\newtcolorbox{sommaire}{enhanced,colback=white,colframe=popink,boxrule=2.2pt,arc=10pt,
  drop shadow=popink,left=18pt,right=18pt,top=13pt,bottom=13pt,before skip=6pt,after skip=20pt,
  before upper={\poppill{Sommaire}{poppurple}{white}\par\vspace{9pt}\popsemi\fontsize{10.6}{14.5}\selectfont\color{popink}}}

% Signature de fin de cours — poussée tout en bas de la dernière page
\newcommand{\findecours}{%
  \par\vfill\nopagebreak
  \begin{center}\popsquiggle{poporange}\end{center}\vspace{4pt}
  \signatureonbuch
}
\AtBeginDocument{\def\llbracket{\mathopen{[\mkern-3.5mu[}}\def\rrbracket{\mathclose{]\mkern-3.5mu]}}} % intervalles d'entiers (glyphe ⟦ absent de la police)
% Glyphes absents de Fira Math : redéfinis à partir de glyphes disponibles
\AtBeginDocument{%
  \def\setminus{\mathbin{\backslash}}\let\smallsetminus\setminus
  \def\vdots{\mathord{\vbox{\baselineskip3.2pt\lineskiplimit0pt\kern4pt\hbox{.}\hbox{.}\hbox{.}}}}%
  \def\ddots{\mathinner{\mkern1mu\raise6.5pt\hbox{.}\mkern2mu\raise3.6pt\hbox{.}\mkern2mu\raise0.7pt\hbox{.}\mkern1mu}}%
  \def\triangle{\mathord{\scalebox{0.9}{$\symup{\Delta}$}}}%
  \def\nabla{\mathord{\raisebox{\depth}{\scalebox{1}[-1]{$\symup{\Delta}$}}}}%
  \def\checkmark{\popcheck{popdark}}%
  \def\star{\mathbin{\ast}}\def\bigstar{\mathord{\ast}}%
  \def\diamond{\mathbin{\scalebox{0.6}{$\Diamond$}}}%
  \def\bigcup{\mathop{\mathchoice{\raisebox{-0.3em}{\scalebox{1.7}{$\cup$}}}{\scalebox{1.25}{$\cup$}}{\cup}{\cup}}\displaylimits}%
  \def\bigcap{\mathop{\mathchoice{\raisebox{-0.3em}{\scalebox{1.7}{$\cap$}}}{\scalebox{1.25}{$\cap$}}{\cap}{\cap}}\displaylimits}%
  \def\lesseqqgtr{\mathrel{\lessgtr}}\def\bigoplus{\mathop{\oplus}\displaylimits}%
}
\providecommand{\ssi}{si et seulement si} % abréviation fréquente
\AtBeginDocument{\providecommand{\wideparen}{\overparen}} % arc de cercle

%% FIGFIX-BEGIN v5 — correctifs globaux des figures (docs/onbuch-generation/tools/figfix)
\usepackage{adjustbox}\usepackage{etoolbox}\tcbuselibrary{raster}
% toute figure TikZ/pgfplots plus large que la ligne est réduite à la largeur disponible
\BeforeBeginEnvironment{tikzpicture}{\begin{adjustbox}{max width=\linewidth}}
\AfterEndEnvironment{tikzpicture}{\end{adjustbox}}
% légendes pgfplots lisibles par-dessus les courbes
\pgfplotsset{every axis/.append style={legend style={fill=white,fill opacity=0.92,text opacity=1,draw=black!15,inner sep=3pt}}}
% étiquettes d'arêtes (syntaxe edge["..."]) avec fond blanc
\tikzset{every edge quotes/.append style={fill=white,inner sep=1.2pt}}
%% FIGFIX-END
\begin{document}

\pagegarde{Amélioration de la conservation des fruits de saison}{SVT}{Tle D}{Module IV : La Biotechnologie — Transformation et conservation des fruits de saison}{N17}{4 h}{Cette leçon étudie les causes d'altération des fruits de saison (mangue, tomate) et les principes biologiques et physico-chimiques des techniques de conservation : séchage, stérilisation, confiture, jus et pâtes. Elle montre comment ces procédés, à la portée des ménages et des PME camerounaises, répondent aux besoins alimentaires, économiques et environnementaux des populations.
\vspace{4pt}
\begin{multicols}{2}\small
\begin{itemize}
\item À la fin de cette leçon, je sais identifier les facteurs responsables de l'altération des fruits (micro-organismes, enzymes, oxydation) à partir d'observations et de résultats expérimentaux.
\item À la fin de cette leçon, je sais expliquer les principes biologiques sur lesquels reposent les différents procédés de conservation (destruction ou inhibition des micro-organismes et des enzymes).
\item À la fin de cette leçon, je sais décrire les étapes des techniques de transformation et de conservation de la mangue et de la tomate (séchage, conserve, confiture, jus, pâte).
\item À la fin de cette leçon, je sais interpréter des résultats expérimentaux et des courbes (croissance microbienne, activité enzymatique en fonction de la température, du pH, de la teneur en sucre) pour justifier un procédé de conservation.
\item À la fin de cette leçon, je sais comparer l'efficacité de différentes techniques de conservation appliquées à un même fruit (durée de conservation, qualité nutritionnelle, coût).
\item À la fin de cette leçon, je sais pratiquer, en conditions simples, une technique de transformation d'un fruit de saison (confiture de mangue, séchage de rondelles).
\end{itemize}\end{multicols}}

\begin{sommaire}
\begin{multicols}{2}
\begin{enumerate}
\item I. Les fruits de saison : abondance, pertes et enjeux
\item II. Les facteurs d'altération des fruits
\item III. Les principes scientifiques de la conservation
\item IV. Les techniques de conservation et de transformation : cas de la mangue et de la tomate
\item V. Comparaison de l'efficacité des techniques de conservation
\item VI. Intérêts économiques, nutritionnels et environnementaux
\item Activité d'intégration
\item Méthodes et exercices résolus
\item Exercices
\item Corrigés des exercices
\item Fiche bilan
\end{enumerate}
\end{multicols}
\end{sommaire}

\begin{objectifs}
\begin{itemize}
\item À la fin de cette leçon, je sais identifier les facteurs responsables de l'altération des fruits (micro-organismes, enzymes, oxydation) à partir d'observations et de résultats expérimentaux.
\item À la fin de cette leçon, je sais expliquer les principes biologiques sur lesquels reposent les différents procédés de conservation (destruction ou inhibition des micro-organismes et des enzymes).
\item À la fin de cette leçon, je sais décrire les étapes des techniques de transformation et de conservation de la mangue et de la tomate (séchage, conserve, confiture, jus, pâte).
\item À la fin de cette leçon, je sais interpréter des résultats expérimentaux et des courbes (croissance microbienne, activité enzymatique en fonction de la température, du pH, de la teneur en sucre) pour justifier un procédé de conservation.
\item À la fin de cette leçon, je sais comparer l'efficacité de différentes techniques de conservation appliquées à un même fruit (durée de conservation, qualité nutritionnelle, coût).
\item À la fin de cette leçon, je sais pratiquer, en conditions simples, une technique de transformation d'un fruit de saison (confiture de mangue, séchage de rondelles).
\item À la fin de cette leçon, je sais argumenter l'intérêt économique, nutritionnel et environnemental de la conservation des fruits de saison au Cameroun.
\end{itemize}
\end{objectifs}

\begin{prerequis}
\begin{itemize}
\item Les micro-organismes (bactéries, levures, moisissures) et leurs conditions de développement (notions de Seconde/Première).
\item Les enzymes : nature, rôle catalytique, sensibilité à la température et au pH (Première).
\item Les constituants des aliments : glucides, vitamines, eau (notions de nutrition).
\item La fermentation et ses applications biotechnologiques (début du module IV).
\item Lecture et exploitation de courbes et de tableaux de résultats expérimentaux.
\end{itemize}
\end{prerequis}

\newpage

\coursec{Les fruits de saison : abondance, pertes et enjeux}

Chaque année, entre mars et juin, les marchés de Yaoundé, Douala et Ngaoundéré débordent de mangues vendues à des prix dérisoires : un seau entier s'échange parfois contre quelques centaines de francs CFA. Quelques semaines plus tard, ces mêmes mangues pourrissent en tas au bord des routes, faute d'acheteurs, dégageant une odeur de fermentation qui attire mouches et moustiques. Pourtant, en décembre, une simple mangue coûte jusqu'à dix fois plus cher, et les ménages les plus modestes doivent y renoncer. Comment expliquer qu'une telle abondance se transforme en gaspillage alors que les besoins alimentaires des populations restent immenses ? Quelles techniques, accessibles aux familles et aux petites entreprises camerounaises, permettraient de conserver ces fruits de saison et de les consommer toute l'année ? C'est à ces questions que répond la \cle{biotechnologie} de la transformation et de la conservation des fruits, objet de cette leçon.

\begin{aretenir}[Problématique de la leçon]
Un fruit récolté est un \cle{tissu vivant}, riche en eau et en nutriments : il continue de respirer, il est attaqué par des micro-organismes et dégradé par ses propres enzymes. Le problème central est donc le suivant : \textbf{comment conserver un aliment vivant et riche en eau sans qu'il se dégrade ?} Toute la leçon consistera à identifier les facteurs d'altération, puis les principes et techniques qui permettent de les bloquer.
\end{aretenir}

\courssub{La notion de fruit de saison}

Observe le calendrier des marchés camerounais : les manguiers du Nord, de l'Adamaoua et de l'Ouest fleurissent en saison sèche et donnent leurs fruits presque simultanément entre mars et juin ; les tomates, cultivées en saison sèche sous irrigation dans les bas-fonds, arrivent en abondance entre décembre et avril. Pendant quelques semaines, l'offre explose ; le reste de l'année, le fruit devient rare ou introuvable à l'état frais. C'est ce phénomène que décrit la notion de \emph{fruit de saison}.

\begin{definition}[Fruit de saison]
Un \cle{fruit de saison} est un fruit dont la production est \textbf{concentrée sur une courte période de l'année}, en raison du cycle biologique de la plante (floraison puis fructification déclenchées par des facteurs climatiques : pluies, températures, durée du jour) et des conditions de culture. En dehors de cette période, le fruit est rare, coûteux, ou disponible uniquement sous forme transformée (séchée, en conserve, en jus).
\end{definition}

\begin{exemplebox}[Quelques fruits de saison au Cameroun]
\begin{itemize}
\item La \textbf{mangue} : pleine saison de mars à juin (jusqu'à juillet dans l'Adamaoua).
\item La \textbf{tomate} : forte production en saison sèche (décembre à avril) grâce aux cultures irriguées ; en saison des pluies, les maladies cryptogamiques (mildiou) ravagent les plants et les prix s'envolent.
\item L'\textbf{avocat} (mai à août), la \textbf{goyave}, le \textbf{corossol}, la \textbf{prune} (safoutier) illustrent le même phénomène.
\end{itemize}
\end{exemplebox}

La conséquence économique directe est une \textbf{forte fluctuation des prix} : la loi de l'offre et de la demande fait s'effondrer les cours en pleine saison et les faire grimper hors saison. Le graphique suivant illustre cette variation pour la mangue sur un marché urbain camerounais.

\begin{popfigure}
\begin{center}
\begin{tikzpicture}
\begin{axis}[
    width=0.95\linewidth, height=6.5cm,
    ybar, bar width=8pt,
    ymin=0, ymax=550,
    ylabel={Prix moyen d'une mangue (FCFA)},
    xlabel={Mois de l'année},
    symbolic x coords={Jan,Fév,Mar,Avr,Mai,Juin,Juil,Août,Sep,Oct,Nov,Déc},
    xtick=data,
    x tick label style={rotate=45, anchor=east, font=\small},
    ytick={0,100,200,300,400,500},
    grid=major, grid style={popline!40},
    every axis plot/.append style={fill=poporange, draw=popdark}
]
\addplot coordinates {(Jan,300) (Fév,350) (Mar,150) (Avr,50) (Mai,30) (Juin,60) (Juil,120) (Août,250) (Sep,350) (Oct,400) (Nov,450) (Déc,500)};
\end{axis}
\end{tikzpicture}
\end{center}
\legende{Variation du prix moyen de la mangue au détail sur un marché urbain (ordres de grandeur documentaires). En pleine saison (avril--juin), l'abondance fait chuter le prix ; hors saison (octobre--décembre), la rareté le multiplie par dix environ. Cette fluctuation est la signature économique d'un fruit de saison.}
\end{popfigure}

\courssub{Le caractère périssable des fruits}

Pourquoi la mangue ne peut-elle pas simplement être stockée dans un grenier, comme le maïs ou le niébé ? La réponse tient à deux propriétés biologiques fondamentales des fruits.

\textbf{a) Une très forte teneur en eau.} Les fruits charnus sont constitués de 80 à 95 \% d'eau. Or l'eau libre est indispensable au développement des micro-organismes (bactéries, moisissures, levures) et au fonctionnement des enzymes. Un milieu riche en eau, en sucres et en vitamines, à température ambiante tropicale (\SIrange{25}{35}{\celsius}), est un véritable bouillon de culture naturel.

\begin{exemplebox}[Teneurs en eau de deux fruits de saison]
\begin{center}
\begin{tabularx}{\linewidth}{|l|X|X|}\hline
\rowcolor{popblueL} \textbf{Fruit} & \textbf{Teneur en eau (pour 100 g)} & \textbf{Conséquence} \\ \hline
Mangue & \SI{83}{\percent} (\SI{83}{g}) & Fruit très fragile, sensible aux chocs et aux moisissures \\ \hline
Tomate & \SI{94}{\percent} (\SI{94}{g}) & Fruit extrêmement périssable : quelques jours à l'air libre \\ \hline
Maïs sec (comparaison) & \SI{12}{\percent} (\SI{12}{g}) & Grain stockable des mois en grenier \\ \hline
\end{tabularx}
\end{center}
\end{exemplebox}

\textbf{b) Un tissu vivant qui respire après récolte.} Contrairement à une idée répandue, un fruit cueilli n'est pas « mort » : ses cellules restent vivantes et continuent de \cle{respirer}, c'est-à-dire d'oxyder leurs réserves glucidiques pour produire de l'énergie selon le bilan :

\[ \ce{C6H12O6 + 6O2 -> 6CO2 + 6H2O} + \text{énergie} \]

Cette respiration consomme les sucres du fruit, dégage de la chaleur et de l'eau, et s'accompagne de la maturation (ramollissement des tissus sous l'action d'enzymes, production d'éthylène, hormone gazeuse de maturation). Chez les fruits dits \emph{climactériques} comme la mangue et la tomate, la respiration s'accélère brutalement après la récolte : le fruit mûrit vite, puis vieillit et devient vulnérable aux pourritures. À la chaleur tropicale, ce processus est très rapide : une mangue mûre se conserve à peine 3 à 5 jours à l'air libre.

\begin{definition}[Fruit climactérique]
Un \cle{fruit climactérique} est un fruit dont la maturation après récolte s'accompagne d'un pic de respiration et d'une production massive d'éthylène (hormone de maturation). La mangue, la tomate, la banane, l'avocat et la papaye sont climactériques ; l'ananas, l'orange, la fraise ne le sont pas (non climactériques).
\end{definition}

\begin{definition}[Fruit périssable]
Un \cle{fruit périssable} est un fruit qui se dégrade rapidement après la récolte (en quelques jours à température ambiante), en raison de sa forte teneur en eau, de la poursuite de son métabolisme (respiration, maturation) et de son attaque par des micro-organismes. La mangue et la tomate sont des fruits périssables ; le maïs sec ou le niébé sont des produits non périssables.
\end{definition}

\begin{attention}[Ne pas confondre conservation et transformation]
Piège classique du Baccalauréat :
\begin{itemize}
\item La \cle{conservation} vise à \textbf{garder le fruit proche de son état initial} (aspect, goût, valeur nutritionnelle) en ralentissant ou en bloquant les altérations : réfrigération, mise en conserve stérilisée, séchage.
\item La \cle{transformation} \textbf{modifie durablement le produit} pour obtenir un nouvel aliment stable : confiture, jus, pâte de tomate, mangue séchée. Le fruit d'origine n'est plus reconnaissable en tant que tel.
\end{itemize}
Dans la pratique, les deux démarches se combinent (une confiture est une transformation qui permet une longue conservation), mais il faut savoir les distinguer conceptuellement.
\end{attention}

\courssub{Le constat chiffré : des pertes post-récolte massives}

\begin{experience}[Document : estimation des pertes post-récolte au Cameroun]
\textbf{Document.} Les enquêtes menées dans les bassins de production fruitière (Nord, Adamaoua, Ouest) convergent vers un ordre de grandeur admis à titre documentaire : \textbf{30 à 50 \% des fruits périssables récoltés sont perdus} entre le verger et le consommateur (pertes dites \emph{post-récolte}).\par
\textbf{Observations.} Les pertes se répartissent à tous les maillons de la chaîne :
\begin{itemize}
\item \emph{à la récolte} : fruits tombés, écrasés, blessés par une cueillette brutale ;
\item \emph{au transport} : routes dégradées, camions surchargés, absence de chaîne du froid, chaleur ;
\item \emph{au stockage et au marché} : entrepôt inadapté, exposition au soleil, attente prolongée faute d'acheteurs ;
\item \emph{chez le consommateur} : conservation domestique limitée à quelques jours.
\end{itemize}
\textbf{Interprétation.} Le caractère périssable du fruit (eau, respiration, micro-organismes) aggrave chaque étape : plus le temps passe, plus la dégradation avance. Sans technique de conservation ou de transformation, l'essentiel de la production excédentaire est condamné.\par
\textbf{Conclusion.} La seule façon de réduire ces pertes est de \emph{stabiliser} le fruit dès la récolte, en bloquant les facteurs d'altération : c'est l'objet des biotechnologies de conservation et de transformation.
\end{experience}

\begin{savaistu}[Des manguiers qui nourrissent... les mouches]
Au Cameroun, les régions du Nord et de l'Adamaoua comptent des millions de manguiers, souvent plantés le long des pistes ou dans les concessions familiales. Chaque année, des milliers de tonnes de mangues pourrissent littéralement sous les arbres, faute de circuits de collecte et d'unités de transformation à proximité. Pourtant, la mangue séchée et le jus de mangue s'exportent à prix d'or depuis d'autres pays d'Afrique de l'Ouest ! Depuis quelques années, des coopératives et des petites unités artisanales de séchage solaire et de production de jus se développent (par exemple autour de Garoua et de Ngaoundéré), montrant qu'une biotechnologie simple peut transformer ce gaspillage en revenus.
\end{savaistu}

\courssub{Les conséquences des pertes : gaspillage, revenus, prix}

Les pertes post-récolte ne sont pas qu'un problème technique : elles ont des conséquences lourdes à trois niveaux.

\begin{itemize}
\item \textbf{Gaspillage alimentaire et nutritionnel.} Les fruits sont des sources irremplaçables de vitamines (A, C), de sels minéraux et de fibres. Perdre 30 à 50 \% de la production, c'est priver les populations de nutriments essentiels, alors que les carences en vitamine A restent un problème de santé publique. Le gaspillage aggrave aussi les nuisances urbaines : les montagnes de fruits pourris attirent les mouches et les rongeurs.
\item \textbf{Pertes de revenus pour les producteurs.} En pleine saison, l'effondrement des prix (voir graphique) empêche le producteur de rentabiliser son travail ; les invendus pourrissent. Sans possibilité de stocker ou de transformer, il est contraint de vendre à n'importe quel prix, immédiatement.
\item \textbf{Flambée des prix hors saison.} Faute de produits conservés ou transformés disponibles toute l'année, le consommateur urbain paie le fruit frais hors saison jusqu'à dix fois son prix de pleine saison — ou s'en prive.
\end{itemize}

\begin{popfigure}
\begin{center}
\begin{tikzpicture}[
  box/.style={draw=popdark, fill=popblueL, rounded corners, minimum width=3.1cm, minimum height=1.2cm, align=center, font=\small},
  perte/.style={draw=popdark, fill=poppinkL, rounded corners, minimum width=2.6cm, minimum height=1.2cm, align=center, font=\small},
  transf/.style={draw=popdark, fill=popgreenL, rounded corners, minimum width=2.6cm, minimum height=1.2cm, align=center, font=\small},
  flux/.style={-{Stealth[length=3mm]}, thick, popdark},
  lbl/.style={font=\footnotesize\itshape, fill=white, inner sep=1pt}
]
\node[box, align=center] (rec) at (0,0){\textbf{Récolte}\\100 \% des fruits};
\node[box, align=center] (circ) at (5.2,0){Transport et\\mise en marché};
\node[box, align=center] (conso) at (10.4,1.6){\textbf{Consommation}\\fraîche\\(50 à 70 \%)};
\node[perte, align=center] (pertes) at (10.4,-1.6){\textbf{Pertes}\\(30 à 50 \%)};
\node[transf, align=center] (transfo) at (5.2,-3.2){\textbf{Transformation}\\(séchage, jus, conserve)\\part encore faible};
\draw[flux] (rec) -- (circ);
\draw[flux] (circ) -- (conso) node[lbl, midway, above, sloped]{fruits sains};
\draw[flux] (circ) -- (pertes) node[lbl, midway, below, sloped]{fruits avariés, invendus};
\draw[flux] (rec) -- (transfo) node[lbl, midway, left]{excédent de pleine saison};
\draw[flux, dashed, popgreen] (transfo.east) -- ++(1.2,0) |- (conso.south) node[lbl, pos=0.25, below]{produits stables toute l'année};
\end{tikzpicture}
\end{center}
\legende{Bilan simplifié des flux de fruits de saison au Cameroun. Sans transformation, 30 à 50 \% de la récolte se perdent entre le verger et le consommateur. La transformation de l'excédent de pleine saison (flèche verte) permet de réduire les pertes et d'étaler la disponibilité des produits sur toute l'année.}
\end{popfigure}

\courssub{Vers la solution : la problématique de la conservation}

Résumons le raisonnement qui structure toute la leçon. Le fruit est un tissu vivant, riche en eau et en nutriments ; après récolte, trois grands facteurs le dégradent : les \cle{micro-organismes} (moisissures, levures, bactéries), les \cle{enzymes} du fruit lui-même (respiration, maturation, brunissement) et l'\cle{oxydation} au contact de l'air. Conserver ou transformer un fruit revient donc toujours, d'une manière ou d'une autre, à \textbf{bloquer ces trois facteurs} : éliminer l'eau disponible (séchage, sucre concentré), tuer ou inactiver les micro-organismes et les enzymes (chauffage, stérilisation), ou priver le milieu d'oxygène (conditionnement hermétique). Ces principes seront étudiés en détail dans les sections suivantes.

\begin{exoresolu}[Calculer l'ampleur des pertes et le gain d'une unité de transformation]
\textbf{Énoncé.} Une coopérative de l'Adamaoua récolte 20 tonnes de mangues en pleine saison. On estime que 40 \% de la récolte serait perdue sans transformation. La coopérative installe un séchoir solaire qui permet de transformer la moitié de l'excédent invendable en mangue séchée. Sachant que le séchage retire environ 80 \% de la masse du fruit frais (eau éliminée), calcule : (a) la masse de mangues perdue sans transformation ; (b) la masse de mangue séchée obtenue à partir de la moitié de cet excédent ; (c) la masse totale de fruits finalement perdue grâce au séchoir.
\tcblower
\textbf{Solution détaillée.}
\begin{enumerate}
\item \emph{Pertes sans transformation :}
\[ m_{\text{pertes}} = 20\ \text{t} \times \frac{40}{100} = 8\ \text{tonnes}. \]
\item \emph{Mangue séchée obtenue :} la moitié de l'excédent, soit $8 / 2 = 4$ tonnes de mangue fraîche, est séchée. Le séchage retire 80 \% de la masse ; il reste donc 20 \% :
\[ m_{\text{séchée}} = 4\ \text{t} \times \frac{20}{100} = 0{,}8\ \text{tonne} = 800\ \text{kg}. \]
\item \emph{Pertes finales :} seule l'autre moitié de l'excédent est perdue :
\[ m_{\text{pertes finales}} = 8\ \text{t} - 4\ \text{t} = 4\ \text{tonnes}. \]
\end{enumerate}
\textbf{Conclusion.} Grâce à une technique simple de transformation (séchage solaire), la coopérative réduit ses pertes de 8 à 4 tonnes, soit une réduction de 50 \%, et obtient 800 kg d'un produit stable, stockable et vendable hors saison à un prix bien supérieur à celui du fruit frais de pleine saison.
\end{exoresolu}

\begin{aretenir}[L'essentiel de la section]
\begin{itemize}
\item Un \cle{fruit de saison} est produit en abondance sur une courte période (mangue : mars--juin ; tomate : saison sèche), ce qui provoque l'effondrement des prix en pleine saison et leur flambée hors saison.
\item Un fruit est \cle{périssable} car il contient 80 à 95 \% d'eau (mangue $\approx$ \SI{83}{\percent}, tomate $\approx$ \SI{94}{\percent}) et reste un tissu vivant qui respire et mûrit après la récolte.
\item La mangue et la tomate sont des fruits \cle{climactériques} : leur maturation post-récolte s'accompagne d'un pic respiratoire et d'une production d'éthylène.
\item Au Cameroun, 30 à 50 \% des fruits périssables sont perdus après la récolte : gaspillage alimentaire, pertes de revenus, prix élevés hors saison.
\item Il faut distinguer la \cle{conservation} (garder le fruit proche de son état) de la \cle{transformation} (modifier durablement le fruit : confiture, jus, pâte, séchage).
\item Problématique centrale : conserver un aliment vivant et riche en eau exige de bloquer les micro-organismes, les enzymes et l'oxydation.
\end{itemize}
\end{aretenir}

\coursec{II. Les facteurs d'altération des fruits}

Nous avons vu dans la section précédente que le Cameroun perd chaque année une fraction considérable de sa production fruitière : manguiers surchargés en saison sèche, tomates qui s'écrasent dans les cagettes entre la ferme et le marché. Mais \emph{pourquoi} un fruit se dégrade-t-il ? Avant d'apprendre à conserver, il faut comprendre ce qui détruit. C'est l'objet de cette section : identifier les \cle{facteurs d'altération} et leurs mécanismes, car toute technique de conservation n'est, en réalité, qu'une réponse ciblée à l'un de ces facteurs.

\courssub{Observation de départ : trois mangues, trois destins}

\begin{experience}[Observation comparée de l'altération des mangues]
\textbf{Matériel :} trois mangues mûres de même variété, récoltées le même jour ; trois assiettes ; un couteau.

\textbf{Protocole :}
\begin{enumerate}
\item Mangue A : laissée \textbf{intacte} (épiderme sain).
\item Mangue B : \textbf{blessée} (entaille de \SI{2}{cm} dans l'épiderme, sans toucher profondément la chair).
\item Mangue C : \textbf{écrasée} (chair largement exposée à l'air).
\end{enumerate}
Les trois fruits sont laissés à l'air libre, à la température ambiante d'une cuisine (environ \SI{28}{\celsius}), et observés chaque jour pendant 5 jours.

\textbf{Observations :}
\begin{center}
\begin{tabularx}{\linewidth}{|l|X|X|X|}
\hline
\rowcolor{popblueL} \textbf{Durée} & \textbf{Mangue A (saine)} & \textbf{Mangue B (blessée)} & \textbf{Mangue C (écrasée)} \\
\hline
Jour 1 & Aucun changement & Brunissement autour de l'entaille & Brunissement étendu, jus qui s'écoule \\
\hline
Jour 2 & Aucun changement & Tache brune élargie, odeur douceâtre & Feutre blanchâtre (moisissures), odeur aigre \\
\hline
Jour 3 & Léger ramollissement & Moisissures vertes au niveau de la plaie & Pourriture complète, liquéfaction \\
\hline
Jour 5 & Taches brunes, ramollissement net & Pourriture étendue à la moitié du fruit & Fruit entièrement détruit \\
\hline
\end{tabularx}
\end{center}

\textbf{Interprétation :} la mangue écrasée pourrit la première, puis la blessée, puis la saine. L'épiderme intact constitue une \textbf{barrière protectrice} : dès qu'elle est rompue, la chair (riche en eau et en sucres) est exposée à l'air et aux micro-organismes qu'il transporte. L'altération est donc d'autant plus rapide que le fruit est blessé.

\textbf{Conclusion :} l'altération d'un fruit résulte d'agents extérieurs (micro-organismes, oxygène de l'air) et d'agents internes (enzymes du fruit lui-même), dont l'action est facilitée par toute rupture de l'épiderme.
\end{experience}

\begin{definition}[Altération d'un fruit]
L'\cle{altération} d'un fruit est l'ensemble des transformations irréversibles qui le rendent impropre à la consommation : modifications de l'aspect (brunissement, moisissures), de la texture (ramollissement, liquéfaction), de l'odeur et du goût (rancissement, aigreur), et perte de valeur nutritionnelle (destruction des vitamines). Elle résulte de l'action combinée de trois grands facteurs : les \textbf{micro-organismes}, les \textbf{enzymes} propres du fruit et l'\textbf{oxygène} de l'air.
\end{definition}

\courssub{Facteur 1 : les micro-organismes}

\textbf{Intuition.} Quand une mangue pourrit, on voit apparaître un « duvet » blanc, vert ou noir, et l'on sent une odeur aigre ou de fermentation. Ces signes trahissent le travail d'êtres vivants microscopiques qui se nourrissent du fruit : les \cle{micro-organismes}.

\begin{definition}[Micro-organismes altérant les fruits]
Les principaux micro-organismes responsables de l'altération des fruits sont :
\begin{itemize}
\item les \textbf{moisissures} (champignons filamenteux) : \emph{Aspergillus} (feutre noirâtre ou jaune-vert), \emph{Penicillium} (feutre vert-bleu, typique des agrumes et des mangues stockées) ;
\item les \textbf{levures} (champignons unicellulaires) : elles provoquent la \textbf{fermentation} des sucres du fruit, d'où une odeur d'alcool et un goût aigre ;
\item les \textbf{bactéries} : certaines provoquent des pourritures molles (liquéfaction des tissus) ou des acidifications.
\end{itemize}
\end{definition}

\textbf{D'où viennent ces micro-organismes ?} Ils sont partout dans l'environnement : dans l'\textbf{air} (spores de moisissures en suspension), dans le \textbf{sol} (d'où l'importance de ne pas stocker les fruits directement à terre), sur les \textbf{mains} du cueilleur ou du vendeur, sur les \textbf{ustensiles} (couteaux, cagettes souillées) et dans l'\textbf{eau} de lavage non potable. C'est pourquoi la mangue blessée de notre observation a pourri avant la mangue saine : la plaie a offert une porte d'entrée directe aux germes de l'air et du couteau.

\begin{propriete}[Conditions de développement microbien]
Les micro-organismes ne se développent que si plusieurs conditions sont réunies :
\begin{itemize}
\item de l'\textbf{eau disponible} : c'est le facteur le plus important, mesuré par l'\cle{activité de l'eau} ($a_w$) ;
\item une \textbf{température favorable} : la plupart des germes d'altération se multiplient activement entre \SI{20}{\celsius} et \SI{40}{\celsius} (optimum souvent vers \SIrange{25}{37}{\celsius}) ;
\item un \textbf{pH} convenable : les bactéries préfèrent un milieu peu acide, tandis que moisissures et levures tolèrent l'acidité des fruits ;
\item de l'\textbf{oxygène} pour les formes aérobies (moisissures notamment) ;
\item des \textbf{nutriments} : sucres, acides organiques, minéraux — le fruit en regorge.
\end{itemize}
Supprimer ou réduire l'une de ces conditions suffit à ralentir ou à bloquer le développement microbien : c'est le fondement de toutes les techniques de conservation.
\end{propriete}

\begin{definition}[Activité de l'eau ($a_w$)]
L'\cle{activité de l'eau}, notée $a_w$, est la fraction d'eau d'un aliment réellement \textbf{disponible} pour les micro-organismes et les réactions chimiques. Elle varie de 0 (produit totalement sec) à 1 (eau pure). Un fruit frais comme la mangue a un $a_w$ voisin de \num{0,97} à \num{0,99} : c'est un milieu idéal pour les germes. En dessous de $a_w = \num{0,60}$, aucun micro-organisme ne peut se développer : c'est le principe du \textbf{séchage}. Le sucre en forte concentration (confitures) et le sel abaissent également l'$a_w$.
\end{definition}

\begin{popfigure}
\begin{center}
\begin{tikzpicture}
\begin{axis}[
  width=0.85\linewidth, height=6.5cm,
  xlabel={Température (\si{\degreeCelsius})},
  ylabel={Vitesse de multiplication microbienne (\%)},
  xmin=-10, xmax=80, ymin=0, ymax=110,
  xtick={-10,0,10,20,30,37,50,60,70,80},
  ytick={0,25,50,75,100},
  grid=both, grid style={popline!40},
  axis lines=left,
]
\addplot[popcurve,domain=-10:80,samples=300]
  {100*exp(-((x-32)^2)/(2*14^2)) * (x<62) };
\addplot[poporange,very thick,domain=62:80,samples=50]{0};
\draw[dashed,popmuted] (axis cs:4,0) -- (axis cs:4,105);
\draw[dashed,popmuted] (axis cs:60,0) -- (axis cs:60,105);
\node[popblue,anchor=south,font=\small] at (axis cs:32,102) {Optimum (\SIrange{25}{37}{\celsius})};
\node[popdark,anchor=west,font=\small,align=left] at (axis cs:-9,60) {Froid :\\ralentissement\\(pas de destruction)};
\node[poporange,anchor=west,font=\small,align=left] at (axis cs:63,60) {Chaleur :\\destruction\\(pasteurisation/stérilisation)};
\end{axis}
\end{tikzpicture}
\legende{Évolution de la vitesse de multiplication des micro-organismes d'altération en fonction de la température. La croissance est maximale vers \SIrange{25}{37}{\celsius} (températures courantes sous nos climats !), fortement ralentie au froid (réfrigérateur vers \SI{4}{\celsius}), et les germes sont détruits par la chaleur au-delà de \SI{60}{\celsius} environ (pasteurisation) ou \SI{100}{\celsius} et plus (stérilisation).}
\end{center}
\end{popfigure}

\begin{attention}[Le froid ne tue pas les micro-organismes]
Erreur classique au Baccalauréat : écrire que « le réfrigérateur tue les microbes ». C'est \textbf{faux}. Le froid (\SI{4}{\celsius}) ne fait que \textbf{ralentir} leur multiplication : les germes restent vivants et reprennent leur activité dès le retour à la température ambiante. Seule la \textbf{chaleur} (pasteurisation vers \SI{85}{\celsius}, stérilisation à plus de \SI{100}{\celsius} en conserve) détruit réellement les micro-organismes. De même, la congélation bloque la multiplication sans stériliser l'aliment.
\end{attention}

\courssub{Facteur 2 : les enzymes propres du fruit}

\textbf{Intuition.} Coupe une mangue ou une banane et laisse-la à l'air : en quelques minutes, la chair brunit. Pourtant, aucun microbe n'a eu le temps de se multiplier ! Le responsable est \emph{à l'intérieur} du fruit : ce sont ses propres \cle{enzymes}.

\begin{definition}[Enzymes d'altération des fruits]
Les \textbf{enzymes} sont des protéines catalysant les réactions biochimiques du fruit. Deux familles jouent un rôle majeur dans l'altération :
\begin{itemize}
\item les \textbf{pectinases} : elles dégradent la \textbf{pectine}, ciment qui lie les cellules végétales entre elles. Leur action provoque le \textbf{ramollissement} progressif du fruit (mangue qui devient « molle » puis liquide) ;
\item les \textbf{polyphénoloxydases} (PPO) : en présence d'oxygène, elles oxydent les composés phénoliques de la chair en pigments bruns (mélanines). C'est le \textbf{brunissement enzymatique}, visible quand on coupe une mangue, une banane ou une pomme.
\end{itemize}
Dans le fruit intact, enzymes et substrats sont séparés dans des compartiments cellulaires distincts ; la coupure ou l'écrasement détruit ces compartiments et met les enzymes en contact avec leurs substrats et avec l'oxygène de l'air.
\end{definition}

\begin{exemplebox}[Le brunissement de la mangue coupée et l'astuce du citron]
Une tranche de mangue exposée à l'air à \SI{28}{\celsius} commence à brunir en \textbf{5 à 10 minutes}. Si l'on arrose une seconde tranche de \textbf{jus de citron}, le brunissement est retardé de plusieurs heures. Double explication : le citron \textbf{acidifie} le milieu (les polyphénoloxydases sont peu actives à pH acide, inférieur à 4) et apporte de la \textbf{vitamine C}, un \textbf{antioxydant} qui capte l'oxygène avant qu'il n'attaque les composés phénoliques. C'est exactement le principe utilisé industriellement lors de la préparation des jus et des confitures.
\end{exemplebox}

\begin{propriete}[Inactivation des enzymes]
Les enzymes sont des protéines : elles sont \textbf{dénaturées} (détruites) par la \textbf{chaleur}. Un chauffage de quelques minutes vers \SIrange{80}{100}{\celsius} (blanchiment, cuisson des confitures, traitement thermique des conserves) inactive durablement pectinases et polyphénoloxydases. Le froid ne fait que ralentir leur activité, et l'acidification (citron, vinaigre) la réduit fortement.
\end{propriete}

\courssub{Facteur 3 : l'oxydation par l'oxygène de l'air}

Même sans microbe et avec des enzymes inactivées, un fruit ou son jus peut encore s'altérer : c'est l'action directe du \textbf{dioxygène} de l'air sur certains constituants.

\begin{definition}[Oxydation des constituants du fruit]
L'\cle{oxydation} est l'action chimique du dioxygène ($O_2$) de l'air sur les molécules du fruit :
\begin{itemize}
\item sur les \textbf{lipides} : elle provoque le \textbf{rancissement}, responsable de goûts et d'odeurs désagréables (important pour les fruits gras comme l'avocat ou les produits séchés riches en huile) ;
\item sur la \textbf{vitamine C} (acide ascorbique) : elle la détruit rapidement. Un jus de mangue laissé à l'air peut perdre plus de la moitié de sa vitamine C en quelques heures, d'où une perte nutritionnelle majeure ;
\item sur les \textbf{pigments} et les arômes : décoloration et perte du parfum caractéristique.
\end{itemize}
Ces réactions sont accélérées par la lumière et la chaleur.
\end{definition}

\courssub{Facteur complémentaire : la maturation poursuivie après récolte}

Un fruit cueilli n'est pas un objet inerte : il reste \textbf{vivant} et continue à \textbf{respirer} (il consomme $O_2$ et ses propres sucres, et rejette $CO_2$, de l'eau et de la chaleur). De plus, les fruits dits \emph{climactériques} comme la mangue, la banane et la tomate produisent après récolte une hormone gazeuse, l'\textbf{éthylène}, qui accélère leur propre mûrissement puis leur \textbf{sénescence} (surmaturation, ramollissement, sensibilité accrue aux infections). Un fruit mûr placé dans une cagette accélère donc le mûrissement — puis la pourriture — de tous ses voisins : c'est le phénomène bien connu des vendeuses du marché qui trient quotidiennement leurs fruits.

\courssub{Synthèse : trois ennemis, trois stratégies de lutte}

\begin{popfigure}
\begin{center}
\begin{tikzpicture}[
  facteur/.style={rectangle, rounded corners=4pt, draw=popblue, fill=popblueL, thick, minimum width=3.6cm, minimum height=1.1cm, align=center, font=\small\bfseries},
  manif/.style={rectangle, rounded corners=4pt, draw=poporange, fill=poporangeL, thick, minimum width=3.6cm, minimum height=1.6cm, align=center, font=\small},
  lutte/.style={rectangle, rounded corners=4pt, draw=popgreen, fill=popgreenL, thick, minimum width=3.6cm, minimum height=1.6cm, align=center, font=\small},
  lien/.style={-{Stealth[length=2.5mm]}, thick, popdark}
]
\node[facteur, align=center] (f1) at (0,0){Micro-organismes\\(moisissures, levures,\\bactéries)};
\node[facteur, align=center] (f2) at (6.2,0){Enzymes du fruit\\(pectinases,\\polyphénoloxydases)};
\node[facteur, align=center] (f3) at (12.4,0){Oxygène de l'air\\($O_2$)};
\node[manif, align=center] (m1) at (0,-2.6){Moisissures visibles,\\fermentation,\\pourriture molle};
\node[manif, align=center] (m2) at (6.2,-2.6){Ramollissement,\\brunissement\\à l'air};
\node[manif, align=center] (m3) at (12.4,-2.6){Rancissement,\\destruction de la\\vitamine C};
\node[lutte, align=center] (l1) at (0,-5.2){Destruction (chaleur)\\ou inhibition (froid,\\séchage, sucre, acidité)};
\node[lutte, align=center] (l2) at (6.2,-5.2){Inactivation par la\\chaleur ; acidification\\(jus de citron)};
\node[lutte, align=center] (l3) at (12.4,-5.2){Exclusion de l'air\\(boîtes hermétiques,\\sirop, antioxydants)};
\draw[lien] (f1) -- (m1);
\draw[lien] (f2) -- (m2);
\draw[lien] (f3) -- (m3);
\draw[lien, popgreen!70!black] (m1) -- (l1);
\draw[lien, popgreen!70!black] (m2) -- (l2);
\draw[lien, popgreen!70!black] (m3) -- (l3);
\end{tikzpicture}
\legende{Les trois grands facteurs d'altération des fruits, leurs manifestations visibles (ligne médiane) et les principes de lutte correspondants (ligne inférieure). Toute technique de conservation agit sur au moins un de ces trois facteurs.}
\end{center}
\end{popfigure}

\begin{methode}[Raisonner en « facteur $\rightarrow$ principe de lutte »]
Face à une question de conservation au Baccalauréat, procède toujours en trois étapes :
\begin{enumerate}
\item \textbf{Identifie le facteur d'altération} mis en jeu dans la situation (microbe, enzyme ou oxygène).
\item \textbf{Énonce le principe de lutte} correspondant : microbe $\rightarrow$ destruction (chaleur) ou inhibition (froid, baisse de l'$a_w$ par séchage ou sucre, acidification) ; enzyme $\rightarrow$ inactivation thermique ou acidification ; oxygène $\rightarrow$ exclusion de l'air (récipient hermétique, recouvrement par un sirop) ou antioxydant.
\item \textbf{Relie la technique au principe} : par exemple, la mise en conserve combine destruction microbienne par la chaleur \emph{et} exclusion de l'air par la fermeture hermétique — elle agit donc sur deux facteurs à la fois.
\end{enumerate}
\end{methode}

\begin{aretenir}[L'essentiel de la section]
\begin{itemize}
\item L'\textbf{altération} d'un fruit résulte de l'action combinée de \textbf{trois facteurs} : les micro-organismes, les enzymes du fruit et l'oxygène de l'air.
\item L'épiderme intact est une barrière : un fruit \textbf{blessé ou écrasé} s'altère beaucoup plus vite.
\item Les micro-organismes exigent de l'\textbf{eau disponible} ($a_w$ élevée), une température favorable, des nutriments et souvent de l'oxygène.
\item Le \textbf{froid ralentit} sans tuer ; seule la \textbf{chaleur détruit} micro-organismes et enzymes.
\item Le \textbf{brunissement} est enzymatique (polyphénoloxydases) ; le \textbf{rancissement} et la perte de vitamine C sont des oxydations.
\item Toute technique de conservation se justifie par le ou les facteurs sur lesquels elle agit.
\end{itemize}
\end{aretenir}

\begin{exoresolu}[Analyser une situation d'altération]
\textbf{Énoncé.} Une coopérative de Ngaoundéré transforme des mangues en jus. Le jus, mis en bouteilles non bouchées, est stocké à la température ambiante (\SI{30}{\celsius}). Au bout de deux jours, les bouteilles dégagent une odeur de fermentation et le jus devient trouble. (1) Identifier le facteur d'altération principal et les conditions qui l'ont favorisé. (2) Proposer deux mesures, en précisant le principe de lutte de chacune.
\tcblower
\textbf{Solution détaillée.}

(1) L'odeur de fermentation et le trouble révèlent le développement de \textbf{micro-organismes}, en particulier des \textbf{levures} qui transforment les sucres du jus en alcool et en $CO_2$. Les conditions favorables étaient toutes réunies :
\begin{itemize}
\item eau disponible très élevée ($a_w \approx \num{0,97}$ dans un jus) ;
\item température de \SI{30}{\celsius}, proche de l'optimum de multiplication microbienne ;
\item sucres et nutriments abondants ;
\item apport continu d'oxygène et de germes de l'air, car les bouteilles n'étaient pas bouchées.
\end{itemize}

(2) Mesures correctives :
\begin{itemize}
\item \textbf{Chauffer le jus} (pasteurisation vers \SI{85}{\celsius} pendant quelques minutes) : principe = \emph{destruction} des micro-organismes et \emph{inactivation} des enzymes par la chaleur ;
\item \textbf{Boucher hermétiquement les bouteilles} immédiatement à chaud : principe = \emph{exclusion de l'air} (lutte contre l'oxygène et contre la recontamination) ;
\item à défaut, \textbf{réfrigérer} vers \SI{4}{\celsius} : principe = \emph{ralentissement} de la multiplication microbienne (sans destruction des germes).
\end{itemize}
\end{exoresolu}

\begin{exercice}[Pour t'entraîner]
Un commerçant de Douala constate que ses tomates entières se conservent une semaine, alors que les tomates coupées s'altèrent en un jour. Explique cette différence en mobilisant les trois facteurs d'altération, puis indique pourquoi la transformation des tomates en pâte bouillie, mise en bocaux fermés, permet une conservation de plusieurs mois.
\end{exercice}

\begin{corrige}[Exercice]
La tomate coupée a perdu sa barrière protectrice (l'épiderme) : les micro-organismes de l'air, des mains et du couteau pénètrent directement dans la chair riche en eau et en nutriments ; la coupure met en contact les enzymes (pectinases, polyphénoloxydases) avec leurs substrats et l'oxygène, d'où ramollissement et brunissement rapides ; enfin, la surface exposée s'oxyde directement. La pâte de tomate bouillie a subi une chaleur qui \textbf{détruit les micro-organismes} et \textbf{inactive les enzymes} ; la fermeture hermétique du bocal \textbf{exclut l'oxygène} et empêche toute recontamination. Les trois facteurs d'altération étant neutralisés, la conservation de plusieurs mois devient possible.
\end{corrige}

\begin{savaistu}[Le lac Nyos et... la conservation ?]
Si le froid ne tue pas les microbes, la chaleur, elle, est une arme ancienne et universelle : bien avant Pasteur, les populations camerounaises faisaient bouillir les jus et séchaient au soleil mangues, tomates et feuilles de légumes. Le séchage solaire, pratiqué dans les savanes du Nord, abaisse l'$a_w$ en dessous de \num{0,60} : c'est de la microbiologie appliquée sans le savoir ! Comprendre les facteurs d'altération permet aujourd'hui d'améliorer ces techniques traditionnelles avec rigueur — c'est précisément l'objet des sections suivantes.
\end{savaistu}

\coursec{III. Les principes scientifiques de la conservation}

Nous avons vu dans la section précédente que les fruits de saison s'altèrent sous l'action combinée de trois grands facteurs : les \cle{micro-organismes} (bactéries, levures, moisissures), les \cle{enzymes} propres du fruit (responsables du brunissement, du ramollissement, des modifications de goût) et les \cle{oxydations} chimiques au contact de l'air. Une question surgit alors naturellement, celle que se posent chaque année les ménagères de Ngaoundéré au cœur de la saison des mangues : \emph{comment empêcher ces agents de dégradation d'agir ?} Toute la science de la conservation tient dans la réponse à cette question. C'est ce que nous allons établir maintenant, en nous appuyant sur des expériences simples et quantitatives.

\courssub{Le principe général : créer des conditions défavorables}

\begin{definition}[Principe général de la conservation]
Conserver un fruit, c'est créer autour de lui et en lui des \cle{conditions défavorables} aux micro-organismes et aux enzymes, soit en les \textbf{détruisant} (action létale : chaleur forte), soit en \textbf{bloquant leur activité} sans les tuer (\cle{anabiose} : froid, dessiccation, concentration en sucre, acidité), soit en les \textbf{empêchant d'accéder} au produit (conditionnement hermétique). Un procédé efficace agit toujours sur au moins l'un des trois facteurs d'altération : micro-organismes, enzymes, oxygène.
\end{definition}

L'idée intuitive est simple : un micro-organisme, comme tout être vivant, a besoin d'\textbf{eau disponible}, d'une \textbf{température} convenable, d'un \textbf{pH} proche de la neutralité et souvent d'\textbf{oxygène}. Retirer l'un de ces éléments, ou le rendre insupportable, suffit à stopper la dégradation. Quant aux enzymes, ce sont des protéines : toute condition qui détruit leur structure spatiale (chaleur) ou ralentit leur fonctionnement (froid) supprime leur action. Voyons maintenant chaque levier en détail.

\courssub{L'action de la chaleur : destruction des micro-organismes et inactivation des enzymes}

\textbf{Observation.} Un jus de tomate fraîchement pressé fermente et tourne en 2 à 3 jours à la température ambiante de Yaoundé (environ \SI{25}{\celsius}). Le même jus, porté à ébullition puis versé bouillant dans une bouteille fermée, se conserve des mois. Pourquoi ?

\begin{experience}[Effet de la température de chauffage sur la survie microbienne]
\textbf{Protocole.} On prépare un jus de tomate homogène que l'on répartit en 5 lots. Chaque lot est chauffé 10 minutes à une température différente, refroidi, puis ensemencé sur milieu de culture gélosé (boîtes de Pétri) et incubé 48 h à \SI{30}{\celsius}.

\begin{center}
\begin{tabularx}{\linewidth}{|c|X|X|}
\hline
\rowcolor{popblueL} \textbf{Lot} & \textbf{Traitement thermique} & \textbf{Résultat après 48 h} \\
\hline
A (témoin) & Aucun chauffage & Culture très abondante (innombrables colonies) \\
\hline
B & \SI{40}{\celsius} pendant 10 min & Culture abondante \\
\hline
C & \SI{60}{\celsius} pendant 10 min & Quelques colonies \\
\hline
D & \SI{75}{\celsius} pendant 10 min & Aucune colonie visible \\
\hline
E & \SI{120}{\celsius} pendant 15 min (autoclave) & Aucune colonie, même après 7 jours \\
\hline
\end{tabularx}
\end{center}

\textbf{Interprétation.} Le développement microbien diminue quand la température de traitement augmente : la chaleur détruit les cellules microbiennes par \cle{dénaturation} de leurs protéines et de leurs enzymes. À \SI{75}{\celsius}, les formes végétatives et les germes pathogènes sont éliminés ; à \SI{120}{\celsius} sous pression, même les \textbf{spores} (formes de résistance) sont détruites.

\textbf{Conclusion.} Il existe un couple \textbf{température--durée} au-delà duquel plus aucun micro-organisme ne se développe : c'est le fondement scientifique de la pasteurisation et de la stérilisation.
\end{experience}

\begin{definition}[Pasteurisation et stérilisation]
La \cle{pasteurisation} est un chauffage modéré (typiquement \SIrange{62}{85}{\celsius} pendant quelques secondes à quelques minutes) qui détruit les \textbf{germes pathogènes} et la plupart des formes végétatives, sans détruire toutes les spores : le produit doit être conservé au frais et consommé assez rapidement (jus, lait). La \cle{stérilisation} est un chauffage plus intense (\SIrange{110}{121}{\celsius} en autoclave ou en stérilisateur) qui détruit \textbf{tous les micro-organismes et leurs spores} : le produit, maintenu en récipient hermétique (conserve, bocal), est stable à température ambiante pendant des mois, voire des années.
\end{definition}

La chaleur n'agit pas seulement sur les microbes : elle agit aussi sur les \textbf{enzymes du fruit}. C'est ce que montre l'expérience quotidienne du brunissement : une tranche de mangue ou d'avocat exposée à l'air brunit en quelques heures (action de la \emph{polyphénol oxydase}), alors qu'une tranche plongée 2 minutes dans l'eau bouillante garde sa couleur. Cette opération s'appelle le \cle{blanchiment} (ou ébouillantage bref) : on la pratique systématiquement avant le séchage ou la congélation.

\begin{propriete}[Dénaturation thermique des enzymes]
L'activité d'une enzyme croît avec la température jusqu'à un \textbf{optimum} (vers \SIrange{35}{45}{\celsius} pour les enzymes végétales), puis s'effondre brutalement : au-delà d'une \textbf{température seuil} (environ \SIrange{60}{70}{\celsius}), la chaleur détruit la structure spatiale de la protéine enzymatique. Cette \cle{dénaturation} est \textbf{irréversible} : même refroidie, l'enzyme ne retrouve jamais son activité. C'est pourquoi le blanchiment protège durablement le fruit du brunissement, contrairement à la simple réfrigération.
\end{propriete}

\begin{popfigure}
\begin{center}
\begin{tikzpicture}
\begin{axis}[
  width=0.85\linewidth, height=6cm,
  xlabel={Température (\si{\celsius})},
  ylabel={Activité enzymatique (\%)},
  xmin=0, xmax=100, ymin=0, ymax=110,
  xtick={0,20,40,60,80,100},
  ytick={0,25,50,75,100},
  grid=both, grid style={popline!40},
  legend style={at={(0.02,0.98)}, anchor=north west, font=\small}
]
\addplot[popcurve, domain=0:100, samples=200]
  {100*exp(-((x-40)/28)^2)};
\addlegendentry{Activité de la polyphénol oxydase}
\draw[dashed, poporange, thick] (axis cs:65,0) -- (axis cs:65,105)
  node[pos=0.95, anchor=south west, font=\small, text=poporange]{seuil de dénaturation ($\approx$ \SI{65}{\celsius})};
\node[font=\small, text=popdark] at (axis cs:40,105) {optimum};
\end{axis}
\end{tikzpicture}
\end{center}
\legende{Activité d'une enzyme de brunissement en fonction de la température. L'activité est maximale vers \SI{40}{\celsius}, puis chute par dénaturation irréversible au-delà de \SIrange{60}{70}{\celsius} : c'est le principe du blanchiment.}
\end{popfigure}

\begin{attention}[Chaleur et qualité du produit]
Un chauffage trop long ou trop fort détruit aussi les \textbf{vitamines} (notamment la vitamine C, thermosensible) et modifie le goût. L'art du conserveur consiste à choisir le couple température--durée \emph{juste suffisant} : assez pour tuer les germes et inactiver les enzymes, pas plus. C'est pourquoi on préfère souvent « température élevée, temps court ».
\end{attention}

\courssub{L'action du froid : un ralentissement, pas un arrêt}

Contrairement à la chaleur, le froid ne détruit presque rien : il \textbf{ralentit}. La \cle{réfrigération} (\SIrange{0}{8}{\celsius}) diminue la vitesse des réactions enzymatiques et la multiplication microbienne ; la \cle{congélation} (\SI{-18}{\celsius}) les réduit à presque rien, car l'eau devient glace et n'est plus disponible. Mais dès le dégel, micro-organismes survivants et enzymes reprennent leur activité : le froid est une \textbf{anabiose réversible}. De plus, la congélation éclate les tissus végétaux (cristaux de glace), ce qui ramollit le fruit : on congèle donc plutôt des purées ou des jus que des mangues entières.

\courssub{L'action de la déshydratation : retirer l'eau disponible}

\textbf{Observation.} Au marché de Maroua, on vend toute l'année des tranches de mangue séchées au soleil qui ne moisissent jamais, alors qu'une mangue fraîche pourrit en une semaine. Pourtant, le séchage ne tue pas les microbes : il les prive d'eau.

\begin{experience}[Effet du séchage sur le développement des moisissures]
\textbf{Protocole.} On découpe des tranches de mangue de même taille en trois lots : lot 1 frais (non séché), lot 2 partiellement séché (2 h au soleil), lot 3 bien séché (jusqu'à texture souple et sèche). Les tranches sont placées dans des boîtes fermées humides, à température ambiante, et observées pendant 15 jours.

\textbf{Résultats.} Lot 1 : moisissures abondantes dès le 3\textsuperscript{e} jour. Lot 2 : moisissures vers le 7\textsuperscript{e} jour. Lot 3 : aucune moisissure après 15 jours.

\textbf{Interprétation.} Ce qui compte n'est pas la quantité totale d'eau, mais l'eau \emph{disponible} pour les micro-organismes, mesurée par l'\cle{activité de l'eau} ($a_w$), rapport entre la pression de vapeur d'eau du produit et celle de l'eau pure ($0 \leq a_w \leq 1$ ; $a_w = 1$ pour l'eau pure, environ 0,99 pour un fruit frais, moins de 0,6 pour un fruit bien séché). En dessous d'un certain seuil de $a_w$, les micro-organismes ne peuvent plus prélever l'eau nécessaire à leur métabolisme.
\end{experience}

\begin{popfigure}
\begin{center}
\begin{tikzpicture}
\begin{axis}[
  width=0.85\linewidth, height=6.5cm,
  xlabel={Activité de l'eau $a_w$},
  ylabel={Vitesse de croissance microbienne (\%)},
  xmin=0.5, xmax=1, ymin=0, ymax=110,
  xtick={0.5,0.6,0.7,0.8,0.9,1},
  ytick={0,25,50,75,100},
  grid=both, grid style={popline!40},
  legend style={at={(0.02,0.98)}, anchor=north west, font=\small}
]
\addplot[popcurve, domain=0.88:1, samples=100] {100*(1-exp(-25*(x-0.88)))};
\addlegendentry{Bactéries (seuil $\approx$ 0,90)}
\addplot[popcurve, poporange, domain=0.68:1, samples=100] {100*(1-exp(-12*(x-0.68)))};
\addlegendentry{Moisissures (seuil $\approx$ 0,70)}
\draw[dashed, poppurple, thick] (axis cs:0.6,0) -- (axis cs:0.6,105)
  node[pos=0.9, anchor=south west, font=\small, text=poppurple]{fruit bien séché ($a_w < 0{,}6$)};
\end{axis}
\end{tikzpicture}
\end{center}
\legende{Croissance microbienne en fonction de l'activité de l'eau. En dessous de $a_w \approx 0{,}9$, les bactéries ne se développent plus ; les moisissures, plus résistantes, exigent $a_w < 0{,}7$. Un séchage poussé ($a_w < 0{,}6$) met le produit à l'abri de tous les germes.}
\end{popfigure}

\courssub{L'action du sucre et du sel : la plasmolyse par osmose}

Pourquoi une confiture ne moisit-elle pas alors qu'elle contient encore beaucoup d'eau ? Parce que cette eau est \emph{piégée} par le sucre. Placée dans un sirop très concentré, une cellule microbienne perd son eau par \cle{osmose} : l'eau sort de la cellule vers le milieu extérieur plus concentré en solutés. La membrane cytoplasmique se décolle de la paroi : c'est la \cle{plasmolyse}. Déshydratée de l'intérieur, la cellule ne peut plus se multiplier. Le sel agit de la même manière (salaisons, saumures).

\begin{popfigure}
\begin{center}
\begin{tikzpicture}[scale=1]
% Cellule normale
\draw[fill=popgreenL, thick, popdark] (0,0) ellipse (1.6 and 1.1);
\draw[fill=popgreen!50, thick, popdark] (0,0) ellipse (1.35 and 0.9);
\draw[fill=poppurple!40, popdark] (0,0) ellipse (0.45 and 0.3);
\node[font=\small, text=popdark] at (0,-1.6) {Cellule en milieu isotonique};
\node[font=\footnotesize] at (2.1,0.9) {paroi};
\node[font=\footnotesize] at (2.3,-0.6) {cytoplasme};
\node[font=\footnotesize] at (0.9,0.1) {noyau};
% Flèche
\draw[-{Stealth[length=4mm]}, very thick, poporange] (2.6,0) -- (4.4,0)
  node[midway, above, font=\small, text=popdark]{sirop concentré};
\node[font=\footnotesize, text=popdark, align=center] at (3.5,-0.5) {sortie d'eau\\ par osmose};
% Cellule plasmolysée
\draw[fill=popgreenL, thick, popdark] (7,0) ellipse (1.6 and 1.1);
\draw[fill=popgreen!50, thick, popdark, decorate, decoration={bumps, amplitude=2pt, segment length=8pt}] (7,0) ellipse (1.0 and 0.62);
\draw[fill=poppurple!40, popdark] (7,0) ellipse (0.35 and 0.22);
\node[font=\small, text=popdark] at (7,-1.6) {Cellule plasmolysée};
\draw[-{Stealth}, popblue, thick] (7,0.85) -- (7,1.35);
\draw[-{Stealth}, popblue, thick] (6.2,0.55) -- (5.85,0.95);
\draw[-{Stealth}, popblue, thick] (7.8,0.55) -- (8.15,0.95);
\node[font=\footnotesize, text=popblue] at (8.6,1.3) {eau};
\end{tikzpicture}
\end{center}
\legende{Plasmolyse d'une cellule microbienne placée dans un sirop concentré. L'eau sort de la cellule par osmose (flèches bleues) ; le cytoplasme se rétracte et se décolle de la paroi : la cellule, déshydratée, ne peut plus se multiplier.}
\end{popfigure}

\begin{exemplebox}[Le seuil sucrier des confitures]
Une confiture est microbiologiquement stable lorsque sa \textbf{teneur finale en sucre} atteint environ \textbf{60 à 65 \%} (soit 60 à 65 degrés Brix). En dessous de 55 \%, des moisissures et levures osmotolérantes peuvent encore se développer en surface. Cette concentration se mesure facilement au \cle{réfractomètre}, petit appareil optique utilisé par les unités artisanales de transformation de mangues de l'Adamaoua. Astuce traditionnelle : le sucre ajouté au départ (environ 800 g à 1 kg de sucre par kg de fruit) compense l'eau évaporée pendant la cuisson pour atteindre ce seuil final.
\end{exemplebox}

\courssub{L'action de l'acidité et l'exclusion de l'oxygène}

\textbf{L'acidité.} La plupart des bactéries pathogènes ou d'altération se développent mal lorsque le \cle{pH} descend en dessous de 4,5. C'est pourquoi les fruits naturellement acides (tomate, pH $\approx$ 4,2 ; citron, pH $\approx$ 2,3 ; mangue, pH $\approx$ 4) se conservent mieux et demandent des traitements thermiques moins sévères que les légumes peu acides. On peut renforcer cette barrière en \textbf{ajoutant du jus de citron} à une purée, ou en laissant des \textbf{bactéries lactiques} fermenter le produit (choucroute, lait caillé, boissons fermentées) : ces bactéries transforment les sucres en acide lactique, abaissant le pH et bloquant leurs concurrentes.

\textbf{L'exclusion de l'oxygène.} De nombreux germes d'altération (moisissures, bactéries aérobies) et les oxydations chimiques (rancissement, brunissement) exigent de l'oxygène. On le supprime par :
\begin{itemize}
\item le \textbf{conditionnement hermétique} (bocaux, conserves) : le chauffage chasse l'air, puis le refroidissement crée un vide qui plaque le couvercle ;
\item la \textbf{couche d'huile de couverture} sur une pâte ou une purée, qui isole la surface de l'air ;
\item l'\textbf{emballage sous vide} ou sous sachet bien fermé pour les fruits séchés.
\end{itemize}

\begin{attention}[Stérilisé ne veut pas dire éternel]
Un bocal stérilisé n'est stable que \textbf{tant qu'il reste fermé}. Dès l'ouverture, l'air (et donc l'oxygène) et les micro-organismes de l'environnement pénètrent : le produit redevient périssable et doit être conservé au réfrigérateur et consommé en quelques jours. De même, un bocal dont le couvercle « bombe » ou qui « pétille » à l'ouverture a fermenté : il faut le jeter sans hésiter (risque de botulisme, intoxication mortelle).
\end{attention}

\courssub{Synthèse : relier facteur d'altération, procédé et principe}

Le tableau suivant rassemble le savoir central de cette leçon : chaque procédé n'est que l'application concrète d'un principe scientifique visant un facteur d'altération précis.

\begin{center}
\rowcolors{2}{popcream}{white}
\begin{tabularx}{\linewidth}{|X|X|X|}
\hline
\rowcolor{popblueL} \textbf{Facteur d'altération visé} & \textbf{Procédé} & \textbf{Principe scientifique} \\
\hline
Micro-organismes et enzymes & Stérilisation (conserve, \SI{115}{}--\SI{121}{\celsius}) & Destruction par dénaturation thermique irréversible \\
\hline
Germes pathogènes & Pasteurisation (\SI{62}{}--\SI{85}{\celsius}) & Destruction des formes végétatives, spores épargnées \\
\hline
Enzymes (brunissement) & Blanchiment avant séchage ou congélation & Dénaturation irréversible des enzymes au-delà de \SI{60}{}--\SI{70}{\celsius} \\
\hline
Micro-organismes et enzymes & Réfrigération, congélation & Ralentissement réversible des réactions (anabiose) \\
\hline
Micro-organismes & Séchage (soleil, séchoir) & Abaissement de l'$a_w$ sous les seuils de croissance \\
\hline
Micro-organismes & Confiture, saumure, salaison & Plasmolyse osmotique par concentration en sucre ou en sel ($\approx$ 60--65 \% de sucre) \\
\hline
Bactéries & Acidification (citron, fermentation lactique) & Abaissement du pH sous 4,5 \\
\hline
Germes aérobies et oxydations & Bocal hermétique, huile de couverture, vide & Exclusion de l'oxygène \\
\hline
\end{tabularx}
\end{center}

\begin{methode}[Choisir le bon procédé pour un fruit donné]
\begin{enumerate}
\item \textbf{Identifier le facteur d'altération dominant} du fruit considéré : la tomate, très acide et très aqueuse, craint surtout les moisissures et levures ; la mangue, riche en enzymes de brunissement, brunit vite à la coupe.
\item \textbf{Choisir le principe adapté} : chaleur pour un produit liquide destiné à la conserve ; séchage si l'on veut un produit léger et transportable ; sucre si l'on vise une confiture.
\item \textbf{Combiner les barrières} : les meilleurs procédés cumulent plusieurs obstacles (ex. confiture = chaleur + sucre + acidité + bocal hermétique). C'est le concept des « technologies combinées ».
\item \textbf{Vérifier l'efficacité} : contrôle du pH, de la concentration en sucre (réfractomètre), de l'étanchéité du bocal.
\end{enumerate}
\end{methode}

\begin{exoresolu}[Justifier un choix de procédé]
Une coopérative de femmes de l'Extrême-Nord récolte chaque année un surplus de tomates en saison sèche. Les tomates pourrissent en moins d'une semaine faute de réfrigérateurs. Proposer deux procédés de conservation adaptés et justifier chacun par le principe scientifique correspondant.
\tcblower
\textbf{Procédé 1 : la conserve de purée de tomate stérilisée.} La tomate est acide (pH $\approx$ 4,2), ce qui facilite le traitement thermique. On concentre la pulpe par cuisson (ce qui détruit germes et enzymes par dénaturation thermique et abaisse l'$a_w$), on la verse bouillante dans des bocaux que l'on stérilise 30 à 45 min à \SI{100}{\celsius}. Principe : destruction microbienne par la chaleur + exclusion de l'oxygène par fermeture hermétique + barrière acide naturelle du fruit. Le produit est stable des mois à température ambiante.

\textbf{Procédé 2 : la tomate séchée.} On tranche les tomates, on les blanchit brièvement (inactivation des enzymes de brunissement), puis on les sèche au soleil sur claies propres jusqu'à $a_w < 0{,}7$. Principe : abaissement de l'activité de l'eau sous le seuil de croissance des moisissures. Le produit, léger et stable, se conserve en sachet fermé.

Dans les deux cas, aucun réfrigérateur n'est nécessaire : les procédés reposent sur des principes indépendants du froid, adaptés au contexte local.
\end{exoresolu}

\begin{aretenir}[L'essentiel de la section]
Conserver, c'est créer des conditions défavorables : \textbf{détruire} (chaleur : pasteurisation, stérilisation), \textbf{bloquer} (froid, dessiccation, sucre, sel, acidité) ou \textbf{isoler} (emballage hermétique). La chaleur dénature irréversiblement enzymes et protéines microbiennes ; le froid ne fait que ralentir ; le séchage abaisse l'$a_w$ ; le sucre et le sel provoquent la plasmolyse osmotique ; l'acidité (pH $< 4{,}5$) inhibe les bactéries ; l'exclusion de l'oxygène bloque moisissures et oxydations. Le choix d'un procédé part toujours de l'identification du facteur d'altération dominant du fruit.
\end{aretenir}

Ces principes étant établis, nous pouvons maintenant les mettre en œuvre concrètement : la section suivante décrira pas à pas les techniques de conservation et de transformation appliquées à nos deux fruits d'étude, la mangue et la tomate.

\coursec{IV. Les techniques de conservation et de transformation : cas de la mangue et de la tomate}

Dans la section précédente, nous avons établi les trois grands principes scientifiques de la conservation : \textbf{abaisser l'activité de l'eau} ($a_w$) pour priver les micro-organismes de l'eau disponible, \textbf{détruire les micro-organismes et inactiver les enzymes} par la chaleur, et \textbf{maintenir le produit à l'abri} de toute recontamination (emballage hermétique, hygiène). Il est temps maintenant de voir comment ces principes se traduisent en \cle{techniques concrètes}, applicables dès le village ou l'atelier scolaire, à deux fruits emblématiques du Cameroun : la \cle{mangue}, dont l'abondance éphémère sature les marchés de mars à juillet, et la \cle{tomate}, produite massivement dans l'Ouest et le Nord mais très périssable.

\courssub{Les étapes communes préalables : l'hygiène, première barrière}

Quelle que soit la technique choisie, toute transformation commence par une série d'opérations préparatoires dont la rigueur conditionne la qualité finale du produit.

\begin{enumerate}
\item \textbf{Le tri} : on élimine les fruits pourris, tachés, perforés ou attaqués par les mouches et les larves. Un seul fruit pourri peut contaminer toute une cuvée : c'est le principe de la \emph{barrière initiale}.
\item \textbf{Le lavage} : à l'eau potable, éventuellement additionnée d'une solution chlorée (eau de Javel diluée, quelques gouttes par litre, suivie d'un rinçage), pour réduire la charge microbienne de surface (terre, spores, levures).
\item \textbf{L'épluchage et le dénoyautage} : pour la mangue, on retire l'épiderme (riche en latex irritant) et le noyau ; pour la tomate, on peut ébouillanter le fruit quelques secondes pour faciliter le pelage.
\item \textbf{Le découpage} : rondelles ou lamelles pour le séchage, morceaux pour la confiture, broyage pour les jus et les pâtes.
\item \textbf{Le blanchiment éventuel} (voir définition ci-dessous), surtout avant séchage.
\end{enumerate}

\begin{definition}[Blanchiment et extrait sec]
Le \cle{blanchiment} (ou ébouillantage) consiste à plonger brièvement les morceaux de fruits ou de légumes dans l'eau bouillante (1 à 5 min) puis à les refroidir rapidement. Il \textbf{inactive les enzymes} responsables du brunissement et du ramollissement, fixe la couleur et réduit la flore microbienne de surface.\par
L'\cle{extrait sec} est la masse de matière sèche (sucres, acides, protéines, minéraux, fibres) contenue dans \SI{100}{g} de produit, exprimée en pourcentage. La tomate fraîche contient environ 6 \% d'extrait sec ; une pâte de tomate double concentré en contient 28 à 30 \%.
\end{definition}

\begin{attention}[L'hygiène n'est pas un détail]
Mains propres, ustensiles lavés et ébouillantés, plans de travail désinfectés, eau potable : la contamination post-traitement est la première cause d'échec des conserves artisanales. Un produit parfaitement stérilisé mais rempli dans un bocal sale se gâtera en quelques jours.
\end{attention}

\courssub{Le séchage : conserver en retirant l'eau}

\textbf{Principe.} Le séchage abaisse l'activité de l'eau ($a_w$) du fruit d'environ 0,98 (fruit frais) à moins de 0,60 : en dessous de ce seuil, aucune bactérie et presque aucune moisissure ne peut se développer. L'eau passe du fruit vers l'air ambiant grâce à la chaleur et à la ventilation ; plus l'air est chaud, sec et renouvelé, plus le séchage est rapide et sûr.

\textbf{Techniques.}
\begin{itemize}
\item Le \textbf{séchage solaire sur claies} : les rondelles sont disposées en couche fine sur des claies (grillage ou vannerie) surélevées, couvertes d'un filet anti-insectes, et retournées régulièrement. Simple et gratuit, mais lent (2 à 4 jours), dépendant de la météo et exposé aux poussières.
\item Le \textbf{séchoir solaire amélioré (tunnel)} : un capteur peint en noir absorbe le rayonnement solaire et chauffe l'air, qui circule par convection à travers les claies placées sous un abri transparent ; l'air humide s'échappe par une cheminée. Le produit est protégé de la pluie, de la poussière et des insectes, et le séchage est plus rapide (températures internes de 45 à \SI{60}{\celsius}).
\end{itemize}

\begin{popfigure}
\begin{center}
\begin{tikzpicture}[scale=0.95, font=\small]
% sol
\draw[thick, popdark] (-0.5,0) -- (10.5,0);
% capteur noir incliné
\fill[popdark!85] (0,0.2) -- (3.2,0.2) -- (3.2,1.1) -- (0,0.5) -- cycle;
\draw[thick, popdark] (0,0.2) -- (3.2,0.2) -- (3.2,1.1) -- (0,0.5) -- cycle;
% chambre de séchage
\draw[thick, popblue] (3.2,0.2) -- (8.2,0.2) -- (8.2,2.6) -- (3.2,2.6) -- cycle;
\draw[thick, popblue!60, fill=popblueL!40] (3.2,2.6) -- (8.2,2.6) -- (7.6,3.1) -- (3.8,3.1) -- cycle;
% claies
\foreach \y in {0.8,1.5,2.2}{
  \draw[poporange, thick] (3.5,\y) -- (7.9,\y);
  \foreach \x in {3.8,4.6,5.4,6.2,7.0,7.8}{\draw[poporange] (\x,\y) -- (\x,\y+0.12);}
}
% rondelles
\foreach \x/\y in {4.0/0.92, 5.0/0.92, 6.0/0.92, 7.0/0.92, 4.5/1.62, 5.5/1.62, 6.5/1.62, 7.5/1.62}{
  \fill[popgold] (\x,\y) ellipse (0.22 and 0.07);
}
% cheminée
\draw[thick, popblue] (7.6,3.1) -- (7.6,4.2) -- (8.2,4.2) -- (8.2,3.1);
\draw[->, thick, popgreen] (7.9,4.2) -- (7.9,5.0) node[right]{air humide};
% soleil
\draw[fill=poppinkL, draw=poppink, thick] (0.4,4.6) circle (0.45);
\foreach \a in {0,45,90,135,180,225,270,315}{
  \draw[poppink, thick] ({0.4+0.6*cos(\a)},{4.6+0.6*sin(\a)}) -- ({0.4+0.95*cos(\a)},{4.6+0.95*sin(\a)});
}
% rayons vers capteur
\draw[->, thick, poppink] (1.0,4.1) -- (1.4,1.0);
\draw[->, thick, poppink] (1.6,4.3) -- (2.4,1.05);
% flux d'air
\draw[->, very thick, popgreen] (-0.6,0.6) -- (0.3,0.6) node[midway, below left=-2pt]{air frais};
\draw[->, very thick, popgreen] (1.2,0.75) -- (3.0,0.9);
\draw[->, very thick, popgreen] (4.0,1.0) -- (7.4,1.0);
\draw[->, very thick, popgreen] (6.0,1.9) -- (7.5,2.6);
% labels
\node[align=center, popdark] at (1.6,-0.4) {capteur solaire\\ (surface noire)};
\node[align=center, popblue] at (5.7,-0.4) {chambre de séchage\\ (claies + rondelles)};
\node[popblue] at (5.7,3.45) {paroi transparente};
\end{tikzpicture}
\end{center}
\legende{Schéma en coupe d'un séchoir solaire de type tunnel. Le rayonnement solaire chauffe le capteur noir ; l'air chaud, plus léger, traverse les claies chargées de rondelles de mangue, se charge de vapeur d'eau et s'échappe par la cheminée, entraînant l'air frais en continu.}
\end{popfigure}

\textbf{Applications.} Les \emph{rondelles de mangue séchées} (tranches de 5 à \SI{8}{mm}, éventuellement blanchies, séchées jusqu'à une texture souple de cuir) et la \emph{tomate séchée} (demi-fruits épépinés, légèrement salés) se conservent plusieurs mois. \textbf{Conditionnement} : dès refroidissement, en sachets ou bocaux \textbf{hermétiques}, à l'abri de l'humidité et de la lumière ; un produit sec qui reprend l'eau de l'air moisit rapidement.

\begin{savaistu}[Pourquoi blanchir avant de sécher ?]
Les rondelles de mangue non blanchies brunissent au séchage : la \emph{polyphénol oxydase}, une enzyme, oxyde les composés phénoliques en pigments bruns en présence de l'oxygène de l'air. Un blanchiment de 2 à 3 min à l'eau bouillante \textbf{dénature cette enzyme} : les rondelles conservent alors leur belle couleur orangée, signe commercial de qualité sur les marchés de Yaoundé et de Douala.
\end{savaistu}

\courssub{La mise en conserve et la stérilisation : conserver par la chaleur en milieu clos}

\textbf{Principe.} On enferme le fruit dans un récipient hermétique (bocal en verre, boîte métallique), puis on le chauffe suffisamment pour détruire les micro-organismes et leurs spores. Le récipient étant étanche, aucune recontamination n'est possible : le produit se conserve \emph{à température ambiante} pendant des mois, voire des années.

\begin{methode}[Les étapes d'une conserve]
\begin{enumerate}
\item \textbf{Préparer} le fruit (tri, lavage, épluchage, découpage) et stériliser les bocaux et couvercles à l'eau bouillante.
\item \textbf{Remplir} les bocaux en tassant légèrement, en laissant 1 à \SI{2}{cm} de vide sous le couvercle.
\item \textbf{Ajouter éventuellement un liquide de couverture} : sirop de sucre (mangues au sirop) ou saumure/purée (coulis de tomate), qui chasse l'air et améliore la conservation.
\item \textbf{Fermer hermétiquement} (joint en bon état, capsule neuve).
\item \textbf{Traiter thermiquement} en autocuiseur ou stérilisateur : environ \SI{100}{\celsius} pour les produits acides (tomate, $pH < 4{,}5$), 110 à \SI{115}{\celsius} pour les produits peu acides, pendant une durée adaptée au volume du bocal (20 à 40 min).
\item \textbf{Refroidir} rapidement, vérifier la fermeture (couvercle bombé vers le bas, « pop » à l'ouverture), \textbf{étiqueter} (contenu, date).
\end{enumerate}
\end{methode}

\begin{attention}[Temps et température : non négociables]
Un bocal mal fermé ou un traitement thermique insuffisant expose au développement de germes dangereux, notamment des bactéries sporulées anaérobies productrices de toxines (risque de \cle{toxi-infection alimentaire} grave). Signes d'alerte à l'ouverture : couvercle bombé, jet de gaz, odeur anormale, liquide trouble $\Rightarrow$ \textbf{jeter sans goûter}. Respecter scrupuleusement les couples temps/température des recettes validées.
\end{attention}

\textbf{Applications.} \emph{Mangues au sirop} : quartiers de mangue fermes recouverts d'un sirop à 30 \% de sucre, stérilisés 25 min à \SI{100}{\celsius}. \emph{Coulis de tomate} : tomates cuites, passées au moulin, réduites, mises en bouteilles ou bocaux et stérilisées 30 min à \SI{100}{\celsius}.

\courssub{La confiture : le sucre comme conservateur}

\textbf{Principe.} La confiture combine \textbf{deux barrières} : le sucre en forte concentration abaisse l'$a_w$ par \emph{osmose} (l'eau est « captée » par le sucre et devient indisponible pour les micro-organismes), et la cuisson détruit la flore et les enzymes. La \cle{pectine}, polysaccharide naturel des fruits, forme en présence de sucre et d'acidité un réseau qui \textbf{gélifie} la préparation au refroidissement.

\begin{methode}[Protocole simplifié de confiture de mangue (club de sciences ou maison)]
\begin{enumerate}
\item Peler et dénoyauter des mangues bien mûres mais saines ; broyer la pulpe et la peser.
\item Ajouter environ \SI{1}{kg} de sucre pour \SI{1}{kg} de pulpe (au minimum \SI{700}{g}), plus le jus d'un citron (apport d'acidité, favorable à la gélification et à la conservation).
\item Cuire à feu moyen en remuant régulièrement ; écumer si nécessaire.
\item \textbf{Test de gélification} : déposer une goutte de confiture sur une assiette froide ; si elle se fige et ne coule pas, la cuisson est terminée (extrait sec $\approx$ 65 \%).
\item Mettre en pots préalablement ébouillantés, remplir à ras bord, fermer immédiatement et retourner les pots 5 min (auto-stérilisation du couvercle), puis laisser refroidir et étiqueter.
\end{enumerate}
\end{methode}

\begin{popfigure}
\begin{center}
\begin{tikzpicture}[font=\small, node distance=0.55cm,
  etape/.style={rectangle, rounded corners=3pt, draw=popblue, fill=popblueL!35, thick, align=center, minimum width=3.1cm, minimum height=0.85cm},
  ctrl/.style={rectangle, rounded corners=3pt, draw=poporange, fill=poporangeL!45, thick, align=center, minimum width=3.1cm, minimum height=0.85cm},
  fleche/.style={->, very thick, popdark}]
\node[etape, align=center] (n1){Mangues fraîches\\ saines et mûres};
\node[etape, below=of n1, align=center] (n2){Tri, lavage,\\ épluchage, dénoyautage};
\node[etape, below=of n2, align=center] (n3){Broyage de la pulpe\\ et pesée};
\node[ctrl, below=of n3, align=center] (n4){Cuisson avec sucre\\ (1 kg/kg) + citron};
\node[ctrl, below=of n4, align=center] (n5){Test de gélification\\ sur assiette froide};
\node[etape, below=of n5, align=center] (n6){Mise en pots\\ ébouillantés, à ras bord};
\node[etape, below=of n6, align=center] (n7){Fermeture, retournement\\ 5 min, refroidissement};
\node[etape, below=of n7, fill=popgreenL!50, draw=popgreen, align=center] (n8){Pot étiqueté :\\ confiture de mangue};
\foreach \a/\b in {n1/n2, n2/n3, n3/n4, n4/n5, n5/n6, n6/n7, n7/n8}{
  \draw[fleche] (\a) -- (\b);
}
\node[right=0.4cm of n4, align=left, text=popmuted] {barrières : chaleur\\ + sucre (osmose)};
\node[right=0.4cm of n5, align=left, text=popmuted] {extrait sec\\ $\approx$ 65 \%};
\end{tikzpicture}
\end{center}
\legende{Organigramme de fabrication de la confiture de mangue, du fruit frais au pot étiqueté. Les étapes encadrées en orange sont les points critiques de maîtrise de la conservation.}
\end{popfigure}

\courssub{Les jus et nectars : une conservation plus courte}

\textbf{Principe.} Le jus est extrait par pressage ou broyage, puis \textbf{filtré}. On y ajoute éventuellement du sucre (nectar de mangue : pulpe diluée + sirop) et de l'\emph{acide citrique} (abaissement du $pH$, barrière supplémentaire). Le jus est ensuite \textbf{pasteurisé} (chauffage de 75 à \SI{85}{\celsius} pendant quelques minutes, qui détruit les germes végétatifs sans cuire le produit) et \textbf{embouteillé à chaud} dans des bouteilles ébouillantées, bouchées immédiatement.

\begin{attention}[Jus $\neq$ confiture]
Un jus pasteurisé contient encore beaucoup d'eau ($a_w$ élevée) : sa conservation est \textbf{plus courte} que celle d'une confiture (quelques semaines à quelques mois, contre plus d'un an). Une fois la bouteille ouverte, le jus doit être conservé au frais et consommé en 2 à 3 jours.
\end{attention}

\courssub{Les pâtes et concentrés : le cas de la pâte de tomate}

\textbf{Principe.} La concentration élimine l'eau par \textbf{évaporation} au cours d'une cuisson prolongée : l'extrait sec augmente, l'$a_w$ chute, et le produit devient stable après conditionnement et stérilisation.

\begin{popfigure}
\begin{center}
\begin{tikzpicture}[font=\small, node distance=0.5cm,
  etape/.style={rectangle, rounded corners=3pt, draw=poppink, fill=poppinkL!40, thick, align=center, minimum width=3.0cm, minimum height=0.8cm},
  fleche/.style={->, very thick, popdark}]
\node[etape] (t1) {Tomates fraîches};
\node[etape, below=of t1] (t2) {Tri et lavage};
\node[etape, below=of t2, align=center] (t3){Broyage et\\ passage au tamis};
\node[etape, below=of t3, align=center] (t4){Cuisson-concentration\\ (évaporation de l'eau)};
\node[etape, below=of t4, align=center] (t5){Conditionnement\\ (bocaux, sachets)};
\node[etape, below=of t5, fill=popgreenL!50, draw=popgreen, align=center] (t6){Stérilisation puis\\ étiquetage : pâte de tomate};
\foreach \a/\b in {t1/t2, t2/t3, t3/t4, t4/t5, t5/t6}{
  \draw[fleche] (\a) -- (\b);
}
\node[right=0.4cm of t4, align=left, text=popmuted] {extrait sec :\\ 6 \% $\rightarrow$ 28-30 \%};
\end{tikzpicture}
\end{center}
\legende{Organigramme de fabrication de la pâte de tomate. L'étape clé est la cuisson-concentration, qui porte l'extrait sec de 6 \% à environ 28-30 \%.}
\end{popfigure}

\begin{exemplebox}[Un calcul de concentration éclairant]
La tomate fraîche contient environ 6 \% d'extrait sec ; la pâte double concentré en contient 28 à 30 \%. Pour obtenir \SI{1}{kg} de pâte à 30 \%, il faut environ \SI{5}{kg} de tomates fraîches : on évapore donc près de \textbf{4/5 de l'eau} initiale. C'est cette énorme dépense d'énergie (bois, gaz) qui explique le prix des concentrés — et l'intérêt des séchoirs solaires comme source de chaleur gratuite.
\end{exemplebox}

\begin{exoresolu}[Rendement d'une unité artisanale de pâte de tomate]
Une coopérative de femmes de Bafoussam transforme \SI{250}{kg} de tomates fraîches à 6 \% d'extrait sec en pâte à 30 \% d'extrait sec. Quelle masse de pâte obtient-elle ? Quelle masse d'eau a été évaporée ?
\tcblower
La matière sèche se conserve au cours de l'évaporation (seule l'eau s'élimine) :
\[
m_{\text{sec}} = 250 \times \frac{6}{100} = \SI{15}{kg}.
\]
Cette matière sèche représente 30 \% de la pâte finale :
\[
m_{\text{pâte}} = \frac{15}{0{,}30} = \SI{50}{kg}.
\]
La masse d'eau évaporée est donc :
\[
m_{\text{eau}} = 250 - 50 = \SI{200}{kg}.
\]
\textbf{Conclusion} : la coopérative obtient \SI{50}{kg} de pâte et a évaporé \SI{200}{kg} d'eau, soit exactement 4/5 de la masse initiale.
\end{exoresolu}

\courssub{Autres voies de transformation (compléments)}

\begin{itemize}
\item \textbf{La fermentation} : des levures transforment les sucres en éthanol (vin de mangue, de goyave) ; des bactéries acétiques peuvent ensuite oxyder l'éthanol en acide acétique (vinaigre). L'alcool et l'acidité constituent des barrières microbiennes efficaces.
\item \textbf{La congélation} : elle bloque l'activité microbienne et enzymatique sans détruire les germes. Elle conserve très bien la qualité nutritionnelle, mais exige une \textbf{chaîne du froid} ininterrompue ($\leq -\SI{18}{\celsius}$) : au Cameroun, les coupures d'électricité et le coût de l'énergie limitent fortement son usage domestique et artisanal.
\end{itemize}

\begin{aretenir}[L'essentiel de la section]
Toutes les techniques reposent sur les mêmes principes, combinés différemment : le \textbf{séchage} retire l'eau ($a_w$ basse) ; la \textbf{conserve} détruit les germes par la chaleur en milieu hermétique ; la \textbf{confiture} associe sucre (osmose) et cuisson ; les \textbf{jus} sont pasteurisés mais restent fragiles ; les \textbf{pâtes} concentrent par évaporation. Dans tous les cas, l'hygiène préalable, le respect des couples temps/température et l'emballage hermétique conditionnent le succès.
\end{aretenir}

\coursec{V. Comparaison de l'efficacité des techniques de conservation}

Nous venons de décrire, dans la section précédente, les principales techniques de transformation et de conservation appliquées à la mangue et à la tomate : séchage, mise en conserve avec stérilisation, transformation en confiture, en jus ou en pâte. Une question pratique se pose alors immédiatement au producteur, à la ménagère ou à l'entrepreneur camerounais : \emph{quelle technique choisir ?} Faut-il sécher les mangues excédentaires de la saison de mars à juin, les mettre en conserve, ou en faire de la confiture ? Pour répondre de manière rigoureuse, il ne suffit pas de se fier aux habitudes : il faut \cle{comparer} les procédés selon des \cle{critères objectifs} et mesurables. C'est l'objet de cette section, qui mobilise la démarche d'investigation scientifique : formuler des critères, recueillir des données, les exploiter dans des tableaux et des graphiques, puis conclure en connaissance de cause.

\courssub{Les critères de comparaison des techniques de conservation}

\begin{definition}[Critères d'évaluation d'un procédé de conservation]
L'efficacité d'une technique de conservation s'apprécie selon plusieurs \cle{critères} complémentaires :
\begin{itemize}
\item la \cle{durée de conservation} : temps pendant lequel le produit reste comestible et sain ;
\item la \cle{qualité nutritionnelle} : conservation des vitamines (notamment la vitamine C), des sels minéraux, des pigments (caroténoïdes) et des arômes ;
\item la \cle{qualité organoleptique} : goût, couleur, odeur, texture du produit fini, appréciés par les sens du consommateur ;
\item le \cle{coût} et l'\cle{énergie requise} : investissement en matériel, combustible ou électricité, main-d'œuvre ;
\item l'\cle{accessibilité technologique} : simplicité de mise en œuvre, disponibilité locale du matériel et des compétences.
\end{itemize}
\end{definition}

\begin{attention}[Durée de conservation et valeur nutritionnelle ne vont pas de pair]
Une erreur fréquente consiste à croire que « plus ça se conserve longtemps, meilleur c'est ». C'est faux : une longue durée de conservation ne signifie \textbf{pas} une meilleure valeur nutritionnelle. La confiture de mangue se conserve plus d'un an, mais sa cuisson prolongée détruit une grande partie de la vitamine C, et l'ajout massif de sucre en fait un aliment très énergétique. Inversement, le fruit frais, champion nutritionnel, ne se conserve que quelques jours. Il faut donc toujours préciser \emph{quel critère} on privilégie avant de juger une technique.
\end{attention}

\courssub{Document 1 : durées indicatives de conservation}

Le document suivant rassemble des durées de conservation admises à titre documentaire, dans des conditions correctes d'hygiène et de stockage (à l'abri de la chaleur, de l'humidité et de la lumière).

\begin{center}
\begin{tabularx}{\linewidth}{|l|X|X|}
\hline
\rowcolor{popblueL} \textbf{Mode de conservation} & \textbf{Durée indicative} & \textbf{Principe limitant l'altération} \\
\hline
Fruit frais à l'air ambiant & quelques jours (2 à 7 jours) & aucun : micro-organismes et enzymes actifs \\
\hline
Fruit réfrigéré (4 à \SI{8}{\celsius}) & 1 à 2 semaines & le froid ralentit métabolismes et croissance microbienne \\
\hline
Fruit séché & plusieurs mois (6 à 12 mois) & eau insuffisante pour les micro-organismes \\
\hline
Confiture & 1 an et plus & forte concentration en sucre + cuisson \\
\hline
Conserve stérilisée & 1 an et plus (2 à 3 ans) & stérilisation + étanchéité (anaérobiose) \\
\hline
\end{tabularx}
\end{center}

On observe une hiérarchie claire : plus le traitement est intense (dessiccation, sucre, chaleur + étanchéité), plus la durée de conservation s'allonge. Mais ce document ne dit rien de la \emph{qualité} du produit conservé. Il faut un second document, portant sur la teneur en vitamine C.

\courssub{Document 2 : la vitamine C, témoin de la qualité nutritionnelle}

\begin{propriete}[Thermosensibilité et oxydabilité de la vitamine C]
La \cle{vitamine C} (acide ascorbique) est une vitamine \cle{hydrosoluble}, \cle{thermolabile} (détruite par la chaleur) et facilement \cle{oxydable} (détruite au contact du dioxygène de l'air, surtout en présence de lumière et de chaleur). Sa teneur dans un produit transformé est donc un excellent \cle{indicateur} de l'agressivité du procédé de conservation : plus un traitement combine chaleur, durée et exposition à l'air, plus la perte en vitamine C est importante.
\end{propriete}

Des dosages réalisés sur un même lot de mangues (variété locale, stade de maturité identique) soumises à différents traitements donnent les résultats suivants (valeurs moyennes, exprimées en milligrammes de vitamine C pour \SI{100}{g} de produit consommé) :

\begin{center}
\begin{tabularx}{\linewidth}{|l|X|X|}
\hline
\rowcolor{popblueL} \textbf{Produit} & \textbf{Vitamine C (mg/100 g)} & \textbf{Cause principale de la perte} \\
\hline
Mangue fraîche & 28 & --- (référence) \\
\hline
Mangue séchée & 12 & chaleur modérée prolongée + oxygène de l'air \\
\hline
Mangue en conserve (sirop) & 15 & chaleur de stérilisation, mais milieu peu oxygéné \\
\hline
Confiture de mangue & 8 & cuisson prolongée à l'air libre \\
\hline
\end{tabularx}
\end{center}

\begin{popfigure}
\begin{center}
\begin{tikzpicture}
\begin{axis}[
    ybar,
    bar width=0.55cm,
    width=0.85\linewidth,
    height=6.5cm,
    ymin=0, ymax=32,
    ylabel={Vitamine C (mg/100 g)},
    symbolic x coords={Fraîche, Séchée, Conserve, Confiture},
    xtick=data,
    nodes near coords,
    nodes near coords align={vertical},
    every node near coord/.append style={font=\small},
    ymajorgrids=true,
    grid style={popline!40},
    axis line style={popdark},
]
\addplot[fill=popgreen!70, draw=popgreen] coordinates {(Fraîche,28)};
\addplot[fill=poporange!70, draw=poporange] coordinates {(Séchée,12)};
\addplot[fill=popblue!50, draw=popblue] coordinates {(Conserve,15)};
\addplot[fill=poppink!70, draw=poppink] coordinates {(Confiture,8)};
\end{axis}
\end{tikzpicture}
\end{center}
\legende{Teneur en vitamine C de la mangue selon le mode de conservation (mg pour \SI{100}{g} de produit). Le fruit frais sert de référence ; les traitements thermiques (séchage, stérilisation, cuisson en confiture) détruisent une fraction variable de la vitamine, la cuisson prolongée à l'air libre étant la plus destructive.}
\end{popfigure}

\begin{exemplebox}[La vitamine C réduite de plus de moitié par la cuisson prolongée]
Dans l'exemple chiffré ci-dessus, la mangue fraîche contient \SI{28}{mg} de vitamine C pour \SI{100}{g}. Après transformation en confiture (cuisson de 30 à 60 minutes à ébullition, à l'air libre, avec ajout de sucre), il n'en reste que \SI{8}{mg} pour \SI{100}{g}, soit une perte de :
\[ \frac{28 - 8}{28} \times 100 \approx 71\,\% \]
La vitamine C, thermolabile et oxydable, est donc réduite de \textbf{plus de moitié} (ici plus des deux tiers) lors d'une cuisson prolongée. La confiture reste appréciable pour son goût, son énergie (sucres) et sa longue conservation, mais elle ne peut plus être considérée comme une source de vitamine C.
\end{exemplebox}

\textbf{Interprétation.} Le classement des pertes (confiture $>$ séchage $>$ conserve) s'explique par la combinaison de deux facteurs : la \emph{température} et la \emph{durée} du chauffage d'une part, l'\emph{exposition à l'oxygène} d'autre part. La conserve subit une chaleur intense mais brève (stérilisation) et surtout en milieu privé d'air une fois le bocal fermé, ce qui limite l'oxydation ; le séchage solaire expose le fruit longtemps à l'air chaud ; la confiture cumule cuisson longue et contact avec l'air.

\courssub{Étude de cas : un même lot de tomates soumis à trois techniques}

\begin{experience}[Suivi comparatif de l'altération de tomates conservées]
\textbf{Problème.} Une coopérative de maraîchers de la région de l'Ouest récolte en pleine saison un excédent de tomates qu'elle ne peut écouler. Quelle technique limite le mieux les pertes ?

\textbf{Hypothèse.} Les trois techniques (séchage, conserve stérilisée, pâte concentrée) diffèrent par leur efficacité à retarder l'altération.

\textbf{Protocole.} Un même lot de tomates mûres et saines est réparti en trois échantillons de 40 unités équivalentes :
\begin{itemize}
\item échantillon A : tomates coupées et séchées, conservées en sachets fermés ;
\item échantillon B : tomates mises en bocaux et stérilisées \SI{100}{\celsius} pendant \SI{90}{min} ;
\item échantillon C : tomates cuites et concentrées en pâte, conditionnées en pots fermés.
\end{itemize}
Chaque mois, on examine les échantillons et on note la proportion d'unités altérées (moisissures, fermentation, odeur anormale, bombeage des bocaux).

\textbf{Résultats} (pourcentage cumulé d'échantillons altérés) :
\begin{center}
\begin{tabular}{|c|c|c|c|}
\hline
\rowcolor{popblueL} \textbf{Temps (mois)} & \textbf{A : séchage} & \textbf{B : conserve} & \textbf{C : pâte} \\
\hline
1 & $0\,\%$ & $0\,\%$ & $0\,\%$ \\
\hline
3 & $5\,\%$ & $0\,\%$ & $2\,\%$ \\
\hline
6 & $15\,\%$ & $2\,\%$ & $8\,\%$ \\
\hline
9 & $30\,\%$ & $5\,\%$ & $15\,\%$ \\
\hline
12 & $45\,\%$ & $8\,\%$ & $25\,\%$ \\
\hline
\end{tabular}
\end{center}
\end{experience}

\begin{popfigure}
\begin{center}
\begin{tikzpicture}
\begin{axis}[
    width=0.85\linewidth,
    height=7cm,
    xlabel={Temps (mois)},
    ylabel={Échantillons altérés (\%)},
    xmin=0, xmax=13,
    ymin=0, ymax=55,
    xtick={0,3,6,9,12},
    ytick={0,10,20,30,40,50},
    legend pos=north west,
    legend style={font=\small},
    ymajorgrids=true,
    grid style={popline!40},
    axis line style={popdark},
]
\addplot[color=poporange, mark=*, thick] coordinates {(1,0) (3,5) (6,15) (9,30) (12,45)};
\addplot[color=popblue, mark=square*, thick] coordinates {(1,0) (3,0) (6,2) (9,5) (12,8)};
\addplot[color=popgreen, mark=triangle*, thick] coordinates {(1,0) (3,2) (6,8) (9,15) (12,25)};
\legend{Séchage, Conserve stérilisée, Pâte concentrée}
\end{axis}
\end{tikzpicture}
\end{center}
\legende{Évolution de la proportion d'échantillons de tomates altérés au cours du temps pour trois techniques de conservation. La conserve stérilisée (courbe bleue) est la plus efficace à long terme ; le séchage (courbe orange) est le plus sensible aux reprises d'humidité et aux contaminations.}
\end{popfigure}

\textbf{Interprétation.} Après 12 mois, la conserve stérilisée n'enregistre que $8\,\%$ d'altérations, contre $25\,\%$ pour la pâte et $45\,\%$ pour le séchage. La stérilisation détruit les micro-organismes et leurs spores, et l'étanchéité du bocal empêche toute recontamination : c'est la technique la plus sûre à long terme. La pâte, moins stérilisée et riche en eau résiduelle, s'altère davantage une fois les pots manipulés. Le séchage, efficace tant que le produit reste sec, est vulnérable à l'humidité ambiante (saison des pluies) : les tomates séchées reprennent de l'eau, ce qui permet la reprise de la croissance microbienne (moisissures).

\textbf{Conclusion.} L'hypothèse est validée : à durée de stockage égale, la conserve stérilisée protège le mieux, mais elle exige matériel (bocaux, autocuiseur ou stérilisateur) et énergie. Le séchage, moins performant sur un an, reste le plus accessible. Le choix dépend donc du contexte.

\courssub{Construire et exploiter un tableau comparatif multicritères}

\begin{methode}[Justifier le choix d'un procédé par une analyse multicritères]
\begin{enumerate}
\item \textbf{Lister les critères} pertinents pour la situation : durée de conservation, qualité nutritionnelle, qualité organoleptique, coût, énergie, accessibilité.
\item \textbf{Attribuer à chaque technique} une évaluation par critère, en s'appuyant sur des données (durées mesurées, dosages de vitamines, coûts réels) et non sur des impressions.
\item \textbf{Identifier le ou les critères prioritaires} imposés par la situation (ex. : pas d'électricité au village $\Rightarrow$ éliminer les procédés énergivores ; vente lointaine $\Rightarrow$ privilégier la longue conservation).
\item \textbf{Conclure en justifiant} : citer explicitement les données qui motivent le choix et reconnaître les inconvénients acceptés.
\end{enumerate}
\end{methode}

Appliquons cette méthode à la mangue :

\begin{center}
\begin{tabularx}{\linewidth}{|l|X|X|X|X|}
\hline
\rowcolor{popblueL} \textbf{Critère} & \textbf{Séchage} & \textbf{Conserve} & \textbf{Confiture} & \textbf{Jus} \\
\hline
Durée de conservation & plusieurs mois & 1 an et plus & 1 an et plus & courte (quelques jours) \\
\hline
Vitamine C préservée & moyenne & moyenne à bonne & faible & bonne si consommé vite \\
\hline
Qualité organoleptique & texture ferme, goût concentré & proche du frais & très sucrée & proche du frais \\
\hline
Coût et énergie & faibles (soleil) & élevés (bocaux, stérilisation) & moyens (sucre, cuisson) & moyens \\
\hline
Accessibilité & très accessible & technologie requise & accessible & accessible \\
\hline
\end{tabularx}
\end{center}

\begin{exoresolu}[Choisir un procédé pour une coopérative villageoise]
\textbf{Énoncé.} Une coopérative d'un village de l'Adamaoua, sans électricité, récolte chaque année un excédent de mangues qu'elle vend en partie sur les marchés urbains. Elle souhaite réduire ses pertes et disposer d'un produit vendable toute l'année. En exploitant les documents de cette section, proposer et justifier un choix de technique de conservation.
\tcblower
\textbf{Solution détaillée.}
\begin{enumerate}
\item \emph{Critères prioritaires.} L'absence d'électricité élimine les procédés très énergivores ; la vente toute l'année exige une durée de conservation de plusieurs mois au moins ; la coopérative dispose de peu de capital.
\item \emph{Analyse.} La conserve stérilisée offre la meilleure durée (1 an et plus) et préserve correctement la vitamine C (\SI{15}{mg/100 g}), mais exige bocaux, stérilisateur et combustible : coût élevé et accessibilité faible. La confiture se conserve longtemps mais perd plus de $70\,\%$ de sa vitamine C et demande beaucoup de sucre (coût). Le séchage solaire conserve plusieurs mois, garde une partie de la vitamine C (\SI{12}{mg/100 g}), concentre les arômes, et ne coûte presque rien en énergie puisque l'Adamaoua bénéficie d'une saison sèche ensoleillée.
\item \emph{Conclusion justifiée.} Le \textbf{séchage solaire} est le choix le plus adapté : il répond au critère prioritaire (pas d'électricité), assure plusieurs mois de conservation et un produit transportable vers les marchés urbains. On acceptera ses limites : vulnérabilité à l'humidité (stockage en sachets hermétiques) et perte partielle de vitamine C. À terme, si la coopérative dégage des bénéfices, elle pourra investir dans la conserve pour viser une conservation supérieure à un an.
\end{enumerate}
\end{exoresolu}

\courssub{Complémentarité des techniques et limites des procédés}

\begin{aretenir}[Aucune technique n'est « meilleure » en absolu]
Les techniques de conservation sont \cle{complémentaires}, pas rivales. Le choix dépend :
\begin{itemize}
\item du \textbf{fruit} (la tomate, très aqueuse, se prête bien à la pâte et à la conserve ; la mangue, parfumée et sucrée, au séchage et à la confiture) ;
\item des \textbf{moyens disponibles} (énergie, matériel, compétences, capital) ;
\item de l'\textbf{usage visé} (consommation familiale rapide, vente locale, exportation, sécurisation alimentaire hors saison).
\end{itemize}
Une stratégie efficace combine souvent plusieurs procédés sur un même excédent de récolte.
\end{aretenir}

Chaque procédé comporte aussi des \cle{limites} qu'il faut connaître pour un usage raisonné :

\begin{itemize}
\item \textbf{Pertes vitaminiques} : tous les traitements thermiques détruisent une partie de la vitamine C (jusqu'à plus de $70\,\%$ pour la confiture) ; le produit conservé ne remplace jamais totalement le fruit frais sur le plan nutritionnel.
\item \textbf{Ajout de sucre} : les confitures contiennent souvent autant de sucre ajouté que de fruit ; leur consommation excessive favorise caries, obésité et diabète, un enjeu de santé publique croissant dans les villes camerounaises.
\item \textbf{Coût énergétique} : la stérilisation des conserves exige une chaleur intense et prolongée (bois, gaz ou électricité), ce qui pèse sur le budget et, si le bois de chauffe est mal géré, sur les ressources forestières.
\item \textbf{Dépendance au climat et à l'hygiène} : le séchage solaire devient difficile en saison des pluies et expose le produit aux poussières, insectes et moisissures si les règles d'hygiène sont négligées.
\end{itemize}

\begin{savaistu}[Le séchage solaire : le moins coûteux, mais pas sans risques]
Le \cle{séchage solaire} est la technique de conservation la moins coûteuse : il ne demande ni combustible ni électricité, seulement le soleil, abondant sous nos latitudes. C'est pourquoi il est pratiqué depuis des générations au Cameroun (mangues, tomates, piments, feuilles de légumes). Mais il présente deux faiblesses : il \textbf{dépend du climat} (impossible ou très lent pendant les longues saisons des pluies du Sud forestier) et il \textbf{expose aux contaminations} si l'hygiène est négligée — poussières, mouches, moisissures productrices de toxines dangereuses (aflatoxines). Les séchoirs solaires améliorés (caisses vitrées, claies surélevées couvertes de moustiquaire) corrigent une grande partie de ces défauts pour un coût modeste : un bel exemple de biotechnologie appropriée au contexte local.
\end{savaistu}

\begin{attention}[Pièges de rédaction au Baccalauréat]
\begin{itemize}
\item Ne jamais écrire « la conserve est la meilleure technique » sans préciser \emph{pour quel critère} : elle est la meilleure pour la durée et la sécurité sanitaire, pas pour le coût ni l'accessibilité.
\item Ne pas confondre \emph{durée de conservation} et \emph{valeur nutritionnelle} : les confitures se conservent très longtemps mais sont pauvres en vitamine C et riches en sucre.
\item Dans l'exploitation d'un graphique, toujours citer des \textbf{valeurs} (ex. : « à 12 mois, $45\,\%$ d'altérations pour le séchage contre $8\,\%$ pour la conserve ») avant de conclure.
\end{itemize}
\end{attention}

En définitive, comparer les techniques de conservation, c'est mettre la science au service d'une décision concrète : réduire les pertes post-récolte tout en préservant au mieux la qualité des aliments. Cette comparaison éclaire naturellement la dernière étape de notre leçon : l'évaluation des intérêts économiques, nutritionnels et environnementaux de la transformation des fruits de saison.

\coursec{VI. Intérêts économiques, nutritionnels et environnementaux}

Après avoir comparé l'efficacité technique des différentes méthodes de conservation (séchage, stérilisation, confiturage) sur la qualité microbiologique et organoleptique des fruits, il est légitime de s'interroger sur la portée réelle de ces biotechnologies pour nos sociétés. La conservation n'est pas une fin en soi : elle répond à une \textbf{famille de situations} centrale du programme, les \emph{problèmes liés à la couverture des besoins alimentaires et thérapeutiques}. Au Cameroun comme dans de nombreux pays tropicaux, la saisonnalité extrême de la production (mangue, tomate, ananas, papaye) crée un paradoxe : l'abondance gaspillée en pleine saison côtoie la pénurie onéreuse en hors-saison. Cette section démontre comment la maîtrise des facteurs d'altération (micro-organismes, enzymes, oxydation) se traduit en bénéfices concrets, mesurables et durables pour la nutrition, l'économie, l'environnement et le tissu social.

\courssub{Intérêt nutritionnel : sécurité alimentaire et lutte contre les carences}

\paragraph{Observation et constat.}
Sur les marchés de Douala, Yaoundé ou Bafoussam, le prix d'un panier de mangues varie du simple au décuple entre la pleine saison (avril-juin) et la saison sèche (décembre-février). En pleine saison, les fruits pourrissent au pied des manguiers ou sur les étals par excès d'offre ; en hors-saison, ils deviennent un produit de luxe inaccessible aux ménages modestes. Or, la mangue est une source majeure de \cle{provitamine A} (\beta-carotène), de vitamine C et de fibres alimentaires. La tomate, quant à elle, apporte du \cle{lycopène} (antioxydant puissant) et de la vitamine C.

\paragraph{Démarche d'investigation : impact de la conservation sur la densité nutritionnelle.}
Considérons une expérience comparative menée au Centre de Recherche sur les Aliments et la Nutrition (CRAN) de l'Université de Ngaoundéré. On suit la teneur en vitamine C (mg/100 g) de mangues fraîches stockées à température ambiante (\SI{28}{\celsius}) versus des mangues séchées (solaire, \SI{55}{\celsius}) et des mangues en conserve (stérilisées \SI{100}{\celsius}, 20 min) sur 90 jours.

\begin{popfigure}
\begin{tikzpicture}
\begin{axis}[
    width=\linewidth, height=8cm,
    xlabel={Durée de stockage (jours)}, ylabel={Vitamine C (mg/100 g)},
    xmin=0, xmax=95, ymin=0, ymax=50,
    xtick={0,15,30,45,60,75,90}, ytick={0,10,20,30,40,50},
    legend style={at={(0.5,-0.15)},anchor=north,legend columns=3},
    grid=major, grid style={popline},
]
\addplot[popcurve, mark=*, mark size=2pt, popblue] coordinates {
    (0,36) (5,28) (10,15) (15,5) (20,0)
};
\addlegendentry{Fraîche (ambiante)}

\addplot[popcurve, mark=square*, mark size=2pt, poporange] coordinates {
    (0,30) (15,28) (30,26) (45,24) (60,22) (75,20) (90,18)
};
\addlegendentry{Séchée (solaire, 12\% H$_2$O)}

\addplot[popcurve, mark=triangle*, mark size=2pt, popgreen] coordinates {
    (0,25) (15,24) (30,23) (45,22) (60,21) (75,20) (90,19)
};
\addlegendentry{Conserve (sirop, stérilisée)}
\end{axis}
\end{tikzpicture}
\legende{Évolution de la teneur en vitamine C de la mangue selon le mode de conservation (données inspirées de travaux CRAN/IRAD). La mangue fraîche perd toute sa vitamine C en 20 jours ; le séchage et la stérilisation stabilisent le taux sur 3 mois.}
\end{popfigure}

\paragraph{Interprétation.}
La courbe montre que la mangue fraîche, non transformée, devient nutritionnellement « vide » en vitamine C au bout de 2 à 3 semaines à cause de l'oxydation enzymatique et microbienne. Le séchage (abaissement de l'\cle{activité de l'eau} $a_w < 0,6$) et la stérilisation (destruction enzymes + micro-organismes, exclusion air) \textbf{figent le capital nutritionnel}. La perte initiale lors du traitement thermique (blanchiment, stérilisation) est compensée par la stabilité à long terme.

\begin{aretenir}[Disponibilité nutritionnelle annuelle]
La transformation et la conservation des fruits de saison permettent de :
\begin{itemize}
    \item \textbf{Lisser l'apport en micronutriments} (vitamines A, C, B9, minéraux, fibres) sur 12 mois, brisant le cycle « abondance/carence » saisonnier.
    \item \textbf{Prévenir les carences spécifiques} : xérophtalmie (vitamine A), scorbut (vitamine C), anémie (folates), notamment chez l'enfant de 6 à 59 mois et la femme enceinte.
    \item \textbf{Diversifier l'alimentation} : jus de tomate, confiture de mangue, pâte de tomate, mangue séchée enrichissent le régime à base de tubercules/céréales (manioc, maïs, riz) souvent pauvres en micronutriments.
\end{itemize}
\end{aretenir}

\begin{attention}[Piège fréquent : « La conserve n'a plus de vitamines »]
\emph{Faux.} Si le traitement thermique détruit une partie des vitamines thermosensibles (vitamine C : -20 à -40\%, folates), il \textbf{stoppe la dégradation continue} qui aurait lieu sur fruit frais. Au bout de 2 mois, la conserve en contient \textbf{davantage} que le fruit frais pourri. De plus, le lycopène de la tomate devient \emph{plus biodisponible} après cuisson (cis-isomérisation). Au Bac, nuancez : « Perte initiale limitée, stabilité à long terme, biodisponibilité parfois accrue. »
\end{attention}

\courssub{Intérêt économique : création de valeur, revenus et stabilisation des prix}

\paragraph{Le gaspillage post-récolte : un manque à gagner massif.}
Au Cameroun, les pertes post-récolte pour la mangue et la tomate sont estimées entre \textbf{30\% et 50\%} de la production (rapports MINADER/FAO). Pour une production nationale de mangue d'environ 200 000 tonnes/an, cela représente 60 000 à 100 000 tonnes perdues, soit des milliards de FCFA brûlés. La transformation convertit ce \textbf{déchet économique} en \textbf{produit à valeur ajoutée}.

\begin{exemplebox}[Exemple chiffré : la valeur ajoutée de la mangue séchée]
Prenons le cas concret d'une coopérative de la région de l'Adamaoua (zone de Ngaoundéré) en pleine saison des mangues (variété Kent/Keitt).
\begin{itemize}
    \item \textbf{Matière première :} 100 kg de mangues fraîches achetées au producteur à \textbf{100 FCFA/kg} $\rightarrow$ Coût = \textbf{10 000 FCFA}.
    \item \textbf{Rendement séchage :} 100 kg frais $\xrightarrow{\text{épluchage, découpe, séchage solaire 2-3 j}}$ $\approx$ 15 à 18 kg de mangue séchée (humidité finale 12-15\%).
    \item \textbf{Conditionnement :} Sachets 100 g (coût emballage + étiquetage $\approx$ 50 FCFA/sachet).
    \item \textbf{Prix de vente local (supermarchés, boutiques bio) :} \textbf{1 500 à 2 000 FCFA / 100 g}.
    \item \textbf{Chiffre d'affaires brut :} 160 sachets $\times$ 1 750 FCFA = \textbf{280 000 FCFA}.
    \item \textbf{Marge brute :} 280 000 - (10 000 + 8 000 emballage + 20 000 main d'œuvre/énergie/logistique) $\approx$ \textbf{242 000 FCFA}.
\end{itemize}
\textbf{Conclusion :} Le kilogramme de mangue fraîche, vendu 100 FCFA, génère \textbf{environ 1 750 FCFA/kg équivalent produit fini} après transformation. La valeur est multipliée par \textbf{17}. C'est le moteur économique de la biotechnologie de conservation.
\end{exemplebox}

\paragraph{Structuration de la filière et création d'emplois.}
Cette valeur ajoutée ne profite pas qu'au transformateur. Elle irrigue toute la \textbf{chaîne de valeur} :

\begin{popfigure}
\begin{tikzpicture}[
    node distance=1.8cm and 2.5cm,
    box/.style={rectangle, draw=popdark, thick, fill=popcream, rounded corners, minimum width=2.8cm, minimum height=1.2cm, align=center, font=\small},
    arrow/.style={->, >=Stealth, thick, popblue},
    val/.style={rectangle, draw=popgreen, fill=popgreenL, rounded corners, minimum width=2.2cm, minimum height=0.8cm, align=center, font=\scriptsize, text=popdark}
]
% Nœuds principaux
\node[box, align=center] (prod){PRODUCTEUR\\ (Cueillette, tri, transport)};
\node[box, right=of prod, align=center] (transfo){UNITÉ DE TRANSFORMATION\\ (GIC, PME, Coopérative)\\ Séchage / Stérilisation / Confiture};
\node[box, right=of transfo, align=center] (distrib){DISTRIBUTION\\ (Grossistes, Supermarchés,\\ Export, Marchés locaux)};
\node[box, right=of distrib, align=center] (conso){CONSOMMATEUR\\ (Ménages, Restaurants,\\ Industries agroalimentaires)};

% Flèches principales
\draw[arrow] (prod) -- node[above, font=\small] {Vente matière première} (transfo);
\draw[arrow] (transfo) -- node[above, font=\small] {Produits finis conditionnés} (distrib);
\draw[arrow] (distrib) -- node[above, font=\small] {Mise en marché} (conso);

% Bulles de valeur ajoutée
\node[val, above=0.8cm of prod, align=center] (v1){\textbf{+ Valeur 1}\\ Revenus producteurs\\ Prix kg stabilisé};
\node[val, above=0.8cm of transfo, align=center] (v2){\textbf{+ Valeur 2}\\ Emplois transformation\\ Technicité, hygiène,\\ Conditionnement};
\node[val, above=0.8cm of distrib, align=center] (v3){\textbf{+ Valeur 3}\\ Marge commerciale\\ Logistique, froid,\\ Accès marchés urbains};
\node[val, above=0.8cm of conso, align=center] (v4){\textbf{+ Valeur 4}\\ Disponibilité produit\\ Qualité, sécurité,\\ Prix abordable hors-saison};

\draw[dashed, popmuted] (prod) -- (v1);
\draw[dashed, popmuted] (transfo) -- (v2);
\draw[dashed, popmuted] (distrib) -- (v3);
\draw[dashed, popmuted] (conso) -- (v4);

% Feedback loop
\draw[arrow, poporange, bend left=40] (conso) to node[above, font=\small, poporange] {Demande, contrats, prix plancher} (prod);
\end{tikzpicture}
\legende{Chaîne de valeur de la mangue transformée au Cameroun. Chaque maillon crée de la valeur (bulles vertes). La flèche orange représente la contractualisation (ex: agriculture contractuelle) qui sécurise le producteur.}
\end{popfigure}

\paragraph{Stabilisation des prix : lissage de l'offre.}
Sans transformation, l'offre est une \textbf{fonction impulsionnelle} : pic massif en 2 mois, quasi-nul le reste de l'année. Le prix suit une hyperbole inverse. La transformation (séchage, conserve) permet de \textbf{stocker la biomasse} sous forme stable et de la réinjecter sur le marché régulièrement.
\begin{itemize}
    \item \textbf{Producteur :} Prix plancher garanti en pleine saison (débouché transformation) $\rightarrow$ incitation à produire plus/qualité.
    \item \textbf{Consommateur :} Prix plafonné en hors-saison (disponibilité conserve/séché) $\rightarrow$ accès nutritionnel maintenu.
    \item \textbf{État :} Moins d'importations de concentré de tomate ou de jus de fruits (économie de devises).
\end{itemize}

\begin{savaistu}[Exportation : la mangue séchée camerounaise en Europe]
Des coopératives comme \textbf{COOPROMA} (Adamaoua) ou \textbf{COOPAF} (Ouest), appuyées par des ONG (SNV, GIZ) et l'IRAD, ont obtenu la certification \textbf{Bio (Ecocert)} et \textbf{Commerce Équitable}. Elles exportent plusieurs tonnes/an de mangue séchée (variétés Brooks, Kent) vers la France, l'Allemagne, la Belgique. Le prix FOB (départ Douala) atteint \textbf{3 500 - 4 500 FCFA/kg} pour la mangue séchée bio, contre 100-150 FCFA/kg pour la fraîche. C'est la preuve que la biotechnologie locale (séchage solaire amélioré, hygiène HACCP) répond aux standards internationaux les plus exigeants.
\end{savaistu}

\paragraph{Visualisation de l'impact d'une unité de transformation sur le devenir de la récolte.}
Le diagramme ci-dessous compare le sort d'une récolte type (100 tonnes de mangues) dans une zone \textbf{sans} unité de transformation vs \textbf{avec} une unité opérationnelle (capacité 5 t/jour, 60 jours/saison).

\begin{popfigure}
\begin{tikzpicture}
\begin{scope}
\begin{axis}[
    width=0.48\linewidth, height=7cm,
    title={SANS unité de transformation},
    title style={font=\small\bfseries, yshift=0.3cm},
    ybar, bar width=25pt,
    ymin=0, ymax=100,
    ylabel={Part de la récolte (\%)},
    symbolic x coords={Fraîche,Pertes,Transformée},
    xtick=data,
    x tick label style={font=\tiny, align=center},
    nodes near coords,
    every node near coord/.append style={font=\small},
]
\addplot[fill=popgreen] coordinates {(Fraîche,40)};
\addplot[fill=poporange] coordinates {(Pertes,50)};
\addplot[fill=popblue] coordinates {(Transformée,10)};
\end{axis}
\end{scope}
\begin{scope}[xshift=0.52\linewidth]
\begin{axis}[
    width=0.48\linewidth, height=7cm,
    title={AVEC unité de transformation},
    title style={font=\small\bfseries, yshift=0.3cm},
    ybar, bar width=25pt,
    ymin=0, ymax=100,
    ylabel={Part de la récolte (\%)},
    symbolic x coords={Fraîche,Pertes,Transformée},
    xtick=data,
    x tick label style={font=\tiny, align=center},
    nodes near coords,
    every node near coord/.append style={font=\small},
]
\addplot[fill=popgreen] coordinates {(Fraîche,30)};
\addplot[fill=poporange] coordinates {(Pertes,10)};
\addplot[fill=popblue] coordinates {(Transformée,60)};
\end{axis}
\end{scope}
\end{tikzpicture}
\legende{Devenir d'une récolte de 100 t de mangues. L'unité de transformation détourne 50 t vers la valeur ajoutée (bleu) et réduit les pertes (orange) de 50 t à 10 t. Données modélisées d'après études MINADER/FAO 2020.}
\end{popfigure}

\courssub{Intérêt environnemental : réduction du gaspillage et économie circulaire}

\paragraph{Le fruit « laid » ou « trop mûr » : une ressource, pas un déchet.}
En pleine saison, 30 à 40\% des fruits sont écartés du circuit commercial frais pour défauts esthétiques (taches, forme irrégulière, coups) ou sur-maturité. Jadis jetés en décharge à ciel ouvert (ex: décharge de Nkolfoulou à Yaoundé, PK10 à Douala), ils génèrent :
\begin{itemize}
    \item \textbf{Émissions de gaz à effet de serre :} Méthane ($CH_4$) par fermentation anaérobie, dioxyde de carbone ($CO_2$).
    \item \textbf{Lixiviats polluants :} Jus acides, sucrés, riches en matière organique, contaminant nappes et cours d'eau (Wouri, Mfoundi).
    \item \textbf{Prolifération vecteurs :} Mouches, rongeurs, moustiques (paludisme, dengue).
\end{itemize}

\paragraph{Valorisation par la biotechnologie de conservation.}
Les techniques étudiées (séchage, confiture, jus, pâte) sont \textbf{agnostiques à l'esthétique} : elles n'exigent que la \textbf{qualité sanitaire et gustative} (absence de pourriture avancée, de mycotoxines).
\begin{itemize}
    \item \textbf{Séchage :} Idéal pour fruits très mûrs (forte teneur en sucre, saveur concentrée).
    \item \textbf{Confiture / Pâte :} Utilise fruits éclatés, trop mûrs (pectine naturelle libérée).
    \item \textbf{Jus / Nectar :} Valorise les « jus de coupe » et fruits non calibrés.
\end{itemize}

\begin{propriete}[Bilan matière et énergie : l'économie circulaire au niveau local]
Une unité de transformation semi-industrielle (500 kg/jour) met en œuvre une \textbf{boucle quasi fermée} :
\begin{align*}
\text{Entrées :} &\quad \text{Fruits « non commercialisables frais»} + \text{Énergie solaire/biomasse} + \text{Eau} \\
\text{Sorties Produits :} &\quad \text{Mangue séchée, Confiture, Jus, Pâte de tomate (Valeur marchande)} \\
\text{Sorties « Déchets » valorisés :} &\quad 
\begin{cases}
\text{Épluchures, noyaux, pépins} \rightarrow \text{Compost (amendement sols acides)} \text{ OU } \text{Aliment bétail/porcs} \\
\text{Eaux de lavage/blanchiment} \rightarrow \text{Lagunage / Irrigation fertilisante} \\
\text{Bagasse (tomate)} \rightarrow \text{Extraction lycopène / Colorant naturel / Biogaz}
\end{cases}
\end{align*}
\textbf{Résultat :} Taux de valorisation massique $> 95\%$. L'empreinte carbone du kg de mangue séchée locale est \textbf{10 à 20 fois inférieure} à celui du jus de fruit importé (concentré + eau + transport maritime + chaîne du froid).
\end{propriete}

\begin{aretenir}[Intérêt environnemental synthétique]
La conservation des fruits de saison s'inscrit pleinement dans les \textbf{Objectifs de Développement Durable (ODD)} :
\begin{itemize}
    \item \textbf{ODD 12.3 :} Réduire de moitié le gaspillage alimentaire mondial d'ici 2030.
    \item \textbf{ODD 13 :} Lutte contre le changement climatique (évitement $CH_4$ décharges, substitution importations).
    \item \textbf{ODD 15 :} Préservation des sols par le compostage des résidus de transformation.
\end{itemize}
\end{aretenir}

\courssub{Intérêt social : emploi rural, autonomisation et sécurité alimentaire des ménages}

\paragraph{Un levier pour l'emploi des femmes et des jeunes.}
Au Cameroun, la transformation artisanale et semi-industrielle des fruits est \textbf{à 80\% féminine} (Groupements d'Initiative Commune - GIC, associations). C'est une activité :
\begin{itemize}
    \item \textbf{Compatible} avec les charges domestiques (séchage solaire à domicile, cuisson confiture).
    \item \textbf{Peu capitalistique} au démarrage (bassines, couteaux, séchoirs solaires améliorés $\approx$ 50 000 - 200 000 FCFA).
    \item \textbf{Génératrice de revenus quotidiens/hebdomadaires} (vente directe, marchés hebdomadaires), contrairement aux cultures de rente (cacao, café) payées une fois l'an.
\end{itemize}
Les jeunes y trouvent des débouchés techniques (maintenance séchoirs solaires, conditionnement, marketing digital, logistique moto-taxi « bendskin » pour la distribution).

\paragraph{Sécurité alimentaire des ménages transformateurs.}
L'autoconsommation d'une partie de la production transformée (confiture pour le petit-déjeuner des enfants, tomate séchée pour la sauce en hors-saison) améliore directement la \textbf{diversité alimentaire} du ménage (score HDDS - Household Dietary Diversity Score), indicateur clé de la sécurité nutritionnelle.

\courssub{Bilan synthétique : la biotechnologie au service du développement durable}

\paragraph{La leçon en une image mentale.}
La biotechnologie de la conservation des fruits (Module IV) est le \textbf{point de convergence} entre :
\begin{enumerate}
    \item \textbf{Les savoirs fondamentaux :} Microbiologie (croissance bactérienne, spores, $a_w$, pH, $T^\circ$), Enzymologie (PPO, PME, inactivation thermique), Chimie (oxydation, réaction de Maillard, gélification pectine).
    \item \textbf{Les savoir-faire techniques :} Maîtrise des procédés (séchage, stérilisation, concentration osmotique, acidification), Hygiène (BPH, HACCP), Conditionnement.
    \item \textbf{Les besoins concrets des populations :} Manger mieux (nutrition), Gagner mieux (économie), Vivre mieux (environnement, social).
\end{enumerate}

\begin{methode}[Rédaction d'une argumentation « Économique + Nutritionnel + Environnemental » (Type Bac)]
\textbf{Consigne type :} « En vous appuyant sur l'exemple de la mangue ou de la tomate, montrez que la transformation des fruits de saison répond aux enjeux du développement durable. » (6-8 points)

\textbf{Structure en 3 paragraphes distincts (1 paragraphe = 1 pilier) :}

\textbf{§1. Argument économique (Création de valeur / Revenus / Prix).}
\begin{itemize}
    \item \textit{Amorce :} « La transformation convertit des pertes post-récolte (30-50\%) en produits à forte valeur ajoutée. »
    \item \textit{Preuve chiffrée :} « 1 kg mangue fraîche à 100 FCFA $\rightarrow$ 15-18 kg mangue séchée vendue 1 500-2 000 FCFA/100g (VA $\times$ 15-20). »
    \item \textit{Mécanisme :} « Stabilisation prix producteur (débouché garanti) et prix consommateur (disponibilité hors-saison). Création d'emplois ruraux (GIC femmes, jeunes), substitution importations (concentré tomate, jus). »
\end{itemize}

\textbf{§2. Argument nutritionnel (Couverture besoins / Lutte carences).}
\begin{itemize}
    \item \textit{Amorce :} « Elle assure la disponibilité annuelle en micronutriments essentiels. »
    \item \textit{Preuve scientifique :} « Courbes de dégradation vitamine C : fruit frais $\rightarrow$ 0 mg/100g en 20 j ; conserve/séché $\rightarrow$ stable 18-25 mg/100g à 90 j. »
    \item \textit{Impact :} « Prévention xérophtalmie (vit A mangue), scorbut (vit C), anémie (folates). Diversification régime tubercules/céréales. »
\end{itemize}

\textbf{§3. Argument environnemental (Économie circulaire / Climat).}
\begin{itemize}
    \item \textit{Amorce :} « Elle valorise les fruits « laids » et sur-mûrs, détournant des décharges des tonnes de biomasse. »
    \item \textit{Preuve :} « Réduction émissions $CH_4$ (anaérobie décharge) et lixiviats. Valorisation résidus (épluchures $\rightarrow$ compost/biogaz ; noyaux $\rightarrow$ combustible/briquettes). »
    \item \textit{Bilan :} « Empreinte carbone locale $\ll$ importations. Contribution ODD 12, 13, 2. »
\end{itemize}

\textbf{Conclusion transversale :} « La biotechnologie de conservation n'est pas une simple technique de cuisine ; c'est un \textbf{levier de résilience territoriale} qui articule science (microbiologie/enzymologie), économie (chaîne de valeur) et écologie (circularité). »
\end{methode}

\begin{exoresolu}[Comparer les intérêts de deux techniques pour un même fruit]
\textbf{Énoncé :} Un GIC de la région du Centre hésite entre investir dans un \textbf{séchoir solaire amélioré} (capacité 50 kg/jour) ou un \textbf{autoclave pour conserves} (capacité 200 bocaux/jour) pour transformer la tomate. Comparez les deux options selon les critères : Investissement initial, Coût de fonctionnement, Qualité nutritionnelle (lycopène, vit C), Durée de conservation, Accessibilité marché (urbain/rural), Impact environnemental.

\textbf{Solution détaillée :}

\begin{center}
\begin{tabularx}{\linewidth}{|l|X|X|}
\hline
\textbf{Critère} & \textbf{Séchage solaire (Tomate séchée)} & \textbf{Stérilisation autoclave (Concentré/Conserve)} \\
\hline
Investissement initial & \textbf{Faible} (500 000 - 1 500 000 FCFA : structure bois/béton, filets, capteur solaire) & \textbf{Élevé} (3 000 000 - 8 000 000 FCFA : autoclave inox, chaudière, cuves, réseau vapeur/eau) \\
\hline
Coût fonctionnement & \textbf{Quasi nul} (Énergie solaire gratuite, main d'œuvre familiale) & \textbf{Significatif} (Bois/Gaz/Electricité pour vapeur, eau traitée, maintenance pression, main d'œuvre qualifiée) \\
\hline
Qualité nutritionnelle & \textbf{Lycopène :} Concentré (biodisponibilité $\uparrow$). \textbf{Vit C :} Perte 40-60\% (chaleur longue + oxydation résiduelle). \textbf{Réhydratation nécessaire.} & \textbf{Lycopène :} Très biodisponible (cis-isomères). \textbf{Vit C :} Perte 20-30\% (court, anaérobie). \textbf{Prêt à l'emploi.} \\
\hline
Durée conservation & \textbf{12-18 mois} (si $a_w < 0,6$, emballage barrière Alu/PE, $T^\circ < 25^\circ$C) & \textbf{24-36 mois} (stérilité commerciale, emballage hermétique, stable $T^\circ$ ambiante) \\
\hline
Accessibilité marché & \textbf{Rural / Peri-urbain :} Léger, transport facile, pas de chaîne du froid. \textbf{Urbain :} Produit « premium » (apéro, cuisine gastronomique). & \textbf{Urbain / Supermarchés / Export :} Format standard (bocal/boîte), usage cuisine quotidienne (sauce). Poids lourd (eau + verre/métal) = frein transport rural. \\
\hline
Impact environnemental & \textbf{Très faible} (Énergie renouvelable, 0 émission directe, résidus $\rightarrow$ compost). & \textbf{Modéré} (Énergie fossile/bois pour vapeur, eau consommation, emballages verre/métal recyclables mais lourds). \\
\hline
\end{tabularx}
\end{center}

\textbf{Synthèse pour le GIC :}
\begin{itemize}
    \item \textbf{Choisir Séchage} si : Capital limité, zone rurale sans électricité stable, main d'œuvre féminine disponible, cible marché « snacking sain / export bio / cuisine moderne », objectif empreinte carbone minimale.
    \item \textbf{Choisir Stérilisation} si : Accès électricité/bois bon marché, volume matière première très gros (> 500 kg/jour), cible marché masse urbaine (sauce tomate quotidienne), recherche durée de vie max et produit « prêt à cuisiner ».
    \item \textbf{Stratégie optimale (observée au Cameroun) :} \textbf{Diversification.} Séchage pour la tomate « cerise/ronde » à haute teneur en sec (rendement séchage meilleur) + Stérilisation pour la tomate « Roma/longue » juteuse (meilleur rendement jus/concentré). Cela sécurise les débouchés et lisse les risques techniques.
\end{itemize}
\end{exoresolu}

\begin{attention}[Vigilance sanitaire : condition \textit{sine qua non} de la durabilité]
Tous ces intérêts (économiques, nutritionnels, environnementaux, sociaux) s'effondrent si la \textbf{maîtrise sanitaire} fait défaut.
\begin{itemize}
    \item \textbf{Risque biologique :} \textit{Clostridium botulinum} (conserves mal stérilisées, pH $> 4,6$), mycotoxines (aflatoxines sur mangues séchées mal séchées/stockées humides), \textit{Salmonella} (hygiène manipulation).
    \item \textbf{Conséquence :} Intoxications graves (botulisme = urgence vitale), destruction image de marque filière, interdiction export, pertes financières totales.
    \item \textbf{Règle d'or :} \textbf{Jamais de raccourci sur l'hygiène (BPH) et la validation du procédé (Barèmes temps/T°/pH/$a_w$).} Un produit « bio » et « local » mais toxique est un échec total de la biotechnologie.
\end{itemize}
\end{attention}

\begin{savaistu}[Innovation camerounaise : le séchoir hybride solaire-biomasse]
Des chercheurs de l'ENSP (Yaoundé) et de l'Université de Dschang ont mis au point des \textbf{séchoirs hybrides} : capteur solaire jour + brûleur biomasse (coques noix de cajou, briquettes sciure) nuit/jour nuageux. Ils garantissent une température constante (50-\SI{60}{\celsius}) et un flux d'air contrôlé, réduisant le temps de séchage de 3-4 jours à 12-18 h, limitant l'oxydation (couleur orange vif préservée, vitamine C $\uparrow$ 30\%) et excluant la contamination poussières/insectes. Une biotechnologie \textbf{adaptée au contexte énergétique local}.
\end{savaistu}

\newpage

\coursec{Activité d'intégration}

\courssub{Objectifs de l'activité}
\begin{itemize}
    \item Mobiliser les connaissances sur les \cle{facteurs d'altération} (micro-organismes, enzymes, oxydation) et les \cle{principes de conservation} (destruction/inhibition des agents, élimination de l'eau, barrière physique) pour analyser une situation réelle.
    \item Interpréter des données expérimentales (développement microbien, teneur en vitamine C) et économiques (coûts, durée de conservation) pour \cle{comparer l'efficacité} de trois techniques : séchage solaire, confiture, conserve stérilisée.
    \item Rédiger une \cle{fiche-conseil technique} argumentée, adaptée au contexte d'une coopérative de Garoua, incluant le protocole détaillé de la technique retenue et les \cle{règles d'hygiène} (BPH - Bonnes Pratiques d'Hygiène) indispensables.
\end{itemize}

\courssub{Situation : La coopérative « Manguiers du Bénoué » à Garoua}
\textbf{Contexte.} La coopérative « Manguiers du Bénoué » regroupe 40 producteurs dans la région du Nord (Garoua). Chaque saison (avril--juillet), elle récolte environ \SI{20}{\tonne} de mangues (variété \emph{Kent} et \emph{Amélie}). Faute de routes praticables en saison des pluies et d'unités de transformation locales, \textbf{près de 50 \% de la production (\SI{10}{\tonne}) est perdue} (pourriture sur l'arbre, au sol ou lors du stockage temporaire). La coopérative souhaite investir dans une unité de transformation artisanale améliorée pour valoriser ces pertes. Trois options techniques sont envisagées : le \textbf{séchage solaire amélioré}, la \textbf{confiture artisanale} et la \textbf{conserve stérilisée (bocaux en verre)}.

\textbf{Document 1 : Climat de Garoua en période de récolte (avril--juillet) -- Moyennes 2015-2023 (Source : ASECNA / IRAD).}
\begin{center}
\begin{tabularx}{\linewidth}{|c|X|X|X|X|}\hline
\rowcolor{popblueL}
\textbf{Mois} & \textbf{Temp. moy. (\si{\celsius})} & \textbf{Temp. max (\si{\celsius})} & \textbf{Humidité relative moy. (\%)} & \textbf{Ensoleillement (h/jour)} \\ \hline
Avril & 32,5 & 39,0 & 55 & 8,2 \\ \hline
Mai & 31,0 & 37,5 & 70 & 7,5 \\ \hline
Juin & 29,5 & 35,0 & 80 & 6,8 \\ \hline
Juillet & 27,8 & 32,5 & 85 & 5,5 \\ \hline
\end{tabularx}
\end{center}
\vspace{0.5em}
\noindent\textit{Note : La saison des pluies s'installe progressivement à partir de mai. L'humidité relative élevée favorise le développement fongique.}

\textbf{Document 2 : Expérience comparative -- Développement de moisissures sur rondelles de mangue (variété Kent, \SI{1}{\cm} d'épaisseur).}
\emph{Protocole :} \SI{50}{\gram} de rondelles par modalité, incubées à \SI{30}{\celsius} (température moyenne jour Garoua), HR \SI{75}{\%}. Observation quotidienne du pourcentage de surface colonisée par \emph{Aspergillus niger} et \emph{Rhizopus sp.} pendant 7 jours.
\begin{center}
\begin{tabularx}{\linewidth}{|c|X|X|X|}\hline
\rowcolor{popblueL}
\textbf{Jour} & \textbf{Témoin : Fraîche (non traitée)} & \textbf{Modalité A : Blanchie \SI{3}{\min} à \SI{90}{\celsius} + Séchée (\SI{15}{\%} H₂O)} & \textbf{Modalité B : Séchée directe (\SI{15}{\%} H₂O, sans blanchiment)} \\ \hline
J0 & 0 \% & 0 \% & 0 \% \\ \hline
J1 & 15 \% & 0 \% & 5 \% \\ \hline
J2 & 45 \% & 0 \% & 12 \% \\ \hline
J3 & 85 \% (pourriture) & 0 \% & 25 \% \\ \hline
J4 & 100 \% & 0 \% & 40 \% \\ \hline
J5 & -- & 0 \% & 60 \% \\ \hline
J7 & -- & 5 \% (contamination secondaire) & 90 \% \\ \hline
\end{tabularx}
\end{center}
\vspace{0.5em}
\noindent\textit{Rappel : Le blanchiment (\SI{90}{\celsius}, \SI{3}{\min}) vise l'inactivation des enzymes (polyphénoloxydase, pectinestérase) et la destruction de la flore de surface. Le séchage vise une activité de l'eau $a_w < 0,60$ (ici \SI{15}{\%} H₂O résiduelle).}

\textbf{Document 3 : Bilan technico-économique prévisionnel pour \SI{1}{\tonne} de mangues fraîches entrantes (Données coopérative + CIPMEN Garoua, 2023).}
\begin{center}
\begin{tabularx}{\linewidth}{|c|X|X|X|}\hline
\rowcolor{popblueL}
\textbf{Critère} & \textbf{Séchage solaire (bâche surélevée, 3 jours)} & \textbf{Confiture (cuisson \SI{30}{\min}, \SI{60}{\%} sucre)} & \textbf{Conserve stérilisée (bocaux \SI{500}{g}, \SI{100}{\celsius}, \SI{20}{\min})} \\ \hline
Produit fini obtenu & \SI{180}{kg} mangues séchées & \SI{1,4}{t} confiture & \SI{1,1}{t} conserves (égouttées) \\ \hline
Rendement massique & 18 \% & 140 \% (ajout sucre/eau) & 110 \% (ajout sirop) \\ \hline
Durée de conservation (conditions ambiantes Garoua) & 12 à 18 mois ($a_w \approx 0,55$) & 18 à 24 mois ($a_w \approx 0,80$, sucre + pH) & 24 à 36 mois (stérilité commerciale) \\ \hline
Coût intrants (FCFA/t fraîche) & 45 000 (bâches, main d'œuvre, emballage) & 180 000 (sucre, citron, gaz/bois, pots, étiquettes) & 220 000 (bocaux, capsules, énergie stérilisation, main d'œuvre) \\ \hline
Prix de vente local estimé (FCFA/kg ou pot) & 3 500 / kg & 1 200 / pot \SI{300}{g} & 800 / bocal \SI{500}{g} \\ \hline
Chiffre d'affaires prévisionnel (FCFA) & 630 000 & 5 600 000 & 1 760 000 \\ \hline
Marge brute (CA - Intrants) & 585 000 & 5 420 000 & 1 540 000 \\ \hline
Investissement matériel initial (FCFA) & 300 000 (séchoirs, tables) & 800 000 (marmites, balance, scelleuse) & 2 500 000 (autoclave, bocaux, source énergie stable) \\ \hline
\end{tabularx}
\end{center}

\textbf{Document 4 : Évolution de la teneur en vitamine C (acide ascorbique) selon le procédé de transformation.}
\emph{Mesure par dosage au DCCIP (2,6-dichlorophénolindophénol) sur produit fini. Valeurs en mg/100 g de produit fini. Mangue fraîche initiale : \SI{45}{\mg/100g}.}
\begin{popfigure}
\begin{tikzpicture}
\begin{axis}[
    width=\linewidth, height=8cm,
    ybar=5pt, bar width=18pt,
    ylabel={Teneur en Vitamine C (\si{\mg/100g})},
    symbolic x coords={Fraiche,Sechage,Confiture,Conserve},
    xtick=data,
    xticklabels={Fraîche, Séchage, Confiture, Conserve},
    ymin=0, ymax=50,
    ymajorgrids=true,
    grid style={popline},
    nodes near coords, nodes near coords align={vertical},
    every node near coord/.append style={font=\small, popdark},
    legend style={at={(0.5,-0.2)}, anchor=north, legend columns=-1},
]
\addplot[fill=popblue!70] coordinates {(Fraiche,45) (Sechage,22) (Confiture,12) (Conserve,16)};
\legend{Teneur en Vitamine C}
\end{axis}
\end{tikzpicture}
\legende{Figure 1 -- Teneur en vitamine C résiduelle dans les produits finis (mg/100 g). Le séchage conserve mieux la vitamine C que les traitements thermiques longs (confiture, conserve), mais la concentration en sucre de la confiture dilue la teneur pour 100 g de produit.}
\end{popfigure}

\courssub{Consignes}
\begin{enumerate}
    \item \textbf{Analyse des facteurs d'altération (Doc 1 \& 2).} À partir du Document 1, identifiez les deux facteurs climatiques majeurs de Garoua en saison des mangues qui favorisent les pertes post-récolte. Expliquez, en vous appuyant sur les résultats du Document 2 (Témoin), le rôle des micro-organismes et des enzymes dans l'altération rapide de la mangue fraîche.
    \item \textbf{Interprétation de l'expérience (Doc 2).} Comparez l'efficacité du blanchiment + séchage (Modalité A) vs séchage seul (Modalité B) pour limiter le développement des moisissures. Quel facteur d'altération le blanchiment cible-t-il spécifiquement en plus de la flore de surface ? Justifiez la réapparition de moisissures à J7 en Modalité A.
    \item \textbf{Comparaison technico-économique et nutritionnelle (Doc 3 \& 4).} Complétez le tableau de synthèse ci-dessous en utilisant les Documents 3 et 4. Classez les trois techniques selon : (a) la durée de conservation, (b) la marge brute par tonne de frais, (c) l'investissement initial, (d) la conservation de la vitamine C (mg/100g produit fini).
    \begin{center}
    \begin{tabularx}{\linewidth}{|c|X|X|X|}\hline
    \rowcolor{popblueL}
    \textbf{Critère} & \textbf{Séchage} & \textbf{Confiture} & \textbf{Conserve} \\ \hline
    Durée de conservation (mois) & & & \\ \hline
    Marge brute (FCFA/t fraîche) & & & \\ \hline
    Investissement initial (FCFA) & & & \\ \hline
    Vitamine C résiduelle (mg/100g) & & & \\ \hline
    \end{tabularx}
    \end{center}
    \item \textbf{Choix argumenté et fiche-conseil.} La coopérative dispose d'un capital de départ de \SI{1\,000\,000}{FCFA} et d'un accès limité à l'électricité (groupe électrogène coûteux). Elle vise l'autonomie alimentaire (consommation locale) et la vente sur les marchés de Garoua et Maroua. \textbf{Proposez une ou deux technique(s) prioritaire(s)}. Justifiez votre choix en croisant \textbf{au moins 4 critères} (technique, économique, nutritionnel, faisabilité, environnemental). Rédigez la \textbf{fiche-conseil} comprenant : le protocole détaillé (étapes, températures, durées, critères de fin) de la technique retenue et les \textbf{5 points critiques d'hygiène (BPH)} à respecter impérativement.
\end{enumerate}

\courssub{Ressources à mobiliser}
\begin{itemize}
    \item \textbf{Facteurs d'altération :} Micro-organismes (moisissures, levures, bactéries) -- Enzymes endogènes (PPO, pectinestérase) -- Oxydation (air, lumière, chaleur) -- Facteurs physico-chimiques ($a_w$, pH, $E_h$, T°).
    \item \textbf{Principes de conservation :} \cle{Abiosis} (stérilisation, pasteurisation), \cle{Anabiose} (froid, séchage, $a_w$ bas, conservateurs), \cle{Barrière} (emballage hermétique).
    \item \textbf{Seuils clés :} $a_w < 0,60$ arrête moisissures/levures ; $a_w < 0,70$ arrête bactéries. Blanchiment : \SI{85-95}{\celsius}, \SI{1-5}{\min} (inactivation PPO/PME). Traitement thermique conserve (pH mangue $\approx 3,8-4,2$) : pasteurisation \SI{90-95}{\celsius} \SI{20}{\min} souvent suffisante, stérilisation \SI{100}{\celsius} \SI{20}{\min} préconisée pour sécurité maximale.
    \item \textbf{Vitamine C :} Thermosensible, hydrosoluble, sensible à l'oxydation (PPO, Cu/Fe, lumière, air). Perte $\propto$ T° $\times$ temps $\times$ surface de contact air/eau.
    \item \textbf{BPH (Bonnes Pratiques d'Hygiène) :} Matière première saine, eau potable, surfaces/ustensiles nettoyés/désinfectés, hygiène personnel (lavage mains, tenues), traçabilité/étiquetage (DLUO, lot).
\end{itemize}

\begin{exoresolu}[Démarche et corrigé commenté]
\textbf{Question 1 -- Facteurs climatiques et altération du témoin.}
\begin{itemize}
    \item \textbf{Doc 1 :} Deux facteurs majeurs : \textbf{Températures élevées} (moyenne \SI{28-33}{\celsius}, max \SI{39}{\celsius} en avril) et \textbf{Humidité relative élevée} (\SI{70-85}{\%} de mai à juillet). La chaleur accélère le métabolisme microbien et enzymatique (règle Q10 : vitesse $\times 2$ à \SI{3}{\celsius} près par \SI{10}{\celsius}). L'humidité maintient une $a_w$ de surface élevée, condition \emph{sine qua non} de la germination des spores fongiques.
    \item \textbf{Doc 2 (Témoin) :} Colonisation à \SI{85}{\%} en J3, pourriture totale J4. Cela démontre l'action conjuguée : (1) \textbf{Contamination de surface} (spores air, sol, insectes) $\rightarrow$ germination rapide (T°/HR optimales) $\rightarrow$ mycélium visible ; (2) \textbf{Enzymes endogènes} (PPO = brunissement enzymatique, PME = déméthoxylation pectines $\rightarrow$ ramollissement tissus) qui préparent le substrat (libération sucres, acides) et facilitent la pénétration fongique. L'oxydation phénolique (brunissement) dégrade aussi la qualité organoleptique et nutritionnelle (vit C).
\end{itemize}

\textbf{Question 2 -- Efficacité blanchiment + séchage vs séchage seul.}
\begin{itemize}
    \item \textbf{Modalité A (Blanchiment + Séchage) :} \textbf{0 \% moisissure jusqu'à J5}, \SI{5}{\%} à J7 (contamination secondaire post-séchage/manipulation). \textbf{Efficacité excellente.}
    \item \textbf{Modalité B (Séchage direct) :} Départ lent (J1: \SI{5}{\%}, J2: \SI{12}{\%}) mais exponentielle (J7: \SI{90}{\%}). \textbf{Efficacité insuffisante.}
    \item \textbf{Analyse :} Le séchage seul (cible $a_w < 0,60$) inhibe la \emph{croissance} microbienne mais ne \emph{tue} pas les spores/propagules déjà présentes sur la peau et dans les premiers mm de chair. Lors du séchage (3 jours au soleil), l'$a_w$ descend progressivement ; pendant cette phase transitoire ($0,95 \rightarrow 0,60$), les moisissures (xérotolérantes comme \emph{Aspergillus}) se développent. Le blanchiment (\SI{90}{\celsius}, \SI{3}{\min}) apporte un \textbf{effet létal} sur la flore de surface (bactéries, levures, spores fongiques superficielles) \textbf{ET inactive les enzymes} (PPO, PME) responsables du brunissement et du ramollissement. C'est une \textbf{barrière combinée} (thermique + $a_w$).
    \item \textbf{Reprise J7 Modalité A :} Contamination secondaire lors du conditionnement (air ambiant chargé en spores, mains non lavées, emballages non stériles) + possible réveil de spores thermorésistantes (rarissime à \SI{90}{\celsius} 3 min) si $a_w$ remonte (réhumidification par air humide Garoua juillet HR 85\%).
\end{itemize}

\textbf{Question 3 -- Tableau de synthèse comparatif (Docs 3 \& 4).}
\begin{center}
\begin{tabularx}{\linewidth}{|c|X|X|X|}\hline
\rowcolor{popblueL}
\textbf{Critère} & \textbf{Séchage} & \textbf{Confiture} & \textbf{Conserve} \\ \hline
Durée de conservation (mois) & 12--18 & 18--24 & 24--36 \\ \hline
Marge brute (FCFA/t fraîche) & 585 000 & \textbf{5 420 000} & 1 540 000 \\ \hline
Investissement initial (FCFA) & \textbf{300 000} & 800 000 & 2 500 000 \\ \hline
Vitamine C résiduelle (mg/100g) & \textbf{22} & 12 & 16 \\ \hline
\end{tabularx}
\end{center}
\vspace{0.5em}
\noindent\textit{Commentaires :} La confiture génère la plus forte marge brute (effet de levier sucre + prix vente), mais nécessite plus de capital que le séchage. La conserve a la plus longue DVC mais investissement très lourd (autoclave, bocaux, énergie). Le séchage est le meilleur compromis nutritionnel (vit C) et financier (faible CAPEX).

\textbf{Question 4 -- Choix argumenté et Fiche-Conseil.}
\textbf{Choix recommandé : \textbf{Séchage solaire amélioré} en technique principale + \textbf{Confiture} en technique secondaire (valorisation des fruits trop mûrs/non calibrés pour le séchage).}

\textbf{Justification multicritères :}
\begin{enumerate}
    \item \textbf{Faisabilité technique \& Capital (CAPEX) :} Investissement \SI{300\,000}{FCFA} $<$ Capital dispo \SI{1\,000\,000}{FCFA}. Pas d'électricité requise (soleil gratuit). Maîtrisable par le groupe (formation courte). La conserve (\SI{2,5}{M} FCFA, énergie continue) et la confiture (\SI{800\,000}{FCFA}, achat sucre \SI{600\,000}{FCFA}/t) sont trop risquées pour démarrer.
    \item \textbf{Durée de conservation \& Logistique :} 12-18 mois à $a_w \approx 0,55$ permet de lisser la vente sur l'année (prix bas en saison, haut hors saison). Produit stable, léger (transport facile Garoua $\rightarrow$ Maroua/Douala), pas de chaîne du froid.
    \item \textbf{Qualité nutritionnelle :} Meilleure rétention Vitamine C (\SI{22}{mg/100g} vs \SI{12-16}{mg}). Concentration en fibres, minéraux, caroténoïdes (provitamine A) préservée. Pas d'ajout de sucre (intérêt santé publique : diabète, obésité).
    \item \textbf{Environnement \& Déchets :} Zéro déchet d'emballage lourd (sacs PP recyclables vs bocaux verre/couverts métal). Faible empreinte carbone (solaire vs bois/gaz pour confiture/conserve). Valorisation des écorces/noyaux (compost, alimentation animale, extraction pectine/huile).
    \item \textbf{Rentabilité :} Marge brute \SI{585\,000}{FCFA/t} correcte pour un démarrage. ROI (Retour sur Investissement) < 1 saison. La confiture sera introduite en Année 2 avec les bénéfices du séchage pour acheter le sucre en gros.
\end{enumerate}

\textbf{FICHE-CONSEIL TECHNIQUE : SÉCHAGE SOLAIRE AMÉLIORÉ DE LA MANGUE (Variété Kent/Amélie)}
\vspace{0.5em}
\noindent\textbf{Objectif :} Produire \SI{180}{kg} de mangues séchées de qualité commerciale (catégorie I) à partir de \SI{1}{t} fraîches, avec $a_w < 0,60$, couleur orange vive, teneur Vit C $> \SI{20}{mg/100g}$.

\begin{methode}{Protocole détaillé (10 étapes)}
\begin{enumerate}
    \item \textbf{Réception \& Tri (J0, matin) :} Pesée. Tri manuel : élimination fruits pourris, piqués (mouches des fruits), trop mous (stade 4-5). Conservation fruits triés à l'ombre (\SI{<25}{\celsius} si possible) max 4h.
    \item \textbf{Lavage (Eau potable + chlore) :} Bain \SI{5}{\min} eau + \SI{50}{ppm} chlore actif (ex: \SI{15}{ml} eau de Javel \SI{2,6}{\%} pour \SI{10}{L} eau). Rinçage \textbf{obligatoire} 2 bains eau potable claire. Égouttage sur claies propres.
    \item \textbf{Épluchage \& Dénoyautage :} Couteaux inox aiguisés. Épluchage fin (perte < 10\%). Dénoyautage soigneux (pas de fragment de noyau). Travail sur tables inox/HDPE, mains gantées (gants nitrile changés toutes les 2h).
    \item \textbf{Tranchage :} Rondelles régulières \textbf{\SI{0,8 - 1,0}{cm}} d'épaisseur (tranchoir manuel inox). Épaisseur uniforme = séchage homogène. Éviter \SI{<0,5}{cm} (cassant, perte vit C) ou \SI{>1,5}{cm} (cœur humide, moisissure).
    \item \textbf{BLANCHIMENT (ÉTAPE CRITIQUE) :} Plongeoir inox. Eau bouillante (\SI{90-95}{\celsius}). \textbf{Durée : \SI{3}{\min} chrono} départ ré-ébullition. Ratio fruits/eau = 1/3 (maintien T°). \textit{But : Inactivation PPO/PME (brunissement/ramollissement) + pasteurisation surface.}
    \item \textbf{Refroidissement \& Égouttage Express :} Sortie immédiate $\rightarrow$ bain eau potable fraîche (\SI{<20}{\celsius}) \SI{1}{\min} (arrêt cuisson, choc thermique). Égouttage \SI{10}{\min} sur claies (perte eau de surface = gain temps séchage).
    \item \textbf{SULFITAGE (Optionnel mais recommandé pour couleur/Vit C) :} Trempage \SI{5}{\min} dans solution \textbf{\SI{500}{ppm} $SO_2$} (ex: \SI{1,5}{g} métabisulfite de potassium $K_2S_2O_5$ pour \SI{10}{L} eau). \textit{Rôle : Antioxydant (protège Vit C, caroténoïdes), antimicrobien résiduel, inhibe PPO résiduelle.} Rinçage léger \SI{30}{s} eau potable. \textbf{Attention :} Mentionner "sulfites" sur étiquette (allergènes). Si refus sulfites $\rightarrow$ blanchiment \SI{4}{\min} + séchage ultra-rapide.
    \item \textbf{CHARGEMENT SÉCHOIRS SOLAIRES AMÉLIORÉS :} Disposition monocouche sur claies (grille inox/HDPE alimentaire), espacement \SI{2}{cm} entre rondelles. Séchoir type "tunnel" ou "bâche surélevée" (hauteur \SI{1}{m}, orientation Est-Ouest, pente \SI{15}{\degree}), face transparente UV-stabilisée, fond noir absorbant. Chargement max \SI{5}{kg/m^2}.
    \item \textbf{SÉCHAGE (3 à 4 jours selon météo Doc 1) :} 
    \begin{itemize}
        \item \textbf{Jour 1 (Avril/Mai, T° max 37-39°C) :} T° interne séchoir \SI{50-55}{\celsius}. HR sortie < 30\%. Retourner rondelles à mi-journée. Poids cible fin J1 : $\approx$ \SI{40}{\%} poids initial.
        \item \textbf{Jour 2-3 :} T° \SI{45-50}{\celsius}. HR < 25\%. Contrôle visuel/tactile : souple, non collant, pas d'humidité au cœur (coupure nette).
        \item \textbf{Critère FIN :} Poids final \SI{18 \pm 2}{\%} poids frais initial ($a_w \approx 0,55$). Test "claquement" : la rondelle pliée ne casse pas mais ne colle pas.
    \end{itemize}
    \item \textbf{Conditionnement \& Stockage :} Refroidissement à T° ambiante \SI{30}{\min} (éviter condensation). Pesée lots \SI{100}{g}, \SI{250}{g}, \SI{500}{g}. \textbf{Emballage :} Sachets PP/PE alimentaire \SI{80}{\mu m} + barrière aluminium (métallisé) OU pots PP hermétiques. \textbf{Scellage à chaud} (soudeuse \SI{180}{\celsius}). Etiquetage : Nom produit, Ingrédients (Mangue, \emph{éventuel: sulfites}), Poids net, Date fabrication, DLUO (Fab + 12 mois), N° Lot, Adresse coopérative. Stockage local ventilé, sur palette, à l'abri lumière/sol/humidité/rongeurs.
\end{enumerate}
\end{methode}

\begin{attention}{5 Points Critiques BPH (Bonnes Pratiques d'Hygiène) -- NON NÉGOCIABLES}
\begin{enumerate}
    \item \textbf{Eau potable certifiée} pour TOUS les lavages, blanchiment, refroidissement, sulfitage (analyse bactériologique mensuelle, chlore résiduel \SI{0,2-0,5}{mg/L} sortie robinet). Eau de forage/puits non traitée = \textbf{REJET LOT}.
    \item \textbf{Hygiène du personnel :} Lavage mains savon + désinfectant (gel hydroalcoolique) \textbf{avant} entrée zone, \textbf{après} chaque pause/toilettes/manipulation déchets. Tenues propres (blouse, charlotte, gants), pas de bijoux, ongles courts, pas de maladie diarrhéique/ cutanée (certificat médical trimestriel).
    \item \textbf{Traçabilité Matière Première :} Registre entrées : Date, Producteur, Parcelle, Variété, Stade maturité, Traitements phytosanitaires (respect DAR - Délai Avant Récolte). Pas de fruits ramassés au sol (mycotoxines : aflatoxines, patuline).
    \item \textbf{Maîtrise du Blanchiment :} Thermomètre sonde \textbf{obligatoire} dans l'eau (pas au doigt !). Chronomètre. Si T° $< \SI{90}{\celsius}$ ou temps $< \SI{3}{\min}$ $\rightarrow$ \textbf{Relancer le cycle}. Sous-blanchiment = enzymes actives = brunissement + texture molle + risque toxines.
    \item \textbf{Contrôle $a_w$ / Humidité finale :} Hygromètre à sonde (type AquaLab ou capacitif étalonné) sur \textbf{chaque lot} avant conditionnement. $a_w$ \textbf{doit être} $\le 0,60$. Si $> 0,65$ $\rightarrow$ \textbf{Remise en séchage}. Conditionnement à $a_w > 0,70$ = développement moisissures \emph{Aspergillus} (aflatoxines) / levures (fermentation) en quelques jours climat Garoua.
\end{enumerate}
\end{attention}

\textbf{Indicateurs de suivi (Tableau de bord coopérative) :} Rendement massique (cible $\ge 18\%$), Taux rebuts tri (cible $< 5\%$), $a_w$ finale (cible $\le 0,60$), Vit C (cible $\ge \SI{20}{mg/100g}$), Analyse mycotoxines (aflatoxines $< \SI{4}{\mu g/kg}$ total, $< \SI{2}{\mu g/kg}$ B1) 1 fois/saison par labo (LANAVET/IRAD).
\end{exoresolu}

\courssub{Bilan de l'activité}
Cette activité d'intégration a permis de mettre en évidence que :
\begin{itemize}
    \item La \cle{conservation} n'est pas une action unique mais une \textbf{stratégie de barrières combinées} (Hurdle Technology) : thermique (blanchiment), physico-chimique ($a_w$ par séchage, pH/sucre pour confiture), et physique (emballage barrière).
    \item Le \textbf{choix technique} ne se fait pas sur un seul critère (ex: durée de conservation max), mais par \textbf{optimisation sous contrainte} (capital dispo, énergie, main d'œuvre, marché cible, qualité nutritionnelle, environnement). Pour Garoua, le \textbf{séchage solaire amélioré} s'impose comme le "premier pas" techniquement maîtrisable, financièrement accessible, nutritionnellement avantageux et écologiquement vertueux.
    \item La \cle{maîtrise des facteurs d'altération} (Doc 2) valide scientifiquement le protocole : le blanchiment n'est pas une option "qualité" mais une \textbf{nécessité sanitaire} (destruction spores surface + inactivation enzymatique) sans laquelle le séchage seul échoue en climat tropical humide (Modalité B).
    \item La \cle{valorisation des fruits de saison} transforme une perte économique (\SI{10}{t} pourries = 0 FCFA) en revenu durable (\SI{585\,000}{FCFA} marge brute/t pour le séchage), contribuant à la \textbf{sécurité alimentaire} (disponibilité fruits hors saison, vitamine A/C) et à la \textbf{lutte contre la pauvreté} rurale (revenus femmes/producteurs).
    \item La \cle{rigueur des BPH} (5 points critiques) est la condition \emph{sine qua non} de l'accès aux marchés formels (supermarchés, export CEMAC) et de la protection du consommateur (aflatoxines, hygiène).
\end{itemize}
\vspace{0.5em}
\noindent\textbf{Compétences évaluées :} C1 (Mobiliser des connaissances), C2 (Mener une démarche d'investigation), C3 (Communiquer/Argumenter), C4 (S'éduquer/Pratiquer responsable). Cette activité prépare directement aux épreuves pratiques (EPC) et écrites du Baccalauréat (analyse de documents, rédaction protocole, argumentation choix technique).

\newpage

\coursec{Méthodes et exercices résolus}

\coursec{Amélioration de la conservation des fruits de saison}

\courssub{Méthodes et exercices résolus}

\begin{methode}[Décrire une technique de transformation : la confiture de mangue]
    \textbf{Objectif :} Restituer les étapes opératoires et expliquer le rôle scientifique de chaque étape (stabilisation microbiologique, gélification, qualité organoleptique).
    \begin{enumerate}
        \item \textbf{Identifier le principe :} Conservation par association \textbf{chaleur (pasteurisation)} + \textbf{sucre (abaissement $a_w$)} + \textbf{acidité (pH < 4,5)} + \textbf{conditionnement hermétique}.
        \item \textbf{Lister les étapes dans l'ordre chronologique :}
            \begin{itemize}
                \item \textbf{Tri / Lavage :} Élimination des fruits pourris (sources de pectinestérases, moisissures, mycotoxines) et des souillures.
                \item \textbf{Épluchage / Dénoyautage :} Élimination de l'épicarpe (pesticides, cire, levures) et du noyau.
                \item \textbf{Cuisson (avec sucre + jus de citron) :} 
                    \begin{itemize}
                        \item \emph{Ébullition :} Destruction micro-organismes, inactivation enzymes (PPO, POD, pectinestérase), concentration en solubles (cible \SI{65-68}{\%} Brix / \SI{103-105}{\celsius}).
                        \item \emph{Rôle sucre :} Agent de conservation ($a_w \downarrow$ \SI{\approx 0,80-0,84}{}), texture, goût.
                        \item \emph{Rôle acide (citron) :} pH \SI{\approx 3,0-3,5}{} $\rightarrow$ favorise gélification pectine (DM > 50\%), inhibe \textit{Clostridium botulinum}, stabilise couleur (anthocyanes, caroténoïdes), hydrolyse partielle saccharose (anti-cristallisation).
                    \end{itemize}
                \item \textbf{Mise en pots à chaud (\SI{>85}{\celsius}) :} Pots stérilisés, remplissage \SI{5-10}{mm} du bord, fermeture immédiate capsules twist-off neuves. Stérilisation du contenant/fermeture, création vide partiel (dépression) au refroidissement.
                \item \textbf{Retournement / Refroidissement rapide :} \SI{5-10}{min} tête en bas $\rightarrow$ stérilisation couvercle/joint par condensation, élimination O$_2$ résiduel. Remise à l'endroit, refroidissement air/eau forcée.
                \item \textbf{Étiquetage / Stockage :} Lieu frais (\SI{<20}{\celsius}), sec, obscurité. DLUO \SI{12-18}{mois}.
            \end{itemize}
        \item \textbf{Justifier scientifiquement :} Pour chaque étape, citer le facteur d'altération maîtrisé (micro-organismes, enzymes, oxydation, eau libre).
    \end{enumerate}
    \textbf{Réflexe Bac :} Citer la cible \textbf{\SI{65-68}{\%} Brix} (ou \SI{103-105}{\celsius}) et le \textbf{pH < 4,5}. Ne pas confondre « stérilisation » (destruction totale formes végétatives + spores, $T \ge \SI{115}{\celsius}$) et « pasteurisation » (destruction formes végétatives pathogènes/altération, $T < \SI{100}{\celsius}$ ou \SI{100}{\celsius} acide) ; la confiture est une \emph{conserve sucrée acide} (pasteurisation + $a_w$ bas + pH bas).
\end{methode}

\begin{exoresolu}[Description et justification de la fabrication de confiture de mangue]
    \textbf{Énoncé :} Un groupement d'intérêt économique (GIE) de la région du Littoral souhaite valoriser les excédents de mangues (\emph{Mangifera indica}) en confiture. 
    \begin{enumerate}
        \item Énumérer les principales étapes de fabrication de la confiture de mangue, de la réception du fruit au stockage.
        \item Pour l'étape de cuisson, préciser la cible en \emph{degrés Brix} et en \emph{température} à atteindre. Justifier l'ajout de jus de citron.
        \item Expliquer pourquoi le retournement des pots à chaud assure une meilleure conservation.
    \end{enumerate}
    \tcblower
    \textbf{Solution détaillée :}
    \begin{enumerate}
        \item \textbf{Étapes de fabrication :}
            \begin{enumerate}
                \item \textbf{Réception et Tri :} Sélection mangues mûres, saines, fermes. Élimination fruits abîmés (source enzymes pectinolytiques, contaminants microbiens, mycotoxines).
                \item \textbf{Lavage :} Eau potable + détergent alimentaire $\rightarrow$ rinçage. Élimination poussières, résidus phytosanitaires, charge microbienne de surface.
                \item \textbf{Épluchage et Dénoyautage :} Manuels ou mécaniques. Épluchage élimine épicarpe (cire, spores fongiques). Dénoyautage retire noyau (amertume, corps dur).
                \item \textbf{Découpage / Pulpage :} Coupes en dés ou pulpeur-fineur (grille \SI{0,5-1}{mm}) pour pulpe homogène.
                \item \textbf{Cuisson (Concentration) :} Bassine cuivre/inox ou marmite double fond. Ajout \textbf{sucre cristallisé} (ratio pulpe/sucre $\approx$ 1/0,7 à 1/1) et \textbf{jus de citron} (\SI{10-20}{mL/kg} pulpe). Ébullition vigoureuse en remuant. \textbf{Cible : \SI{65-68}{\%} Brix (réfractomètre) ou \SI{103-105}{\celsius} (thermomètre)}. Écumage (élimination mousse = protéines, air, impuretés).
                \item \textbf{Mise en pots :} Remplissage pots verre stérilisés (ébouillantés), chauds, jusqu'à \SI{5-10}{mm} du bord. Fermeture immédiate couvercles neufs (capsules twist-off).
                \item \textbf{Pasteurisation / Retournement :} Retourner pots (\SI{5-10}{min}) $\rightarrow$ stérilisation couvercle/joint par chaleur produit. Remettre à l'endroit, refroidissement rapide $\rightarrow$ création vide (dépression \emph{« pop »}).
                \item \textbf{Étiquetage et Stockage :} Etiquette (nom, ingrédients, date, lot, DLUO). Stockage local frais (\SI{<20}{\celsius}), sec, obscurité.
            \end{enumerate}
        \item \textbf{Cuisson : Cibles et rôle du citron.}
            \begin{itemize}
                \item \textbf{Cible Brix : \SI{65-68}{\%}} (masses solubles totales). Correspond $a_w \approx 0,80-0,84$, inhibant bactéries (seuil $a_w \approx 0,90$), levures (seuil $\approx 0,88$), limitant moisissures (seuil $\approx 0,70$).
                \item \textbf{Cible Température : \SI{103-105}{\celsius}} (élévation ébullition due au sucre). Garantit destruction formes végétatives et concentration correcte.
                \item \textbf{Rôle du jus de citron (Acide citrique) :}
                    \begin{enumerate}
                        \item \textbf{Abaissement pH (3,0 - 3,5) :} Condition \emph{sine qua non} gélification pectines haut DM (> 50\%) naturellement présentes. Sans acide, pas de prise (sirop).
                        \item \textbf{Sécurité microbiologique :} pH < 4,5 empêche germination spores \textit{Clostridium botulinum} (botulisme).
                        \item \textbf{Stabilité couleur :} Stabilise pigments (caroténoïdes, anthocyanes) et limite brunissement enzymatique résiduel (PPO) et non enzymatique (Maillard).
                        \item \textbf{Inversion partielle saccharose :} Limite cristallisation sucre en stockage (formation glucose + fructose).
                    \end{enumerate}
            \end{itemize}
        \item \textbf{Justification du retournement :}
            \begin{itemize}
                \item Produit mis en pot \SI{>85}{\celsius}. Retournement met couvercle/joint en contact produit bouillant \SI{5-10}{min} $\rightarrow$ \textbf{stérilisation thermique couvercle + espace de tête} (air résiduel + condensation vapeur).
                \item Refroidissement $\rightarrow$ condensation vapeur $\rightarrow$ \textbf{dépression (vide partiel)}. Couvercle bombé vers l'intérieur (\emph{« clic »} ouverture = garantie intégrité).
                \item Élimine O$_2$ résiduel (limite oxydation, moisissures aérobies strictes) et assure étanchéité.
            \end{itemize}
    \end{enumerate}
    \textbf{Conclusion :} La confiture de mangue repose sur l'association traitements thermiques (cuisson, retournement), chimiques (sucre, acide) et physiques (conditionnement hermétique) pour garantir stabilité microbiologique à température ambiante (> 12 mois).
\end{exoresolu}

\begin{methode}[Élaborer un protocole de transformation : la pâte de tomate (concentré)]
    \textbf{Objectif :} Rédiger un mode opératoire complet, reproductible, hygiénique, adapté au contexte artisanal/semi-industriel camerounais.
    \begin{enumerate}
        \item \textbf{Définir la matière première :} Variété (Roma, Mongal - chair ferme, peu de graines, haute teneur lycopène/pectines), maturité (rouge vif, fermeté), quantité.
        \item \textbf{Structurer en postes unitaires (Schéma bloc) :} Réception $\rightarrow$ Tri/Lavage $\rightarrow$ Échaudage/Émonder $\rightarrow$ Pulpage/Raffinage $\rightarrow$ Concentration (Cuisson) $\rightarrow$ Conditionnement (Pots/Boîtes) $\rightarrow$ Traitement thermique final (Stérilisation) $\rightarrow$ Refroidissement $\rightarrow$ Contrôle/Stockage.
        \item \textbf{Préciser les paramètres critiques (PCC - Points Critiques de Contrôle / HACCP) :}
            \begin{itemize}
                \item \textbf{Lavage :} Eau chlorée (\SI{50-100}{ppm} chlore actif), temps contact \SI{2-5}{min}.
                \item \textbf{Échaudage :} \SI{90-95}{\celsius}, \SI{30-60}{s} $\rightarrow$ inactivation PPO/POD (brunissement), facilite émondage, expulsion air (meilleur vide).
                \item \textbf{Pulpage :} Taille maille raffineur \SI{0,5-0,8}{mm} (élimine peaux, pépins, fibres).
                \item \textbf{Concentration :} Évaporation sous vide (\SI{60-70}{\celsius}, \SI{0,6-0,8}{bar}) \textbf{OU} cuisson atmosphérique (\SI{100-105}{\celsius}) agitation constante. Cible : \textbf{Double/Triple concentré : \SI{28-30}{\%} Brix (\SI{28-30}{g/100g})}. \emph{Ne pas dépasser \SI{110}{\celsius} (dégradation lycopène, brunissement Maillard)}.
                \item \textbf{Stérilisation finale (si pots/boîtes) :} Autoclave \SI{115-121}{\celsius} / \SI{15-20}{min} (selon diamètre) $\rightarrow$ $F_0 \ge 3$ min. pH tomate $\approx$ 4,2-4,5 (limite) $\rightarrow$ stérilisation obligatoire (pas pasteurisation seule).
            \end{itemize}
        \item \textbf{Rédiger le mode opératoire :} Verbes d'action à l'infinitif, quantités, temps, températures, matériels, consignes sécurité/hygiène.
    \end{enumerate}
    \textbf{Réflexe Bac :} Distinguer \textbf{concentration} (évaporation eau) et \textbf{stérilisation} (traitement thermique final en contenant clos). La tomate est un produit \textbf{faiblement acide (pH > 4,5 possible)} : la stérilisation en autoclave (\SI{>115}{\celsius}) est \textbf{obligatoire} pour les conserves appertisées (boîtes/pots), contrairement à la confiture (pH < 4,5 + $a_w$ bas).
\end{methode}

\begin{exoresolu}[Rédaction d'un protocole HACCP simplifié pour pâte de tomate]
    \textbf{Énoncé :} Dans le cadre d'une unité de transformation villageoise à Dschang (Ouest-Cameroun), vous devez rédiger la fiche technique de fabrication du \textbf{double concentré de tomate (28-30\% Brix)} conditionné en pots verre \SI{200}{g}, stérilisés à l'autoclave. La matière première est la variété \emph{Roma} (rendement jus \SI{55}{\%}, Brix initial \SI{5}{\%}).
    \begin{enumerate}
        \item Proposer un schéma bloc des opérations unitaires en identifiant 3 Points Critiques de Contrôle (PCC/CCP).
        \item Calculer le rendement massique théorique (kg concentré / 100 kg tomates fraîches).
        \item Justifier la nécessité de la stérilisation à \SI{115}{\celsius} au lieu d'une simple pasteurisation à \SI{90}{\celsius}.
    \end{enumerate}
    \tcblower
    \textbf{Solution détaillée :}
    \begin{enumerate}
        \item \textbf{Schéma bloc et PCC :}
            \begin{center}
                \begin{tabular}{|l|l|l|}\hline
                \textbf{Étape} & \textbf{Opérations / Paramètres} & \textbf{PCC ?} \\ \hline
                Réception & Tri variétal (Roma), maturité rouge, T° < \SI{25}{\celsius} & Non \\ \hline
                Lavage & Eau potable + Hypochlorite \SI{50}{ppm}, \SI{5}{min}, égouttage & \textbf{OUI (PCC1 : Résidus chlorés / Charge microbienne initiale)} \\ \hline
                Échaudage & \SI{95}{\celsius}, \SI{45}{s}, refroidissement eau glacée & Non (Qualité) \\ \hline
                Émondage & Manuel / Machine (vapeur) & Non \\ \hline
                Broyage / Pulpage & Raffineur maille \SI{0,6}{mm} (élim. peaux/pépins) & Non \\ \hline
                \textbf{Concentration} & \textbf{Cuisson atmosphérique \SI{100-102}{\celsius}, agitation continue. Cible \textbf{\SI{30}{\%} Brix}.} & \textbf{OUI (PCC2 : Teneur en solides / $a_w$ / Stabilité)} \\ \hline
                Remplissage & Pots \SI{200}{g} stérilisés, T° produit \SI{>90}{\celsius}, espace tête \SI{6}{mm} & Non \\ \hline
                \textbf{Stérilisation} & \textbf{Autoclave \SI{115}{\celsius} (\SI{0,7}{bar} surpression), \SI{20}{min} (centre \SI{>110}{\celsius}).} & \textbf{OUI (PCC3 : Létalité $F_0$ / Stérilité commerciale)} \\ \hline
                Refroidissement & Eau chlorée \SI{30-40}{\celsius} $\rightarrow$ \SI{<40}{\celsius} rapide & Non \\ \hline
                Contrôle / Stockage & Vérif. vide, aspect, étiquette. Stockage \SI{<25}{\celsius} & Non \\ \hline
                \end{tabular}
            \end{center}
        \item \textbf{Calcul du rendement massique (Bilan matière sur solides totaux) :}
            \begin{itemize}
                \item Masse tomates fraîches $M_F = \SI{100}{kg}$.
                \item Masse solides totaux (ST) initiaux = $M_F \times \text{Brix}_F = 100 \times 0,05 = \SI{5}{kg}$ ST.
                \item Perte émondage/pulpage (peaux, pépins, eau) $\approx \SI{15}{\%}$ masse fraîche $\rightarrow$ Pulpe traitée $\approx \SI{85}{kg}$ (conservant \SI{5}{kg} ST).
                \item Concentration : Évaporation eau seule. Masse ST finale = Masse ST initiale (hypothèse pas de perte sucre) = \SI{5}{kg}.
                \item Cible Brix final = \SI{30}{\%} = 0,30.
                \item Masse concentré $M_C = \frac{\text{Masse ST}}{\text{Brix}_C} = \frac{5}{0,30} = \SI{16,67}{kg}$.
                \item \textbf{Rendement massique = $\frac{16,67}{100} \times 100 = \mathbf{16,7\,\%}$ (soit \SI{16,7}{kg} concentré / \SI{100}{kg} tomates fraîches).}
            \end{itemize}
            \emph{Remarque : Rendement industriel réel souvent 14-15\% dû aux pertes sucres (collage, écumes, rejets).}
        \item \textbf{Justification Stérilisation \SI{115}{\celsius} vs Pasteurisation \SI{90}{\celsius} :}
            \begin{itemize}
                \item \textbf{pH de la tomate :} \SIrange{4,2}{4,5} (variable variété, maturité, sol). \textbf{Seuil sécurité botulique = 4,5}. Tomate = produit \textbf{« limite »} (low-acid food borderline).
                \item \textit{Clostridium botulinum} (type A, B, E) ne germe/produit toxine \textbf{que si pH > 4,5 ET $a_w > 0,94$ ET anaérobiose}.
                \item Dans concentré \SI{30}{\%} Brix, $a_w \approx 0,91-0,92$ (inhibiteur pour \textit{C. botulinum}). \textbf{MAIS} : en surface pot, condensation $\rightarrow$ $a_w$ locale $\uparrow$, pH local $\uparrow$ (moins sucre). Risque développement en surface.
                \item \textbf{Règlementation (Codex Alimentarius, Arrêtés MINESEC/MINCOMMERCE) :} Conserves légumes (pH > 4,5 ou limite) $\rightarrow$ \textbf{Stérilisation appertisée} ($T \ge \SI{115}{\celsius}$, $F_0 \ge 3$ min) obligatoire pour stérilité commerciale (destruction spores \textit{C. botulinum}, $D_{121} \approx \SI{0,2}{min}$, $z = \SI{10}{\celsius}$).
                \item Pasteurisation \SI{90}{\celsius} détruit formes végétatives, \textbf{PAS les spores}. Risque botulisme mortel inacceptable.
            \end{itemize}
    \end{enumerate}
    \textbf{Conclusion :} Le protocole intègre la maîtrise des PCC (Hygiène, Concentration, Stérilisation). Le rendement théorique est de \SI{16,7}{\%}. La stérilisation à \SI{115}{\celsius} est une exigence réglementaire et toxicologique absolue pour ce produit faiblement acide conditionné hermétiquement.
\end{exoresolu}

\begin{methode}[Comparer l'efficacité de techniques de conservation : Analyse de données expérimentales]
    \textbf{Objectif :} Interpréter des tableaux/courbes (comptages microbiens, teneur en vitamines, qualité sensorielle, $a_w$, pH) pour classer les techniques.
    \begin{enumerate}
        \item \textbf{Identifier les indicateurs mesurés :}
            \begin{itemize}
                \item \textbf{Microbiologique :} UFC/g (Total, Levures/Moisissures, Coliformes, \textit{Salmonella}, \textit{Staph. aureus}). Seuil : \emph{Non détectable} ou < seuil règlementaire (ex. \SI{e3}{UFC/g} flore totale).
                \item \textbf{Physico-chimique :} $a_w$, pH, \%\ Brix, \% Humidité, Acidité titrable.
                \item \textbf{Nutritionnel :} Vitamine C (mg/100g), $\beta$-carotène, Polyphénols totaux. \textbf{Taux de rétention} $TR = \frac{Q_{\text{nut, final}}}{Q_{\text{nut, initial}}} \times 100$ (sur base matière première initiale).
                \item \textbf{Sensoriel :} Note hédonique (couleur, odeur, texture, goût) / 5 ou / 10.
            \end{itemize}
        \item \textbf{Analyser la cinétique (si données temporelles) :} Pente de dégradation (ordre 0 ou 1). Demi-vie $t_{1/2}$.
        \item \textbf{Comparer au temps $t$ (ex. 3 mois, 6 mois) :} Tableau récapitulatif.
        \item \textbf{Conclure par un compromis :} Meilleure conservation microbiologique $\neq$ meilleure qualité nutritionnelle/sensorielle. \textbf{Ex :} Séchage solaire (bon $a_w$, mauvais Vit C) vs Congélation (excellente Vit C, chaîne de froid coûteuse) vs Confiture (bon microbiologique, perte Vit C thermique, sucre ajouté).
    \end{enumerate}
    \textbf{Réflexe Bac :} Ne pas dire « La technique A est la meilleure ». Dire : « Pour le critère X, la technique A est supérieure (valeur) ; pour le critère Y, la technique B l'est. Au regard de l'objectif (ex. disponibilité vitamine C en saison sèche), la technique B est préconisée malgré... ».
\end{methode}

\begin{exoresolu}[Comparaison séchage solaire vs confiture vs congélation (Mangue)]
    \textbf{Énoncé :} Des mangues de variété \emph{Kent} (Vit C initiale \SI{45}{mg/100g} MF, $a_w=0,99$, pH 4,0) sont traitées selon 3 modes. Résultats après \textbf{6 mois} stockage (\SI{25}{\celsius} séchage/confiture, \SI{-18}{\celsius} congélation). Les teneurs en vitamines sont exprimées en mg/100g \textbf{équivalent matière première fraîche} pour permettre la comparaison des taux de rétention.
    \begin{center}
        \begin{tabular}{|l|c|c|c|}\hline
        \textbf{Paramètre} & \textbf{Séchage solaire (éclats)} & \textbf{Confiture (65°Brix)} & \textbf{Congélation (pulpes)} \\ \hline
        $a_w$ final & 0,45 & 0,82 & 0,99 (gelée) \\ \hline
        Flore totale (UFC/g) & $2 \times 10^3$ & $< 10$ & $< 10$ (après décongélation) \\ \hline
        Levures/Moisissures (UFC/g) & $5 \times 10^2$ & $< 10$ & $< 10$ \\ \hline
        Vitamine C résiduelle (mg/100g eq. MF) & 8,1 & 12,6 & 38,2 \\ \hline
        $\beta$-carotène résiduel (\% initial) & 65\% & 78\% & 92\% \\ \hline
        Note sensorielle globale / 10 & 6,5 (texture cuir) & 8,0 (goût sucré) & 9,0 (frais) \\ \hline
        Coût estimé (FCFA/kg produit fini) & 800 & 1200 & 2500 (énergie + chaîne froid) \\ \hline
        \end{tabular}
    \end{center}
    \begin{enumerate}
        \item Calculer le taux de rétention de la Vitamine C pour chaque technique.
        \item Analyser la stabilité microbiologique au regard des $a_w$ mesurés.
        \item Comparer la qualité nutritionnelle et sensorielle.
        \item Préconiser la technique la plus adaptée pour : (a) Un groupement féminin rural sans électricité ; (b) Un supermarché à Yaoundé cible « premium ».
    \end{enumerate}
    \tcblower
    \textbf{Solution détaillée :}
    \begin{enumerate}
        \item \textbf{Taux de rétention Vitamine C ($TR_{VitC} = \frac{C_{\text{final (eq. MF)}}}{C_{\text{initial}}} \times 100$) :}
            \begin{itemize}
                \item Initial $C_i = \SI{45}{mg/100g}$ MF.
                \item \textbf{Séchage :} $TR = \frac{8,1}{45} \times 100 = \mathbf{18,0\,\%}$. (Perte massive : oxydation thermique/UV + enzymes résiduelles + temps long).
                \item \textbf{Confiture :} $TR = \frac{12,6}{45} \times 100 = \mathbf{28,0\,\%}$. (Perte thermique cuisson + oxydation, mais protection par sucre/acide/pH).
                \item \textbf{Congélation :} $TR = \frac{38,2}{45} \times 100 = \mathbf{84,9\,\%}$. (Meilleure rétention : froid arrête enzymes/oxydation, pas de chaleur).
            \end{itemize}
        \item \textbf{Stabilité microbiologique vs $a_w$ :}
            \begin{itemize}
                \item \textbf{Séchage ($a_w=0,45$) :} $a_w < 0,60$ $\rightarrow$ \textbf{Arrêt total croissance microbienne} (bactéries, levures, moisissures). Flore résiduelle ($10^3$) = contaminants post-séchage (poussières, manipulation) ou survivants thermorésistants. \textbf{Stable.}
                \item \textbf{Confiture ($a_w=0,82$) :} $0,70 < a_w < 0,88$ $\rightarrow$ \textbf{Inhibe bactéries et levures}, limite moisissures (seuil $a_w \approx 0,70$). pH 3,2 + chaleur mise en pot $\rightarrow$ \textbf{Stérilité commerciale} (Flore $<10$). \textbf{Très stable.}
                \item \textbf{Congélation ($a_w=0,99$) :} $a_w$ non limitant. Stabilité due \textbf{uniquement à la température} (\SI{-18}{\celsius} $\rightarrow$ arrêt métabolisme). \textbf{Risque :} Rupture chaîne froid $\rightarrow$ reprise croissance rapide. Flore faible car blanchiment pré-congélation + froid.
            \end{itemize}
        \item \textbf{Qualité Nutritionnelle / Sensorielle :}
            \begin{itemize}
                \item \textbf{Nutrition :} Congélation $\gg$ Confiture $>$ Séchage (Vit C, Caroténoïdes). La confiture apporte sucre ajouté (\SI{60-65}{g/100g}) $\rightarrow$ densité énergétique haute, densité nutritionnelle basse.
                \item \textbf{Sensoriel :} Congélation (proche frais) $>$ Confiture (produit nouveau, apprécié) $>$ Séchage (texture altérée, couleur brunie, goût concentré/oxydé).
            \end{itemize}
        \item \textbf{Préconisation contextuelle :}
            \begin{itemize}
                \item \textbf{(a) Groupement féminin rural (sans électricité, saison sèche) :} \textbf{Séchage solaire amélioré (tunnel/séchoir hybride).}
                    \begin{itemize}
                        \item \emph{Pourquoi :} Investissement faible (bâches, grilles), énergie solaire gratuite, produit stable \textbf{sans chaîne de froid} ($a_w$ bas), transport facile (poids/volume $\downarrow$), génération revenus (vente éclats mangue séchée). 
                        \item \emph{Amélioration :} Blanchiment \SI{90}{\celsius}/\SI{2}{min} + trempage solution \SI{1}{\%} acide citrique + \SI{2}{\%} metabisulfite (ou jus citron) $\rightarrow$ inactivation PPO, protection Vit C/couleur $\rightarrow$ $TR_{VitC}$ passe \SI{18}{\%} $\rightarrow$ \SI{35-40}{\%}.
                    \end{itemize}
                \item \textbf{(b) Supermarché Yaoundé « Premium » :} \textbf{Congélation (IQF - Individually Quick Frozen) ou Purée surgelée.}
                    \begin{itemize}
                        \item \emph{Pourquoi :} Clientèle aisée, demande « naturel », « sans sucre ajouté », « haute qualité nutritionnelle (Vit C 85\%) », texture préservée. Chaîne de froid maîtrisée (logistique supermarché). Prix de vente élevé couvre coûts énergie/transport.
                    \end{itemize}
            \end{itemize}
    \end{enumerate}
    \textbf{Conclusion :} Il n'y a pas de « meilleure » technique absolue. Le choix est un \textbf{arbitrage multi-critères} : Séchage = Résilience/Accessibilité/Coût ; Confiture = Valorisation gustative/Stabilité ambiante ; Congélation = Excellence nutritionnelle/sensorielle/Coût énergétique.
\end{exoresolu}

\begin{methode}[Calculer le rendement et le coût de revient simplifié]
    \textbf{Objectif :} Quantifier la performance technico-économique d'une filière de transformation.
    \begin{enumerate}
        \item \textbf{Rendement massique ($R_m$) :} $R_m (\%) = \frac{\text{Masse Produit Fini (kg)}}{\text{Masse Matière Première Brute Entrante (kg)}} \times 100$.
            \begin{itemize}
                \item Base : MP brute (avec déchets : peaux, noyaux, queues). Intègre pertes inévitables.
                \item Bilan matière : $MP_{\text{brute}} = PF + \text{Pertes (peaux, noyaux, eau évaporée, rejets)}$.
            \end{itemize}
        \item \textbf{Rendement en matière sèche / Nutriment ($R_{nut}$) :} $R_{nut} = \frac{Q_{\text{nut, PF}}}{Q_{\text{nut, MP brute}}} \times 100$. Essentiel pour vitamines thermosensibles.
        \item \textbf{Coût de revient unitaire (CRU) :} $CRU = \frac{\text{Somme Coûts Fixes (amortis) + Coûts Variables}}{\text{Quantité Produit Fini}}$.
            \begin{itemize}
                \item \textbf{CV (Variables) :} MP (achat + transport), Sucre/Additifs, Emballages (pots, capsules, étiquettes), Énergie (bois, gaz, électricité), Main d'œuvre journalière, Eau, Traitement déchets.
                \item \textbf{CF (Fixes) :} Amortissement matériel (séchoir, marmite, pulpeur, autoclave) sur durée vie (ex. 5 ans), Loyer local, Assurance, Salaires permanents.
            \end{itemize}
        \item \textbf{Seuil de rentabilité / Prix de vente :} $PV > CRU + \text{Marge}$. Comparer $PV$ au prix marché.
    \end{enumerate}
    \textbf{Réflexe Bac :} Ne pas oublier les \textbf{pertes de transformation} (épluchage, noyautage, évaporation) dans le dénominateur du rendement (base MP brute). Distinguer \textbf{Coût Variable} (proportionnel au volume) et \textbf{Coût Fixe} (indépendant du volume). L'emballage (pot verre) est souvent le 1er poste de coût variable en confiture.
\end{methode}

\begin{exoresolu}[Analyse technico-économique : Confiture de tomate vs Pâte de tomate]
    \textbf{Énoncé :} Une unité traite \SI{1}{tonne} de tomates Roma (\SI{50}{FCFA/kg} à la ferme) par jour. Deux scénarios :
    \begin{itemize}
        \item \textbf{Scénario A : Confiture de tomate (sucre 70\%).} Rendement massique PF/MP brute = \textbf{45\%}. Pots \SI{300}{g} (\SI{120}{FCFA/unité}). Sucre \SI{400}{FCFA/kg}. Gaz \SI{5000}{FCFA/jour}. MO \SI{3000}{FCFA/jour}. Production jour = \textbf{450 kg} PF.
        \item \textbf{Scénario B : Double concentré (28°Brix).} Rendement massique PF/MP brute = \textbf{16\%}. Boîtes métal \SI{400}{g} (\SI{80}{FCFA/unité}). Gaz \SI{8000}{FCFA/jour} (concentration + stérilisation autoclave). MO \SI{4000}{FCFA/jour}. Production jour = \textbf{160 kg} PF.
    \end{itemize}
    Amortissement matériel commun : \SI{2000}{FCFA/jour}.
    \begin{enumerate}
        \item Calculer le Coût de Revient Unitaire (CRU) en FCFA/kg PF pour les 2 scénarios.
        \item Si prix de vente marché : Confiture \SI{2500}{FCFA/kg}, Concentré \SI{3500}{FCFA/kg}. Calculer la marge nette journalière.
        \item Discuter le choix stratégique au regard de la perte de Vitamine C (Confiture \SI{30}{\%} rétention, Concentré \SI{60}{\%} rétention) et de la durée de conservation (Confiture 18 mois ambiante, Concentré 24 mois ambiante).
    \end{enumerate}
    \tcblower
    \textbf{Solution détaillée :}
    \begin{enumerate}
        \item \textbf{Calcul CRU (FCFA/kg PF) :}
            \textbf{Données communes/jour :} MP = \SI{1000}{kg} $\times$ \SI{50}{FCFA} = \SI{50000}{FCFA}. Amortissement = \SI{2000}{FCFA}. Total Fixes = \SI{2000}{FCFA}.
            
            \textbf{Scénario A (Confiture) :}
            \begin{itemize}
                \item Production jour = \SI{450}{kg} PF (cohérence rendement 45\%).
                \item \textbf{Coût MP au kg PF :} $\frac{\text{Prix MP kg}}{\text{Rendement}} = \frac{50}{0,45} = \SI{111,1}{FCFA/kg PF}$.
                \item \textbf{Sucre :} Recette 70\% sucre / pulpe. 1 kg PF $\approx$ 588g pulpe + 412g sucre. Coût sucre = $0,412 \times 400 = \SI{164,8}{FCFA}$.
                \item \textbf{Emballage :} Pot \SI{300}{g}. Nb pots/kg = $1/0,3 = 3,33$. Coût = $3,33 \times 120 = \SI{400}{FCFA}$.
                \item \textbf{Energie/MO (Variables) au kg PF :} $\frac{5000+3000}{450} = \SI{17,8}{FCFA/kg}$.
                \item \textbf{Fixes au kg PF :} $\frac{2000}{450} = \SI{4,4}{FCFA/kg}$.
                \item \textbf{CRU$_A$ = 111,1 + 164,8 + 400 + 17,8 + 4,4 = \SI{698,1}{FCFA/kg} $\approx$ \SI{698}{FCFA/kg}.}
            \end{itemize}
            
            \textbf{Scénario B (Concentré) :}
            \begin{itemize}
                \item Production jour = \SI{160}{kg} PF (cohérence rendement 16\%).
                \item \textbf{Coût MP au kg PF :} $\frac{50}{0,16} = \SI{312,5}{FCFA/kg PF}$.
                \item \textbf{Additifs :} Sel \SI{1}{\%} (négligeable), Acide citrique \SI{0,1}{\%} (\SI{2000}{FCFA/kg}) $\rightarrow$ \SI{2}{FCFA/kg}. Negligeable.
                \item \textbf{Emballage :} Boîte \SI{400}{g}. Nb boîtes/kg = 2,5. Coût = $2,5 \times 80 = \SI{200}{FCFA}$.
                \item \textbf{Energie/MO au kg PF :} $\frac{8000+4000}{160} = \SI{75}{FCFA/kg}$.
                \item \textbf{Fixes au kg PF :} $\frac{2000}{160} = \SI{12,5}{FCFA/kg}$.
                \item \textbf{CRU$_B$ = 312,5 + 2 + 200 + 75 + 12,5 = \SI{602}{FCFA/kg}.}
            \end{itemize}
            
            \textbf{Résultat :} CRU Confiture = \SI{698}{FCFA/kg} > CRU Concentré = \SI{602}{FCFA/kg}. \emph{Le concentré coûte moins cher à produire au kg car l'emballage (boîte métal) est 2x moins cher que le pot verre, et pas de sucre, malgré le rendement 3x plus faible.}
        \item \textbf{Marge nette journalière :}
            \begin{align*}
                \text{Marge}_A &= (\text{PV}_A - \text{CRU}_A) \times Q_A = (2500 - 698) \times 450 = 1802 \times 450 = \mathbf{\SI{810\,900}{FCFA/jour}} \\
                \text{Marge}_B &= (\text{PV}_B - \text{CRU}_B) \times Q_B = (3500 - 602) \times 160 = 2898 \times 160 = \mathbf{\SI{463\,680}{FCFA/jour}}
            \end{align*}
            La confiture génère une marge absolue journalière \textbf{supérieure} (+75\%) grâce au volume produit 3x plus important, malgré un CRU plus élevé.
        \item \textbf{Discussion stratégique :}
            \begin{itemize}
                \item \textbf{Nutrition :} Concentré (Rétention Vit C 60\%) $\times$ 2 (reconstitution) $\approx$ 120\% Vit C tomate fraîche vs Confiture 30\% (sucre ajouté, chaleur longue). \textbf{Avantage Concentré.}
                \item \textbf{Conservation :} Concentré (24 mois, boîte métal robuste, empilable, transport facile) > Confiture (18 mois, pot verre cassant, lourd). \textbf{Avantage Concentré.}
                \item \textbf{Marché / Débouchés :} Confiture = niche « petit-déjeuner », concurrence forte (importations), sucre ajouté (mal vu santé). Concentré = \textbf{intrant culinaire universel} (sauces, plats quotidiens Cameroun), demande massive, non saisonnière, exportable (CEMAC).
                \item \textbf{Risque / Investissement :} Concentré nécessite autoclave (investissement, maintenance, opérateur formé, risque explosion). Confiture = technologie plus accessible (marmite feu bois).
            \end{itemize}
            \textbf{Conclusion :} Techniquement et nutritionnellement, le \textbf{double concentré (Scénario B) est supérieur}. Économiquement, la confiture (A) dégage plus de marge journalière brute si l'écoulement des 450 kg/jour est assuré. Le choix final dépend de la capacité d'investissement (autoclave) et de l'accès au marché de gros (concentré) vs détail local (confiture).
    \end{enumerate}
\end{exoresolu}

\begin{methode}[Identifier les facteurs d'altération et choisir le principe de conservation adapté]
    \textbf{Objectif :} Relier cause (facteur) $\rightarrow$ conséquence (altération) $\rightarrow$ principe de maîtrise (technique).
    \begin{enumerate}
        \item \textbf{Lister les 3 familles de facteurs :}
            \begin{enumerate}
                \item \textbf{Biotiques (Micro-organismes) :} Bactéries (altération, pathogènes), Levures (fermentation alcoolique), Moisissures (mycotoxines, pourriture). \emph{Conditions :} $a_w$, pH, T°, O$_2$, nutriments.
                \item \textbf{Biochimiques (Enzymes endogènes) :} PPO/POD (brunissement enzymatique), Pectinestérase/Polygalacturonase (ramollissement, dépectinisation), Lipoxygénase (goûts rances), Amylases. \emph{Actives T° ambiante, détruites T° > \SI{80-90}{\celsius} (blanchiment).}
                \item \textbf{Chimiques/Physiques :} Oxydation (air, lumière, métaux $\rightarrow$ perte Vit C, couleur, rancidité), Réaction de Maillard (sucre + AA $\rightarrow$ brunissement non enzymatique, perte lysine), Transpiration (perte eau, flétrissure), Choc mécanique (blessures = portes d'entrée microbes).
            \end{enumerate}
        \item \textbf{Associer le principe de conservation (Barrière / Hurdle Technology) :}
            \begin{center}
                \begin{tabular}{|l|l|l|}\hline
                \textbf{Facteur Cible} & \textbf{Principe (Barrière)} & \textbf{Technique Exemple} \\ \hline
                Micro-organismes (végétatifs) & Chaleur (Pasteurisation \SI{72-95}{\celsius}) & Jus pasteurisé, Confiture (cuisson) \\ \hline
                Spores bactériennes & Chaleur forte (Stérilisation \SI{>115}{\celsius}) + pH < 4,5 ou $a_w < 0,9$ & Conserves appertisées (tomate, haricots) \\ \hline
                Enzymes & Blanchiment (\SI{90-95}{\celsius}, \SI{1-3}{min}) & Préalable Séchage, Congélation, Conserves \\ \hline
                Micro-organismes / Enzymes & Froid (Réfrigération \SI{0-4}{\celsius} / Congélation \SI{-18}{\celsius}) & Fruits frais, Pulpes surgelées \\ \hline
                Micro-organismes (levures/moisissures) & Déshydratation ($a_w \downarrow$) : Séchage ($a_w<0,6$), Sucre ($a_w<0,85$), Sel & Mangue séchée, Confiture, Tomates séchées \\ \hline
                Micro-organismes (bactéries) & Acidification (pH < 4,5) : Lactofermentation, Ajout acide & Cornichons, Confiture (citron), Jus acidifiés \\ \hline
                Oxydation (O$_2$) & Exclusion air (Vide, Gaz inerte, Enrobage, Antioxidants) & Conditionnement sous vide, Sulfitation, Vit C ajoutée \\ \hline
            \end{tabular}
            \end{center}
        \item \textbf{Combiner les barrières (Technologie des barrières / Hurdle Technology) :} Produit stable = Somme effets barrières. Ex: Confiture = Chaleur + $a_w$ (sucre) + pH (acide) + Anaérobiose (vide). Séchage solaire = Chaleur modérée + $a_w$ bas + UV (germicide) + Air (oxydation $\uparrow$ $\rightarrow$ besoin pré-traitement antioxydant).
    \end{enumerate}
    \textbf{Réflexe Bac :} Citer le couple \textbf{Facteur $\leftrightarrow$ Principe $\leftrightarrow$ Technique}. Ex: « Pour limiter le brunissement enzymatique (PPO), on applique un blanchiment (\SI{95}{\celsius}, \SI{2}{min}) avant le séchage ». Ne pas dire « on sèche pour tuer les enzymes » (faux, séchage $\neq$ inactivation enzymatique complète, réactivation possible à l'humidité).
\end{methode}

\begin{exoresolu}[Diagnostic d'altération et proposition de solution (Cas Mangue)]
    \textbf{Énoncé :} Un producteur de mangues \emph{Brooks} à Njombé (Moungo) observe sur son stock récolté il y a 5 jours (stockage ambiant \SI{28}{\celsius}, humidité relative 85\%) :
    \begin{itemize}
        \item 30\% des fruits présentent des taches brunes molles en expansion (odeur fermentée).
        \item 20\% des fruits sont ridés, mous au toucher, sans tache visible.
        \item Le reste semble sain mais la peau noircit au niveau des contacts.
    \end{itemize}
    \begin{enumerate}
        \item Identifier les 3 types d'altération observés (agent causal + mécanisme).
        \item Proposer 3 techniques de conservation adaptées pour valoriser ce stock \emph{avant} perte totale, en justifiant le choix par les facteurs maîtrisés.
        \item Pour la technique « Séchage », préciser le pré-traitement indispensable pour éviter le noircissement des éclats.
    \end{enumerate}
    \tcblower
    \textbf{Solution détaillée :}
    \begin{enumerate}
        \item \textbf{Diagnostic des altérations :}
            \begin{enumerate}
                \item \textbf{Pourriture molle brune + odeur fermentée :} \textbf{Champignons moisissures (ex. \textit{Rhizopus}, \textit{Aspergillus}, \textit{Penicillium}) ou Bactéries (\textit{Erwinia}, \textit{Pseudomonas}).} Mécanisme : Pénétration par blessures (stade climactérique = sensibilité max), sécrétion enzymes pectinolytiques (ramollissement tissus) + métabolisme anaérobie (fermentation sucres $\rightarrow$ \ce{CO2}, éthanol, acides $\rightarrow$ odeur). Facteurs favorisants : HR 85\% ($a_w$ surface élevée), T° \SI{28}{\celsius} (optimum), maturité avancée.
                \item \textbf{Rides, mollesse sans tache :} \textbf{Sénescence / Transpiration excessive (Perte d'eau).} Mécanisme : Respiration climactérique intense (pic éthylène/\ce{CO2}) + Transpiration cuticulaire $\rightarrow$ Perte turgescence cellulaire (flétrissure), consommation réserves (amidon $\rightarrow$ sucres puis respiration), ramollissement enzymatique paroi (pectinestérase). Facteurs : T° élevée (vitesse respiration $\times$ 2-3 par \SI{10}{\celsius}), HR 85\% insuffisante pour stopper transpiration (déficit de saturation), absence barrière physique (cire, film).
                \item \textbf{Noircissement peau aux contacts :} \textbf{Brunissement enzymatique (Polyphénol Oxydase - PPO).} Mécanisme : Chocs/pressions (contacts) $\rightarrow$ rupture vacuoles $\rightarrow$ mise en contact PPO (cytoplasme) + Substrats phénoliques (vacuole) + \ce{O2} (air) $\rightarrow$ Oxydation $\rightarrow$ Quinones $\rightarrow$ Polymérisation $\rightarrow$ Mélanines (brun/noir). Facteur déclencheur : Blessure mécanique.
            \end{enumerate}
        \item \textbf{3 Techniques de valorisation d'urgence (Tri strict préalable obligatoire : éliminer pourris) :}
            \begin{enumerate}
                \item \textbf{Transformation en Confiture / Gelée (Cuisson + Sucre + Acide) :}
                    \begin{itemize}
                        \item \emph{Justification :} \textbf{Chaleur} (\SI{105}{\celsius}) détruit microbes/enzymes (PPO) $\rightarrow$ arrête pourriture et brunissement. \textbf{Sucre} ($a_w \downarrow$) empêche recontamination. \textbf{Acide (citron)} pH < 4,5 sécurise (botulisme) + fixe couleur + gélification. \textbf{Conditionnement hermétique} exclut \ce{O2}.
                        \item \emph{Avantage :} Valorise fruits mous (pulpe), forte valeur ajoutée, stockage ambiant long.
                    \end{itemize}
                \item \textbf{Séchage solaire / hybride (Éclats ou cuir de fruits) :}
                    \begin{itemize}
                        \item \emph{Justification :} \textbf{Déshydratation} ($a_w < 0,60$) arrête \textbf{toute activité microbienne} (pourriture) et \textbf{enzymatique} (sénescence, PPO si inactivée avant). Réduit volume/poids (transport). Stable ambiant.
                        \item \emph{Condition :} Tri drastique (pas de fruit pourri $\rightarrow$ mycotoxines non détruites par séchage).
                    \end{itemize}
                \item \textbf{Pulpe pasteurisée / Jus (Conditionnement chaud en bouteilles/briques) :}
                    \begin{itemize}
                        \item \emph{Justification :} \textbf{Pasteurisation} (\SI{90-95}{\celsius}, \SI{15-30}{s}) détruit formes végétatives + inactivation PPO. \textbf{Acidification} (pH < 4,0 jus citron) + \textbf{Remplissage à chaud} + \textbf{Capsulage} $\rightarrow$ Vide + Anaérobiose. Conservation 6-12 mois ambiant.
                        \item \emph{Avantage :} Processus rapide (flux tendu), valorise gros volumes, produit prêt à boire.
                    \end{itemize}
            \end{enumerate}
        \item \textbf{Pré-traitement indispensable pour Séchage (anti-noircissement PPO) :}
            \begin{itemize}
                \item \textbf{Blanchiment :} Immersion éclats (\SI{1-1,5}{cm}) dans eau \SI{95-98}{\celsius} pendant \textbf{\SI{2-3}{min}}. Refroidissement immédiat eau froide. \emph{Effet :} Inactivation thermique PPO/POD (thermolables). Arrêt brunissement enzymatique.
                \item \textbf{(Alternative/Complément) Trempage antioxydant :} Solution \textbf{\SI{0,5-1}{\%} Acide citrique} + \textbf{\SI{0,1-0,2}{\%} Métabisulfite de sodium (\ce{SO2})} ou \SI{2-3}{\%} Jus de citron + \SI{5}{\%} Sirop sucre (osmose protection), pendant \SI{10}{min} avant égouttage/séchage.
                    \begin{itemize}
                        \item Acide : pH $\downarrow$ $\rightarrow$ PPO moins active + chélation \ce{Cu^{2+}} (cofacteur PPO).
                        \item \ce{SO2} : Inhibiteur puissant PPO + antioxydant (réduit quinones $\rightarrow$ phénols) + antimicrobien.
                    \end{itemize}
                \item \textbf{Résultat attendu :} Éclats séchés couleur jaune-orangé vif (rétention $\beta$-carotène/Vit C améliorée), pas de brunissement au stockage.
            \end{itemize}
    \end{enumerate}
    \textbf{Conclusion :} L'altération est multifactorielle (microbienne, physiologique, biochimique). Le tri est l'étape critique n°1. La confiture et le jus pasteurisé utilisent la chaleur comme barrière principale ; le séchage utilise l'$a_w$. Le blanchiment est \textbf{obligatoire} avant séchage pour la qualité couleur/nutrition.
\end{exoresolu}

\begin{methode}[Argumenter les intérêts économiques, nutritionnels et environnementaux]
    \textbf{Objectif :} Structurer une argumentation écrite (type dissertation Bac) en 3 piliers, avec des indicateurs chiffrés et du contexte camerounais.
    \begin{enumerate}
        \item \textbf{Intérêt Économique (Revenus, Emploi, Souveraineté) :}
            \begin{itemize}
                \item \textbf{Valorisation excédents / Réduction pertes post-récolte :} Pertes mangue/tomate \SI{30-50}{\%} (FAO/MINADER). Transformation = revenu au lieu de déchet.
                \item \textbf{Valorisation économique (Value Addition) :} Prix kg fruit frais \SI{100-300}{FCFA} (saison) $\rightarrow$ Confiture \SI{2500}{FCFA/kg}, Concentré \SI{3500}{FCFA/kg}, Séchée \SI{4000}{FCFA/kg}. Facteur multiplicateur 10-40x.
                \item \textbf{Création d'emplois :} Filière transformation (cueillette, transport, main d'œuvre usine, conditionnement, commercialisation). Jeunes / Femmes (GIE).
                \item \textbf{Substitution importations :} Concentré tomate importé (Chine/Italie) \SI{>100\,000}{t/an} (FCFA milliards). Production locale = économie devises, balance commerciale.
                \item \textbf{Disponibilité hors-saison / Stabilisation prix :} Lissage offre/demande, prix accessible aux consommateurs urbains toute l'année.
            \end{itemize}
        \item \textbf{Intérêt Nutritionnel (Sécurité alimentaire, Santé) :}
            \begin{itemize}
                \item \textbf{Disponibilité micronutriments :} Vitamine C, $\beta$-carotène (Provitamine A), Fibres, Potassium, Polyphénols, Lycopène (tomate). Lutte contre carences (Vit A : cécité, immunité ; Fer : anémie).
                \item \textbf{Rétention maîtrisée :} Congélation/Séchage amélioré/Concentré (court/fort) $>$ Confiture (long/doux + sucre). \textbf{Attention :} Confiture = sucre ajouté (\SI{60-65}{g/100g}) $\rightarrow$ risque surpoids/diabète si surconsommation. \textbf{Message :} « Confiture $\neq$ portion fruit frais » (repère PNNS).
                \item \textbf{Sécurité sanitaire :} Maîtrise hygiène (HACCP) $\rightarrow$ élimination mycotoxines (aflatoxines sur mangue pourrie), pathogènes. Eau potable, emballages sains.
            \end{itemize}
        \item \textbf{Intérêt Environnemental (Durabilité, Économie circulaire) :}
            \begin{itemize}
                \item \textbf{Réduction déchets organiques / Gaz à effet de serre :} Fruits pourris en décharge $\rightarrow$ \ce{CH4} (GES puissant, GWP 28), lixiviats. Transformation = 0 déchet fruit.
                \item \textbf{Valorisation coproduits (Bioraffinerie) :}
                    \begin{itemize}
                        \item Peaux/Noyaux mangue $\rightarrow$ Pectine, Polyphénols, Huile de noyau (cosmétique), Aliment bétail, Compost, Biogaz, Biochar.
                        \item Peaux/Pépins tomate (pomace) $\rightarrow$ Lycopène (colorant/nutraceutique), Fibres, Aliment volaille/porc, Compost.
                    \end{itemize}
                \item \textbf{Efficacité énergétique :} Séchage solaire (énergie renouvelable) vs Congélation (électricité/charbon). Optimisation procédés (évaporateurs multi-effets, récupération chaleur).
                \item \textbf{Emballages :} Verre (recyclable, réutilisable consigne) vs Plastique/Boîte métal (recyclage difficile au Cameroun). Éco-conception.
            \end{itemize}
        \item \textbf{Synthèse Bac :} Utiliser la structure : « Sur le plan économique... Sur le plan nutritionnel... Sur le plan environnemental... En conclusion, la transformation s'inscrit dans une démarche de \textbf{développement durable} (Objectifs ODD 2, 8, 12, 13). »
    \end{enumerate}
\end{methode}

\begin{exoresolu}[Argumentation développement durable : Filière Mangue au Cameroun]
    \textbf{Énoncé :} Le Cameroun produit environ \SI{200\,000}{tonnes} de mangues/an (régions : Littoral, Sud-Ouest, Adamaoua, Nord). Les pertes post-récolte sont estimées à \SI{40}{\%}. L'importation de concentré de tomate coûte \SI{30}{milliards} FCFA/an. 
    Rédigez une note de synthèse (20 lignes max) argumentant l'intérêt de développer des unités locales de transformation de la mangue (confiture, séché, jus) et de la tomate (concentré), structurée selon les 3 piliers du développement durable.
    \tcblower
    \textbf{Solution détaillée (Modèle de rédaction Bac) :}
    
    \textbf{Le développement de la transformation locale des fruits de saison (mangue, tomate) au Cameroun constitue un levier majeur de développement durable, articulant enjeux économiques, nutritionnels et environnementaux.}
    
    \textbf{1. Sur le plan économique, c'est un vecteur de souveraineté et d'emploi.}
    Les pertes post-récolte (\SI{40}{\%} soit \SI{80\,000}{t} mangues/an) représentent un manque à gagner colossal pour les producteurs. La transformation (confiture, séché, jus, concentré) opère une \textbf{valorisation économique massive} : le kg de mangue fraîche (\SI{150}{FCFA} en saison) devient confiture (\SI{2\,500}{FCFA/kg}) ou fruit séché (\SI{4\,000}{FCFA/kg}), multipliant la valeur ajoutée par 15 à 25. Pour la tomate, la production de double concentré local (cible \SI{28-30}{\%} Brix) permettrait de substituer une partie des \SI{30}{milliards} FCFA/an d'importations, préservant les devises et créant des emplois ruraux (GIE féminins, PME agroalimentaires) tout en stabilisant les prix urbains hors saison.
    
    \textbf{2. Sur le plan nutritionnel, c'est une réponse à l'insécurité micronutritionnelle.}
    La mangue et la tomate sont des sources stratégiques de \textbf{Vitamine C}, \textbf{$\beta$-carotène (Provitamine A)} et \textbf{lycopène} (antioxydant). Des techniques adaptées (congélation, séchage solaire amélioré avec blanchiment/sulfitation, concentration courte sous vide) permettent de \textbf{retenir 60 à 85\% de la Vitamine C} contre 20-30\% pour une confiture traditionnelle. Cela assure un apport en micronutriments essentiels pendant la saison sèche (période de soudure), luttant contre l'hypovitaminose A (cécité, mortalité infantile) et renforçant l'immunité (contexte VIH/Sida, Paludisme). La maîtrise de l'hygiène (HACCP, $a_w$, pH, stérilisation) garantit l'innocuité (absence mycotoxines, botulisme).
    
    \textbf{3. Sur le plan environnemental, c'est un modèle d'économie circulaire bas-carbone.}
    La transformation évite la mise en décharge de \SI{80\,000}{t} de déchets organiques fermentescibles, source de \textbf{méthane (\ce{CH4})}, gaz à effet de serre 28x plus puissant que le \ce{CO2}. Elle permet la \textbf{bioraffinerie} : peaux/noyaux de mangue (pectines, huile cosmétique, biogaz) et pomace de tomate (lycopène, aliment bétail, compost) deviennent des ressources. L'usage du \textbf{séchage solaire} (énergie gratuite, décarbonée) pour la mangue et l'optimisation énergétique des concentrateurs (évaporateurs multi-effets) pour la tomate réduisent l'empreinte carbone comparée aux importations (transport maritime/routier) ou à la congélation intensive.
    
    \textbf{En conclusion,} investir dans des unités semi-industrielles de proximité, dotées d'autoclaves pour la tomate (pH limite) et de séchoirs hybrides pour la mangue, appuyées par la formation HACCP et l'accès au financement (Fonds Agricole), est la condition \emph{sine qua non} pour transformer l'abondance saisonnière en sécurité alimentaire durable et en croissance économique inclusive au Cameroun.
\end{exoresolu}

\begin{attention}[Les 5 erreurs qui coûtent des points au Bac]
    \begin{enumerate}
        \item \textbf{Confondre Pasteurisation et Stérilisation (et leurs cibles) :} 
            \emph{Erreur :} « On stérilise la confiture à \SI{100}{\celsius} » ou « On pasteurise la tomate en boîte à \SI{95}{\celsius} ».
            \textbf{Vérité :} Confiture (pH < 4,5 + $a_w$ bas) = \textbf{Pasteurisation} (cuisson \SI{105}{\celsius} + retournement). Tomate (pH $\approx$ 4,3-4,5, $a_w$ élevé) = \textbf{Stérilisation appertisée} \textbf{obligatoire} en autoclave \SI{\ge 115}{\celsius} ($F_0 \ge 3$ min) pour détruire spores \textit{Clostridium botulinum}. Ne pas citer \SI{115}{\celsius} pour la tomate = \textbf{0 point} sur la technique.
        \item \textbf{Oublier le Blanchiment avant Séchage / Congélation :} 
            \emph{Erreur :} « On coupe la mangue et on met au soleil » ou « On congèle la tomate crue ».
            \textbf{Vérité :} Les enzymes (PPO, POD, Pectinestérase, Lipoxygénase) sont \textbf{thermolables mais actives à T° ambiante et même au froid lent}. Sans blanchiment (\SI{90-95}{\celsius}, \SI{1-3}{min}) : \textbf{Brunissement enzymatique} (couleur brune, perte Vit C), \textbf{Ramollissement} (pectinestérase), \textbf{Goûts rances} (lipoxygénase). Le blanchiment est une \textbf{étape critique (PCC)} qualité.
        \item \textbf{Calculer le rendement sur la mauvaise base (MP brute vs MP nette) :} 
            \emph{Erreur :} Rendement = Masse PF / Masse MP nette (épluchée). 
            \textbf{Vérité :} Le rendement industriel/économique standard se calcule sur \textbf{Matière Première Brute (réception)} : $R_m = \frac{M_{PF}}{M_{MP\,\text{brute}}} \times 100$. Cela intègre les pertes inévitables (peaux, noyaux, queues, lavage). Pour la tomate concentré : rendement \SI{16}{\%} sur brute (pas \SI{30}{\%} sur pulpe).
        \item \textbf{Dire « Le séchage tue les microbes » au lieu de « Le séchage inhibe la croissance » :} 
            \emph{Erreur :} « Le séchage stérilise le produit ».
            \textbf{Vérité :} Le séchage ($a_w < 0,6$) est \textbf{bactériostatique et fongistatique} (arrêt métabolisme). Il \textbf{ne détruit pas} les spores (thermorésistantes) ni les mycotoxines déjà produites (thermostables). Si le fruit était pourri \emph{avant} séchage, les toxines restent. \textbf{Tri rigoureux = PCC n°1}. La réhumidification ($a_w \uparrow$) réactive la flore résiduelle.
        \item \textbf{Négliger l'impact du sucre ajouté en confiture sur le bilan nutritionnel :} 
            \emph{Erreur :} « La confiture conserve les vitamines du fruit » ou « C'est une portion de fruit ».
            \textbf{Vérité :} Cuisson longue + pH acide + \ce{O2} $\rightarrow$ \textbf{Perte Vit C > 70\%}. Surtout : \textbf{Sucre ajouté \SI{60-65}{g/100g}}. La confiture est un \textbf{produit sucré}, pas un fruit. Au Bac, il faut nuancer : « Intérêt énergétique et gustatif, apport en fibres/minéraux/caroténoïdes (stables), mais \textbf{faible densité nutritionnelle en Vit C} et \textbf{apport sucre libre élevé} (recommandation OMS < 10\% AET) ». Ne pas dire « bon pour la santé » sans réserve.
    \end{enumerate}
\end{attention}

\newpage

\coursec{Exercices}

\courssub{Je vérifie mes connaissances}

\begin{exercice}[QCM : les facteurs d'altération]
Pour chaque question, choisis la (ou les) bonne(s) réponse(s).
\begin{enumerate}
\item L'altération microbienne d'une mangue mûre est favorisée par :
\begin{enumerate}
\item une activité de l'eau ($a_w$) proche de 1 ;
\item une température comprise entre 20 \textcelsius{} et 40 \textcelsius{} ;
\item un pH très acide (pH $< 3$) ;
\item une teneur élevée en eau.
\end{enumerate}
\item Le brunissement d'une mangue coupée à l'air libre est dû :
\begin{enumerate}
\item à des bactéries lactiques ;
\item à l'action d'enzymes (polyphénoloxydases) en présence de dioxygène ;
\item à la fermentation alcoolique ;
\item à une réaction d'oxydation enzymatique.
\end{enumerate}
\item La stérilisation d'une conserve de tomate vise à détruire :
\begin{enumerate}
\item uniquement les enzymes ;
\item les micro-organismes et leurs spores ;
\item uniquement les levures ;
\item les vitamines.
\end{enumerate}
\item L'ajout de sucre en grande quantité dans une confiture conserve le fruit parce qu'il :
\begin{enumerate}
\item tue directement les bactéries ;
\item abaisse l'activité de l'eau disponible pour les micro-organismes ;
\item augmente le pH ;
\item détruit les spores.
\end{enumerate}
\end{enumerate}
\end{exercice}

\begin{exercice}[Vrai ou faux justifié]
Pour chacune des affirmations suivantes, indique si elle est vraie ou fausse, puis justifie ta réponse en une ou deux phrases.
\begin{enumerate}
\item « Le séchage de la mangue tue tous les micro-organismes présents sur le fruit. »
\item « La pasteurisation d'un jus de mangue à 70 \textcelsius{} pendant 15 min détruit les spores bactériennes. »
\item « Le blanchiment des tomates avant transformation inactive les enzymes responsables du brunissement. »
\item « Une confiture correctement sucrée et fermée hermétiquement peut se conserver plusieurs mois à température ambiante. »
\item « L'oxydation de la vitamine C est accélérée par la lumière et la chaleur. »
\end{enumerate}
\end{exercice}

\begin{exercice}[Définitions et applications]
Définis chacun des termes suivants en une phrase, puis cite pour chacun un exemple d'application concrète à la mangue ou à la tomate :
\begin{enumerate}
\item la pasteurisation ;
\item la stérilisation ;
\item le blanchiment (ou ébouillantage) ;
\item l'activité de l'eau ($a_w$) ;
\item la déshydratation osmotique.
\end{enumerate}
\end{exercice}

\begin{exercice}[Calculs directs : rendements et humidité]
Une coopérative de la localité de Mora (Extrême-Nord) récolte 250 kg de mangues fraîches contenant 82 \% d'eau en masse.
\begin{enumerate}
\item Calcule la masse d'eau et la masse de matière sèche contenues dans ces mangues.
\item Après séchage, les mangues ne contiennent plus que 15 \% d'eau en masse, la matière sèche étant conservée intégralement. Calcule la masse de mangues séchées obtenue.
\item Une ménagère prépare une confiture : elle mélange 4 kg de pulpe de mangue avec 3 kg de sucre ; après cuisson, il reste 5,6 kg de confiture. Calcule la masse d'eau évaporée pendant la cuisson.
\end{enumerate}
\end{exercice}

\courssub{Je m'entraîne}

\begin{exercice}[Logigramme de fabrication de la pâte de tomate]
Une unité artisanale de transformation de tomates à Bafoussam réalise les opérations suivantes : tri et lavage des tomates ; ébouillantage à 95 \textcelsius{} pendant 3 min ; égouttage et broyage ; tamisage (élimination des peaux et pépins) ; concentration par chauffage jusqu'à 28 \% de matière sèche ; remplissage à chaud des bocaux ; stérilisation des bocaux fermés à 100 \textcelsius{} pendant 45 min ; refroidissement et stockage.
\begin{enumerate}
\item Présente ces opérations sous forme d'un logigramme (étapes dans des rectangles reliés par des flèches).
\item Pour chaque étape, indique le facteur d'altération (micro-organismes, enzymes, oxygène de l'air) qu'elle vise à limiter, en justifiant brièvement.
\item Pourquoi le remplissage des bocaux se fait-il « à chaud » ?
\end{enumerate}
\end{exercice}

\begin{exercice}[Comparaison multicritères de trois procédés]
On applique trois techniques de conservation à une même variété de mangue. Le tableau ci-dessous rassemble les résultats obtenus.

\begin{center}
\begin{tabularx}{\linewidth}{|l|X|X|X|}\hline
\rowcolor{popblueL} Critère & Mangue séchée & Confiture de mangue & Conserve de mangue au sirop \\ \hline
Durée de conservation & 6 à 12 mois & 12 mois (pot fermé) & 2 à 3 ans \\ \hline
Teneur en vitamine C conservée & 30 \% & 20 \% & 45 \% \\ \hline
Coût de l'équipement & faible (séchoir solaire) & faible (marmite) & élevé (autoclave) \\ \hline
Consommation d'énergie & très faible & moyenne & élevée \\ \hline
Risque de recontamination après ouverture & faible & élevé & moyen \\ \hline
\end{tabularx}
\end{center}

\begin{enumerate}
\item À partir du tableau, compare l'efficacité des trois procédés selon les critères de durée de conservation et de qualité nutritionnelle.
\item Quel procédé recommanderais-tu à un groupement de femmes d'un village non électrifié ? Argumente ta réponse en mobilisant au moins trois critères du tableau.
\item Rédige une conclusion de cinq lignes maximum montrant qu'il n'existe pas de technique « idéale » mais une technique adaptée à chaque contexte.
\end{enumerate}
\end{exercice}

\begin{exercice}[La confiture qui moisit après ouverture]
Un pot de confiture de mangue préparé en avril reste parfaitement stable pendant 8 mois dans le placard. Une semaine après son ouverture, des taches de moisissure apparaissent à la surface.
\begin{enumerate}
\item Rappelle les deux facteurs qui assuraient la stabilité de la confiture dans le pot fermé.
\item Explique, en mobilisant les facteurs d'altération (micro-organismes, activité de l'eau, oxygène), pourquoi des moisissures se développent après ouverture.
\item Propose deux mesures pratiques permettant de ralentir cette altération après ouverture du pot.
\end{enumerate}
\end{exercice}

\begin{exercice}[Exploitation d'un tableau de survie microbienne]
On inocule des tranches de mangue avec une levure, puis on les soumet à différents traitements. Le tableau donne le nombre de levures vivantes (en unités formant colonies, UFC) par gramme de mangue après 7 jours de stockage à 25 \textcelsius.

\begin{center}
\begin{tabularx}{\linewidth}{|l|X|}\hline
\rowcolor{popblueL} Traitement & Levures (UFC/g) après 7 jours \\ \hline
Témoin (aucun traitement) & $2{,}0 \times 10^{6}$ \\ \hline
Blanchiment 2 min à 90 \textcelsius{} & $5{,}0 \times 10^{4}$ \\ \hline
Blanchiment + séchage ($a_w = 0{,}55$) & $1{,}0 \times 10^{2}$ \\ \hline
Blanchiment + séchage + emballage étanche & $< 10$ \\ \hline
\end{tabularx}
\end{center}

\begin{enumerate}
\item Calcule, pour chaque traitement, le facteur de réduction du nombre de levures par rapport au témoin.
\item Montre que le blanchiment seul ne suffit pas à conserver durablement la mangue.
\item Explique le rôle complémentaire du séchage et de l'emballage étanche.
\end{enumerate}
\end{exercice}

\courssub{Je me prépare au Bac}

\begin{exercice}[Type Bac : activité de l'eau et séchage de la mangue]
Dans la plaine du Logone, les pertes post-récolte de mangues atteignent 40 \% en pleine saison. Pour comprendre comment le séchage permet de conserver ce fruit, on étudie en laboratoire la croissance d'une moisissure (\emph{Aspergillus niger}) sur du gelose additionné de pulpe de mangue, en faisant varier l'activité de l'eau ($a_w$) du milieu. Les résultats sont consignés dans le tableau ci-dessous.

\begin{center}
\begin{tabularx}{\linewidth}{|l|X|X|X|X|X|X|}\hline
\rowcolor{popblueL} $a_w$ du milieu & 0{,}98 & 0{,}95 & 0{,}90 & 0{,}80 & 0{,}70 & 0{,}60 \\ \hline
Diamètre de la colonie après 5 jours (mm) & 42 & 35 & 22 & 9 & 2 & 0 \\ \hline
\end{tabularx}
\end{center}

On rappelle que la mangue fraîche a une $a_w$ d'environ 0{,}98 et que la teneur en eau de la pulpe fraîche est de 82 \% en masse. On admet qu'une $a_w$ inférieure ou égale à 0{,}60 correspond, pour la mangue, à une humidité résiduelle d'environ 15 \% en masse.

\begin{enumerate}
\item Trace, sur papier millimétré, la courbe du diamètre de la colonie en fonction de l'activité de l'eau. Échelles : 1 cm pour 0{,}05 unité de $a_w$ et 1 cm pour 5 mm de diamètre.
\item Analyse la courbe obtenue : comment varie la croissance de la moisissure lorsque l'activité de l'eau diminue ? Détermine graphiquement la valeur de $a_w$ en dessous de laquelle la croissance devient nulle.
\item Explique, au niveau cellulaire, pourquoi une faible activité de l'eau empêche le développement des micro-organismes (on fera intervenir la disponibilité de l'eau et les phénomènes osmotiques).
\item En utilisant les données de l'énoncé, montre que le séchage de la mangue jusqu'à une humidité résiduelle de 15 \% permet sa conservation. Calcule la masse d'eau qu'il faut éliminer à partir de 100 kg de mangues fraîches pour atteindre ce taux d'humidité résiduel (la matière sèche est conservée).
\item Le séchage solaire à l'air libre expose les tranches de mangue aux poussières et aux insectes. Propose une amélioration technique simple, adaptée au contexte rural camerounais, qui rende le séchage plus sûr et plus rapide, en justifiant ta proposition.
\end{enumerate}
\end{exercice}

\begin{exercice}[Type Bac : traitements thermiques du jus de mangue]
Une élève de Terminale D réalise, dans le cadre d'un projet scolaire à l'instar des unités de transformation de Ngaoundéré, trois lots de jus de mangue préparés dans des conditions identiques (mêmes fruits, même matériel lavé) :
\begin{itemize}
\item lot A : jus non chauffé, mis en bouteille fermée ;
\item lot B : jus chauffé à 70 \textcelsius{} pendant 15 min, mis en bouteille fermée à chaud ;
\item lot C : jus porté à 100 \textcelsius{} pendant 20 min dans des bouteilles fermées.
\end{itemize}
Les trois lots sont conservés à température ambiante (environ 28 \textcelsius). Les observations après 10 jours sont les suivantes :

\begin{center}
\begin{tabularx}{\linewidth}{|l|X|X|X|}\hline
\rowcolor{popblueL} Lot & Aspect & Odeur & Dénombrement microbien (UFC/mL) \\ \hline
A & trouble, présence de gaz (bouteille gonflée) & odeur alcoolique et acide & $8 \times 10^{7}$ \\ \hline
B & limpide, pas de gaz & normale & $3 \times 10^{2}$ \\ \hline
C & limpide, pas de gaz & normale, goût légèrement « cuit » & $< 10$ \\ \hline
\end{tabularx}
\end{center}

\begin{enumerate}
\item Formule le problème scientifique que cette expérience permet de résoudre, puis dégage l'hypothèse qu'elle teste.
\item Interprète les observations du lot A : quel type de micro-organismes s'est développé et quel type de fermentation explique la production de gaz et l'odeur ? Écris l'équation bilan de la fermentation alcoolique du glucose.
\item Nomme les procédés de conservation correspondant respectivement aux lots B et C, et précise ce qui les distingue (température, durée, micro-organismes détruits).
\item Explique pourquoi les lots B et C restent sains après 10 jours alors que le lot A s'altère. Pourquoi le lot B présente-t-il encore quelques micro-organismes vivants ?
\item La vitamine C est thermosensible : sa teneur diminue de 15 \% dans le lot B et de 40 \% dans le lot C. Sachant que le jus frais contient 30 mg de vitamine C pour 100 mL, calcule la teneur en vitamine C des lots B et C, puis discute le compromis entre sécurité microbiologique et qualité nutritionnelle.
\item Rédige une conclusion argumentée indiquant quel procédé tu conseillerais à une petite unité artisanale qui vend son jus sur les marchés locaux dans la semaine qui suit la production.
\end{enumerate}
\end{exercice}

\begin{exercice}[Type Bac : transformation de la tomate et filière de conservation]
Chaque année, pendant la grande saison (décembre à mars), les marchés de la région de l'Ouest croulent sous les tomates, dont le prix chute de 500 FCFA à 150 FCFA le kilogramme, tandis qu'en saison sèche le kilogramme atteint 1\,200 FCFA. Une coopérative décide de transformer une partie de sa récolte en pâte de tomate. Le technicien mesure l'évolution de la charge microbienne totale de la pulpe de tomate au cours des étapes de fabrication :

\begin{center}
\begin{tabularx}{\linewidth}{|l|X|}\hline
\rowcolor{popblueL} Étape & Charge microbienne (UFC/g) \\ \hline
Tomates entières après récolte & $1{,}0 \times 10^{5}$ \\ \hline
Après tri et lavage à l'eau propre & $2{,}0 \times 10^{4}$ \\ \hline
Après ébouillantage (95 \textcelsius, 3 min) et broyage & $5{,}0 \times 10^{2}$ \\ \hline
Après concentration (cuisson 30 min) & $5{,}0 \times 10^{1}$ \\ \hline
Après stérilisation des bocaux (115 \textcelsius, 20 min) & $< 1$ \\ \hline
\end{tabularx}
\end{center}

\begin{enumerate}
\item À partir des données économiques de l'énoncé, calcule le gain potentiel par kilogramme pour un producteur qui achèterait des tomates en pleine saison pour revendre de la pâte en saison sèche (on supposera qu'1 kg de pâte nécessite 6 kg de tomates fraîches et se vend l'équivalent de 2\,000 FCFA ; les autres coûts sont estimés à 500 FCFA par kg de pâte).
\item Représente sous forme de logigramme les étapes de fabrication de la pâte de tomate décrites dans le tableau.
\item Analyse le tableau : calcule le facteur de réduction de la charge microbienne à chaque étape et identifie les deux étapes les plus efficaces.
\item Pour chacune de ces deux étapes, précise le facteur d'altération visé (micro-organismes, enzymes, oxygène) et le mécanisme d'action (dénaturation des protéines, destruction des spores, etc.).
\item Explique pourquoi la concentration par cuisson contribue aussi à la conservation par un mécanisme autre que la chaleur (on fera intervenir l'activité de l'eau).
\item Dégage deux intérêts économiques et un intérêt nutritionnel de la transformation locale des tomates pour les populations de la région de l'Ouest.
\end{enumerate}
\end{exercice}

\newpage

\coursec{Corrigés des exercices}

\courssub{Je vérifie mes connaissances}

\begin{corrige}[Exercice 1 — QCM : les facteurs d'altération]
\begin{enumerate}
\item \textbf{Bonnes réponses : 1, 2, 4.}
\begin{itemize}
\item \textbf{1.} Une \cle{activité de l'eau} ($a_w$) proche de 1 (eau très disponible) favorise le développement de la quasi-totalité des bactéries, levures et moisissures.
\item \textbf{2.} La plage \SI{20}{\degreeCelsius}--\SI{40}{\degreeCelsius} correspond à la température optimale de croissance des micro-organismes mésophiles (majoritaires sur fruits).
\item \textbf{3.} \emph{Faux.} Un pH très acide (pH $< 3$) \emph{inhibe} la plupart des bactéries d'altération (seules certaines levures et moisissures tolèrent un pH aussi bas). La mangue mûre a un pH $\approx$ 3,5--4,5, favorable aux moisissures/levures.
\item \textbf{4.} Une teneur élevée en eau (mangue $\approx$ \SI{82}{\%}) assure une $a_w$ élevée, condition nécessaire au métabolisme microbien.
\end{itemize}

\item \textbf{Bonnes réponses : 2, 4.}
\begin{itemize}
\item \textbf{2.} Le \cle{brunissement enzymatique} est catalysé par les \cle{polyphénoloxydases} (PPO) qui oxydent les polyphénols en quinones (brunes) en présence de $\ce{O2}$.
\item \textbf{4.} C'est bien une réaction d'oxydation enzymatique (substrats : polyphénols ; co-substrat : $\ce{O2}$ ; enzyme : PPO).
\item \textbf{1.} Les bactéries lactiques provoquent une fermentation lactique (acidification), pas le brunissement.
\item \textbf{3.} La fermentation alcoolique (levures) produit $\ce{CO2}$ et éthanol, pas de pigments bruns.
\end{itemize}

\item \textbf{Bonne réponse : 2.}
La \cle{stérilisation} (ex. \SI{115}{\degreeCelsius}--\SI{121}{\degreeCelsius} en autoclave, ou \SI{100}{\degreeCelsius} pour produits acides pH $< 4,5$) vise la destruction de \emph{tous} les micro-organismes \emph{et de leurs spores} (stérilité commerciale). La pasteurisation ne détruit que les formes végétatives.

\item \textbf{Bonne réponse : 2.}
Le sucre (saccharose) abaisse l'activité de l'eau ($a_w$) par effet osmotique (liaison de l'eau libre). La plupart des bactéries ne se développent plus si $a_w < 0,90$ ; les moisissures/levures s'arrêtent vers $a_w < 0,60\text{--}0,70$. Le sucre ne tue pas directement (pas bactéricide à dose alimentaire), n'augmente pas le pH (l'abaisse légèrement) et ne détruit pas les spores.
\end{enumerate}
\end{corrige}

\begin{corrige}[Exercice 2 — Vrai ou faux justifié]
\begin{enumerate}
\item \textbf{Faux.} Le \cle{séchage} (déshydratation) abaisse l'activité de l'eau ($a_w$) sous le seuil de croissance des micro-organismes (généralement $a_w < 0,60$ pour les moisissures), ce qui \emph{inhibe} leur multiplication. Il ne tue pas systématiquement tous les micro-organismes : les spores bactériennes et fongiques, ainsi que certaines levures xérotolérantes, survivent à la déshydratation et peuvent germer si l'$a_w$ remonte (réhydratation).
\item \textbf{Faux.} La \cle{pasteurisation} (ex. \SI{70}{\degreeCelsius} / \SI{15}{\minute}) détruit les formes végétatives (bactéries, levures, moisissures) et les agents pathogènes, mais \emph{ne détruit pas les spores bactériennes} (thermorésistantes). Seule une stérilisation (température $> \SI{100}{\degreeCelsius}$, typiquement \SI{115}{\degreeCelsius}--\SI{121}{\degreeCelsius}) assure la destruction des spores.
\item \textbf{Vrai.} Le \cle{blanchiment} (ébouillantage : \SI{90}{\degreeCelsius}--\SI{100}{\degreeCelsius}, \SI{1}{\minute}--\SI{3}{\minute}) provoque la dénaturation thermique des enzymes endogènes (polyphénoloxydases, péroxydases, pectinestérases) responsables du brunissement enzymatique, du ramollissement et des pertes de vitamines. C'est une étape cruciale avant séchage, congélation ou mise en conserve.
\item \textbf{Vrai.} Une \cle{confiture} contient $\approx$ \SI{60}{\%}--\SI{70}{\%} de sucre ($a_w \approx 0,75\text{--}0,80$), un pH acide (pH $\approx$ 3,0--3,5) et a subi un traitement thermique (cuisson $\approx$ \SI{105}{\degreeCelsius}) suivi d'un remplissage à chaud hermétique. Ces barrières combinées (facteurs \emph{hurdles}) assurent une stabilité microbiologique de plusieurs mois à température ambiante tant que le pot n'est pas ouvert.
\item \textbf{Vrai.} La \cle{vitamine C} (acide ascorbique) s'oxyde facilement en acide déhydroascorbique (réversible) puis en produits inactifs (irréversible). Cette oxydation est catalysée par la lumière (photodégradation), la chaleur (accélération cinétique), la présence d'$\ce{O2}$, d'ions métalliques ($\ce{Fe^{2+}}, \ce{Cu^{2+}}$) et favorisée par un pH élevé.
\end{enumerate}
\end{corrige}

\begin{corrige}[Exercice 3 — Définitions et applications]
\begin{enumerate}
\item \textbf{Pasteurisation :} Traitement thermique modéré (généralement \SI{60}{\degreeCelsius}--\SI{90}{\degreeCelsius} pendant quelques secondes à plusieurs minutes) destiné à détruire les micro-organismes pathogènes et les formes végétatives d'altération, tout en préservant au maximum les qualités nutritionnelles et organoleptiques du produit.
\emph{Exemple :} Jus de mangue chauffé à \SI{70}{\degreeCelsius} pendant \SI{15}{\minute} puis mis en bouteille à chaud (conservation quelques semaines/mois au frais).

\item \textbf{Stérilisation :} Traitement thermique sévère (température $\ge \SI{100}{\degreeCelsius}$, généralement \SI{115}{\degreeCelsius}--\SI{121}{\degreeCelsius} sous pression en autoclave, ou \SI{100}{\degreeCelsius} pour produits acides pH $< 4,5$) appliqué à un produit conditionné hermétiquement, visant à détruire \emph{tous} les micro-organismes (formes végétatives \emph{et} spores) pour obtenir la \cle{stérilité commerciale} (conservation années à température ambiante).
\emph{Exemple :} Conserves de tomates entières ou de mangues au sirop en bocaux stérilisés à \SI{100}{\degreeCelsius} (pH acide) ou \SI{115}{\degreeCelsius} (pH neutre) pendant \SI{20}{\minute}--\SI{45}{\minute}.

\item \textbf{Blanchiment (ébouillantage) :} Traitement thermique court (\SI{90}{\degreeCelsius}--\SI{100}{\degreeCelsius}, \SI{1}{\minute}--\SI{5}{\minute}) suivi d'un refroidissement rapide, appliqué aux fruits/légumes frais avant transformation, pour inactiver les enzymes endogènes (PPO, POD, pectinestérases), réduire la charge microbienne de surface et chasser l'oxygène intercellulaire.
\emph{Exemple :} Tomates plongées \SI{3}{\minute} à \SI{95}{\degreeCelsius} avant broyage/tamisage pour inactiver les pectinestérases (éviter la séparation de phases) et les PPO (éviter brunissement) ; mangues blanchies \SI{2}{\minute} avant séchage.

\item \textbf{Activité de l'eau ($a_w$) :} Rapport entre la pression de vapeur d'eau exercée par le produit alimentaire ($p$) et la pression de vapeur d'eau saturante de l'eau pure ($p_0$) à la même température : $a_w = p / p_0$. Elle quantifie l'eau \emph{libre} (disponible pour le métabolisme microbien et les réactions chimiques), par opposition à l'eau liée. $0 \le a_w \le 1$.
\emph{Exemple :} Mangue fraîche : $a_w \approx 0,98$ (très favorable aux microbes) ; Mangue séchée (\SI{15}{\%} humidité) : $a_w \approx 0,55$ (inhibiteur) ; Confiture : $a_w \approx 0,75$.

\item \textbf{Déshydratation osmotique :} Opération de transfert de masse par osmose : immersion de tranches de fruit dans une solution hypertonique (sirop de sucre concentré ou saumure) qui extrait l'eau du fruit (perte d'eau) et fait pénétrer du soluté (gain de solides), réalisant une déshydratation partielle à basse température avant un séchage final.
\emph{Exemple :} Tranches de mangue trempées \SI{4}{\hour}--\SI{6}{\hour} dans un sirop à \SI{60}{\degreeBrix}--\SI{70}{\degreeBrix} (sucre) à \SI{40}{\degreeCelsius}--\SI{50}{\degreeCelsius}, puis égouttées et séchées au séchoir solaire. Réduit le temps de séchage, améliore la couleur et la texture (moins de rétractation).
\end{enumerate}
\end{corrige}

\begin{corrige}[Exercice 4 — Calculs directs : rendements et humidité]
Données : $m_{\text{fraîches}} = \SI{250}{\kg}$ ; $\% \text{eau} = \SI{82}{\%}$ ; $\% \text{MS} = \SI{18}{\%}$.
\begin{enumerate}
\item Masse d'eau : $m_{\text{eau}} = 250 \times 0,82 = \textbf{\SI{205}{\kg}}$.
Masse de matière sèche (MS) : $m_{\text{MS}} = 250 - 205 = \textbf{\SI{45}{\kg}}$ (ou $250 \times 0,18 = \SI{45}{\kg}$).

\item Après séchage : $\% \text{eau} = \SI{15}{\%}$ $\Rightarrow$ $\% \text{MS} = \SI{85}{\%}$.
La MS est conservée : $m_{\text{MS}} = \SI{45}{\kg} = 0,85 \times m_{\text{séchées}}$.
$m_{\text{séchées}} = \dfrac{45}{0,85} \approx \textbf{\SI{52,94}{\kg}}$.
(Vérification : eau résiduelle $= 52,94 \times 0,15 \approx \SI{7,94}{\kg}$ ; total $= 45 + 7,94 = \SI{52,94}{\kg}$.)

\item Bilan de masse confiture :
Masse initiale = pulpe + sucre = $\SI{4}{\kg} + \SI{3}{\kg} = \SI{7}{\kg}$.
Masse finale (confiture) = $\SI{5,6}{\kg}$.
Masse d'eau évaporée = Masse initiale $-$ Masse finale = $7 - 5,6 = \textbf{\SI{1,4}{\kg}}$.
\end{enumerate}
\end{corrige}

\courssub{Je m'entraîne}

\begin{corrige}[Exercice 5 — Logigramme de fabrication de la pâte de tomate]
\begin{enumerate}
\item \textbf{Logigramme :}
\begin{popfigure}
\begin{tikzpicture}[node distance=1.2cm and 0.5cm, >=Stealth,
  etape/.style={rectangle, draw=popblue, thick, fill=popblueL, rounded corners, minimum width=3cm, minimum height=1.1cm, align=center, font=\small},
  fleche/.style={->, thick, popdark}]
\node[etape, align=center] (1){Tri et lavage\\(eau propre)};
\node[etape, below of=1, align=center] (2){Ébouillantage\\\SI{95}{\degreeCelsius}, \SI{3}{\minute}};
\node[etape, below of=2] (3) {Égouttage et broyage};
\node[etape, below of=3, align=center] (4){Tamisage\\(élim. peaux/pépins)};
\node[etape, below of=4, align=center] (5){Concentration\\(cuisson $\to$ \SI{28}{\%} MS)};
\node[etape, below of=5, align=center] (6){Remplissage à chaud\\bocaux};
\node[etape, below of=6, align=center] (7){Stérilisation bocaux\\\SI{100}{\degreeCelsius}, \SI{45}{\minute}};
\node[etape, below of=7, align=center] (8){Refroidissement\\et stockage};
\draw[fleche] (1) -- (2);
\draw[fleche] (2) -- (3);
\draw[fleche] (3) -- (4);
\draw[fleche] (4) -- (5);
\draw[fleche] (5) -- (6);
\draw[fleche] (6) -- (7);
\draw[fleche] (7) -- (8);
\end{tikzpicture}
\legende{Logigramme de fabrication de la pâte de tomate (conserve stérilisée).}
\end{popfigure}

\item \textbf{Facteurs d'altération visés par étape :}
\begin{center}
\begin{tabularx}{\linewidth}{|l|X|X|}\hline
\rowcolor{popblueL} \textbf{Étape} & \textbf{Facteur(s) visé(s)} & \textbf{Justification} \\ \hline
Tri et lavage & Micro-organismes (charge initiale), Corps étrangers & Élimination physique des souillures, terres, parties pourries, réduction flore de surface. \\ \hline
Ébouillantage (\SI{95}{\degreeCelsius}, \SI{3}{\minute}) & \textbf{Enzymes} (PPO, POD, pectinestérases), Micro-organismes (végétatifs) & Dénaturation thermique des enzymes (inactivation brunissement, ramollissement). Destruction bactéries/levures/moisissures de surface. Chasse l'$\ce{O2}$ intercellulaire. \\ \hline
Égouttage / Broyage & --- & Préparation physique (réduction taille, libération pulpe). \\ \hline
Tamisage & Corps étrangers (peaux, pépins) & Élimination parties indésirables (fibres, graines) source d'enzymes résiduelles et défauts texture. \\ \hline
Concentration (cuisson) & \textbf{Eau ($a_w$)}, Micro-organismes, Enzymes résiduelles & Évaporation eau $\to$ baisse $a_w$ (effet hurdle). Chaleur prolongée $\to$ destruction microbienne complémentaire, inactivation enzymes résiduelles. \\ \hline
Remplissage à chaud & \textbf{Micro-organismes (contenant)}, \textbf{Oxygène (tête)} & Stérilisation interne bocal/couvercle par le produit chaud. Chasse l'air de la tête $\to$ vide partiel au refroidissement (limite oxydation, moisissures aérobies). \\ \hline
Stérilisation (\SI{100}{\degreeCelsius}, \SI{45}{\minute}) & \textbf{Micro-organismes + Spores} & Destruction totale flore résiduelle (végétatives + spores) dans le produit \emph{et} le contenant hermétique $\to$ stérilité commerciale. \\ \hline
Refroidissement / Stockage & --- & Arrêt cuisson, condensation vapeur $\to$ vide, contrôle intégrité bocaux. \\ \hline
\end{tabularx}
\end{center}

\item \textbf{Pourquoi remplissage « à chaud » ?}
Trois raisons majeures :
\begin{enumerate}
\item \textbf{Stérilisation du contenant :} Le produit chaud ($\ge \SI{85}{\degreeCelsius}$) stérilise l'intérieur du bocal et du couvercle au contact.
\item \textbf{Création du vide partiel :} Au refroidissement, la vapeur d'eau de la tête se condense, créant une dépression (vide) qui assure l'herméticité (couvercle « bombé » rentrant), empêche le développement de moisissures aérobies et limite l'oxydation (vitamine C, couleur).
\item \textbf{Limitation du choc thermique :} Évite l'éclatement du verre par différence de température brutale (produit froid dans bocal chaud ou inverse).
\end{enumerate}
\end{enumerate}
\end{corrige}

\begin{corrige}[Exercice 6 — Comparaison multicritères de trois procédés]
\begin{enumerate}
\item \textbf{Durée de conservation :} Conserve au sirop (\SI{2}{\an}--\SI{3}{\an}) $>$ Confiture (\SI{12}{\month} pot fermé) $>$ Mangue séchée (\SI{6}{\month}--\SI{12}{\month}). La stérilisation en conserve hermétique offre la plus longue stabilité.
\item \textbf{Qualité nutritionnelle (Vitamine C conservée) :} Conserve au sirop (\SI{45}{\%}) $>$ Mangue séchée (\SI{30}{\%}) $>$ Confiture (\SI{20}{\%}). Le traitement thermique le plus doux (sirop, stérilisation \SI{100}{\degreeCelsius} court) préserve mieux la vitamine C thermosensible que la cuisson prolongée de la confiture (concentration + chaleur) ou le séchage long (oxydation résiduelle + chaleur).

\item \textbf{Procédure recommandée : Le séchage (séchoir solaire).}
\textbf{Arguments (3 critères minimum) :}
\begin{enumerate}
\item \textbf{Coût de l'équipement :} \emph{Faible} (séchoir solaire artisanal réalisable localement avec bois, grillage, bâche transparente), accessible à un groupement villageois sans crédit.
\item \textbf{Consommation d'énergie :} \emph{Très faible} (énergie solaire gratuite), indispensable en zone non électrifiée (pas de facture, pas de groupe électrogène).
\item \textbf{Risque de recontamination après ouverture :} \emph{Faible} (produit sec, $a_w \approx 0,55$, stable à l'air ambiant plusieurs semaines si emballage refermé), contrairement à la confiture (risque élevé moisissures surface) ou conserve (risque moyen).
\item \textbf{Durée de conservation :} \SI{6}{\month}--\SI{12}{\month}, suffisante pour couvrir la période de soudure jusqu'à la prochaine saison.
\end{enumerate}

\item \textbf{Conclusion :} Il n'existe pas de technique de conservation « idéale » universelle. Chaque procédé résulte d'un compromis entre durée de conservation, qualité nutritionnelle, coût d'investissement, consommation énergétique et praticité d'usage. Le séchage solaire s'impose en milieu rural non électrifié par sa sobriété énergétique et son faible coût ; la confiture convient à l'artisanat urbain à faible capacité d'investissement ; la conserve stérilisée, plus coûteuse en énergie et équipement, garantit la plus longue conservation et la meilleure rétention en vitamine C pour une filière commerciale structurée. Le choix technique doit s'adapter au contexte local (ressources, marché, consommation).
\end{enumerate}
\end{corrige}

\begin{corrige}[Exercice 7 — La confiture qui moisit après ouverture]
\begin{enumerate}
\item \textbf{Deux facteurs de stabilité en pot fermé :}
\begin{enumerate}
\item \textbf{Faible activité de l'eau ($a_w \approx 0,75\text{--}0,80$) :} Due à la forte concentration en sucre (effet osmotique), elle inhibe le développement de la plupart des bactéries ($a_{w,\min} \approx 0,90$) et limite les moisissures/levures ($a_{w,\min} \approx 0,60\text{--}0,70$).
\item \textbf{Absence d'oxygène / Stérilité commerciale :} Le traitement thermique (cuisson + remplissage à chaud) a détruit les micro-organismes présents ; le conditionnement hermétique empêche toute recontamination et prive les moisissures aérobies strictes de l'$\ce{O2}$ nécessaire à leur croissance.
\end{enumerate}

\item \textbf{Explication du développement après ouverture :}
L'ouverture du pot brise les deux barrières :
\begin{itemize}
\item \textbf{Apport d'oxygène :} L'air ambiant (\SI{21}{\%} $\ce{O2}$) entre en contact avec la surface de la confiture. Les moisissures sont des champignons \emph{aérobies strictes} : elles ont \emph{absolument besoin} d'$\ce{O2}$ pour respirer et croître. Les spores (présentes dans l'air, sur le couvercle, les ustensiles) peuvent alors germer.
\item \textbf{Réhydratation de surface / $a_w$ locale :} À l'air libre, la surface de la confiture peut absorber l'humidité ambiante (surtout si hygrométrie élevée) ou former de la condensation, augmentant localement l'$a_w$ au-dessus du seuil critique des moisissures ($\approx 0,70$). De plus, la consommation de sucre par les moisissures (métabolisme) produit de l'eau métabolique, entretenant l'humidité.
\item \textbf{Contamination secondaire :} L'utilisation d'une cuillère non propre, le contact avec l'air chargé en spores (\emph{Aspergillus, Penicillium, Rhizopus}) ensemencent le milieu.
\end{itemize}

\item \textbf{Deux mesures pratiques pour ralentir l'altération :}
\begin{enumerate}
\item \textbf{Conserver au réfrigérateur (\SI{4}{\degreeCelsius}) après ouverture :} La basse température ralentit drastiquement la vitesse de germination des spores et la croissance mycélienne (facteur cinétique), prolongeant la durée de vie de plusieurs semaines.
\item \textbf{Refermer hermétiquement le pot immédiatement après chaque utilisation} (couler une fine couche d'alcool alimentaire ou de sucre en surface avant de fermer est un plus) \emph{et utiliser une cuillère propre et sèche} à chaque prélèvement pour éviter d'inoculer des spores ou d'apporter de l'eau.
\end{enumerate}
\end{enumerate}
\end{corrige}

\begin{corrige}[Exercice 8 — Exploitation d'un tableau de survie microbienne]
\begin{enumerate}
\item \textbf{Facteur de réduction ($FR$) = $\dfrac{\text{UFC Témoin}}{\text{UFC Traitement}}$ (Témoin $= \num{2,0e6}$ UFC/g).}
\begin{center}
\begin{tabularx}{\linewidth}{|l|X|X|}\hline
\rowcolor{popblueL} \textbf{Traitement} & \textbf{Levures (UFC/g)} & \textbf{Facteur de réduction} \\ \hline
Témoin & $\num{2,0e6}$ & 1 (référence) \\ \hline
Blanchiment \SI{2}{\minute} à \SI{90}{\degreeCelsius} & $\num{5,0e4}$ & $\dfrac{\num{2,0e6}}{\num{5,0e4}} = \textbf{40}$ \\ \hline
Blanchiment + Séchage ($a_w = 0,55$) & $\num{1,0e2}$ & $\dfrac{\num{2,0e6}}{\num{1,0e2}} = \textbf{20\,000}$ \\ \hline
Blanchiment + Séchage + Emballage étanche & $< 10$ & $> \dfrac{\num{2,0e6}}{10} = \textbf{200\,000}$ \\ \hline
\end{tabularx}
\end{center}

\item \textbf{Blanchiment seul insuffisant :}
Le blanchiment réduit la charge de 40 fois seulement, laissant $\num{5,0e4}$ UFC/g. Ce niveau reste \emph{très élevé} : à \SI{25}{\degreeCelsius} et $a_w$ élevée (mangue fraîche $\approx 0,98$), ces levures survivantes (thermorésistantes ou protégées) vont se multiplier rapidement (temps de génération $\approx$ \SI{90}{\minute}) et atteindre $> \num{1e7}$ UFC/g en quelques heures, causant une fermentation alcoolique (gaz, odeur) et une altération sensorielle en moins de \SI{24}{\hour}--\SI{48}{\hour}. Le blanchiment ne détruit pas les spores et n'abaisse pas l'$a_w$.

\item \textbf{Rôle complémentaire :}
\begin{itemize}
\item \textbf{Séchage ($a_w = 0,55$) :} Abaisse l'activité de l'eau bien en dessous du minimum de croissance des levures ($a_{w,\min} \approx 0,88$ pour la plupart, $\approx 0,60$ pour les xérotolérantes). L'eau devient indisponible (liée aux solutés), le métabolisme s'arrête (plasmolysse, dénaturation enzymatique). Le facteur de réduction passe de 40 à 20\,000 (x500 supplémentaire).
\item \textbf{Emballage étanche :} Maintient l'$a_w$ basse en empêchant la réabsorption de l'humidité ambiante (hygiène + barrière vapeur). Empêche la recontamination par les spores aéroportées, les insectes, les poussières. Assure la stabilité dans le temps (facteur $> 200\,000$).
\end{itemize}
La combinaison \emph{Blanchiment (enzymes + flore végétative) + Séchage ($a_w$) + Emballage (maintien $a_w$ + barrière)} illustre la technologie des \emph{barrières combinées (hurdle technology)}.
\end{enumerate}
\end{corrige}

\courssub{Je me prépare au Bac}

\begin{corrige}[Exercice 9 — Type Bac : activité de l'eau et séchage de la mangue]
\begin{enumerate}
\item \textbf{Représentation graphique (courbe $D = f(a_w)$) :}
\begin{popfigure}
\begin{tikzpicture}
\begin{axis}[
width=\linewidth, height=8cm,
xlabel={Activité de l'eau $a_w$}, ylabel={Diamètre colonie (mm)},
xmin=0.55, xmax=1.0, ymin=0, ymax=45,
xtick={0.6,0.65,0.7,0.75,0.8,0.85,0.9,0.95,1.0},
ytick={0,5,10,15,20,25,30,35,40,45},
grid=major, grid style={popline},
tick label style={font=\small},
label style={font=\small},
axis lines=middle,
enlargelimits={abs=0.02},
]
\addplot[popcurve, very thick, mark=*, mark size=2pt, popblue] coordinates {
(0.98,42) (0.95,35) (0.90,22) (0.80,9) (0.70,2) (0.60,0)
};
\addplot[poporange, dashed, thick] coordinates {(0.60,0) (0.60,-5)};
\node[poporange, anchor=north, font=\small] at (axis cs:0.60,-2) {$a_{w,\text{crit}} \approx 0,60$};
\end{axis}
\end{tikzpicture}
\legende{Croissance d'\emph{Aspergillus niger} en fonction de l'activité de l'eau ($a_w$) sur milieu pulpe de mangue.}
\end{popfigure}
\textit{Échelles respectées : 1 cm pour 0,05 $a_w$ (abscisse), 1 cm pour 5 mm (ordonnée). Points tracés et reliés.}

\item \textbf{Analyse :}
La croissance (diamètre) \emph{diminue de manière non linéaire} lorsque l'$a_w$ diminue. Elle est maximale à $a_w = 0,98$ (\SI{42}{\mm}), chute fortement entre 0,90 et 0,80 (\SI{22}{\mm} $\to$ \SI{9}{\mm}), devient très faible à 0,70 (\SI{2}{\mm}) et \emph{nulle} à $a_w = 0,60$ (\SI{0}{\mm}).
\textbf{Seuil critique graphique :} $a_w \approx \mathbf{0,60}$ (valeur en dessous de laquelle le diamètre = \SI{0}{\mm}). Cela correspond au minimum d'$a_w$ pour \emph{Aspergillus niger} (moisissure xérotolérante).

\item \textbf{Explication cellulaire (niveau moléculaire) :}
Une faible $a_w$ signifie un \emph{potentiel hydrique} ($\Psi_w$) très bas (très négatif) dans le milieu extérieur. Par osmose, l'eau sort du cytoplasme microbien vers le milieu (gradient de potentiel hydrique). Cela entraîne :
\begin{itemize}
\item \textbf{Plasmolysse :} Rétraction de la membrane plasmique par rapport à la paroi (champignons/bactéries), rupture des contacts membrane-paroi, arrêt des transports membranaires.
\item \textbf{Déshydratation cytoplasmique :} Augmentation de la concentration ionique interne $\to$ dénaturation des enzymes, précipitation des protéines, arrêt du métabolisme (glycolyse, synthèse protéique, réplication ADN).
\item \textbf{Perte de turgescence :} Arrêt de l'expansion cellulaire (croissance) et de la division.
\end{itemize}
Les micro-organismes entrent en état de \emph{viabilité non culturable} (dormance) ou meurent. Seules les spores (paroi épaisse, cytoplasme déshydraté, protéines de choc thermique, tréhalose) survivent.

\item \textbf{Preuve de conservation par séchage à \SI{15}{\%} humidité :}
L'énoncé indique : $a_w \le 0,60 \Leftrightarrow$ humidité résiduelle $\approx \SI{15}{\%}$ (masse). Or, d'après la courbe, pour $a_w \le 0,60$, la croissance d'\emph{Aspergillus niger} est \emph{nulle}. Donc une mangue séchée à \SI{15}{\%} d'humidité ($a_w \le 0,60$) ne permet pas le développement de cette moisissure (ni de la plupart des bactéries/levures), assurant sa conservation microbiologique.

\textbf{Calcul de la masse d'eau à éliminer (base \SI{100}{\kg} mangues fraîches) :}
\begin{itemize}
\item Mangues fraîches : \SI{100}{\kg}. Eau = \SI{82}{\kg}. Matière sèche (MS) = \SI{18}{\kg}.
\item Mangues séchées (cible \SI{15}{\%} eau) : MS = \SI{85}{\%} de la masse finale $m_f$.
$\SI{18}{\kg} = 0,85 \times m_f \Rightarrow m_f = \dfrac{18}{0,85} \approx \SI{21,18}{\kg}$.
\item Eau résiduelle dans mangues séchées = $m_f \times 0,15 \approx \SI{3,18}{\kg}$ (ou $21,18 - 18 = \SI{3,18}{\kg}$).
\item \textbf{Eau à éliminer} = Eau initiale $-$ Eau résiduelle = $82 - 3,18 = \textbf{\SI{78,82}{\kg}}$.
\end{itemize}
(Soit une perte de masse totale de $100 - 21,18 = \SI{78,82}{\kg}$, essentiellement de l'eau).

\item \textbf{Amélioration technique : Séchoir solaire ventilé de type « tunnel » ou « caisson » à convection forcée (cheminée solaire ou petit ventilateur PV).}
\textbf{Justification :}
\begin{itemize}
\item \textbf{Hygiène :} Structure close (filet anti-insectes, bâche transparente) $\to$ protection contre poussières, insectes, fientes d'oiseaux, animaux (qualité sanitaire).
\item \textbf{Vitesse :} Effet de serre (T° interne \SI{45}{\degreeCelsius}--\SI{60}{\degreeCelsius}) + ventilation (renouvellement air chaud humide) $\to$ gradient de pression de vapeur maximal $\to$ séchage 2 à 3 fois plus rapide qu'au soleil nu (réduit le risque de développement microbien \emph{pendant} le séchage).
\item \textbf{Qualité :} Moins d'oxydation (temps réduit, protection UV partielle) $\to$ meilleure couleur (caroténoïdes) et rétention vitamine C.
\item \textbf{Contexte rural camerounais :} Matériaux locaux (bambou, bois, bâche agricole, grillage), faible coût, pas d'électricité requise (convection naturelle par cheminée ou petit panneau PV \SI{10}{\watt} pour ventilateur), reproductible par groupements.
\end{itemize}
\end{enumerate}

\vspace{0,5cm}
\noindent\textbf{Barème indicatif (sur 20) :}
\begin{itemize}
\item Q1 : Graphique correct (échelles, points, courbe) : 3 pts.
\item Q2 : Analyse variation + seuil $a_w \approx 0,60$ : 3 pts.
\item Q3 : Explication cellulaire (osmose, plasmolysse, métabolisme) : 4 pts.
\item Q4 : Preuve ($a_w \le 0,60 \leftrightarrow$ \SI{15}{\%} H) + Calcul eau éliminée (\SI{78,8}{\kg}) : 5 pts.
\item Q5 : Proposition technique adaptée + 3 justifications (hygiène, rapidité, qualité, contexte) : 5 pts.
\end{itemize}
\textbf{Points de vigilance :} Unités obligatoires (kg, mm, \%). Courbe : axes légendés, unités, points visibles. Seuil $a_w$ lu sur graphique (pas calculé). Explication cellulaire : vocabulaire précis (potentiel hydrique, osmose, plasmolysse, turgescence). Calcul : raisonnement sur MS invariante.
\end{corrige}

\begin{corrige}[Exercice 10 — Type Bac : traitements thermiques du jus de mangue]
\begin{enumerate}
\item \textbf{Problème scientifique :} Quel est l'effet de l'intensité du traitement thermique (température $\times$ durée) sur la stabilité microbiologique et la qualité nutritionnelle (vitamine C) d'un jus de mangue conditionné en bouteille fermée ?
\textbf{Hypothèse testée :} L'augmentation de la température de traitement (de 0 à \SI{70}{\degreeCelsius} puis \SI{100}{\degreeCelsius}) améliore la destruction microbienne (stabilité) mais dégrade davantage la vitamine C (qualité nutritionnelle).

\item \textbf{Interprétation Lot A (témoin non chauffé) :}
\begin{itemize}
\item \textbf{Micro-organismes :} Développement massif de \textbf{levures} (fermentation alcoolique) et de \textbf{bactéries acétiques/lactiques} (odeur acide). Dénombrement $\num{8e7}$ UFC/mL = altération majeure.
\item \textbf{Type fermentation :} \textbf{Fermentation alcoolique} (levures : \emph{Saccharomyces}, \emph{Candida}) $\to$ production de $\ce{CO2}$ (gaz, bouteille gonflée) et d'éthanol (odeur alcoolique). Puis fermentation acétique (bactéries acétiques : \emph{Acetobacter}) sur l'éthanol $\to$ acide acétique (odeur vinaigrée/acide).
\item \textbf{Équation bilan fermentation alcoolique du glucose :}
\[ \ce{C6H12O6 -> 2 C2H5OH + 2 CO2 + 2 ATP} \]
\end{itemize}

\item \textbf{Procédés et distinctions :}
\begin{center}
\begin{tabularx}{\linewidth}{|l|X|X|X|}\hline
\rowcolor{popblueL} \textbf{Lot} & \textbf{Procédé} & \textbf{Conditions} & \textbf{Micro-organismes détruits} \\ \hline
B & \textbf{Pasteurisation} & \SI{70}{\degreeCelsius}, \SI{15}{\minute}, remplissage à chaud & Formes végétatives : bactéries, levures, moisissures. \textbf{PAS les spores}. \\ \hline
C & \textbf{Stérilisation} (ou « pasteurisation stérilisante » pour pH $< 4,5$) & \SI{100}{\degreeCelsius}, \SI{20}{\minute}, \emph{in flacon fermé} & Formes végétatives \textbf{ET spores} (grâce à l'acidité pH $< 4,5$ + \SI{100}{\degreeCelsius}). \\ \hline
\end{tabularx}
\end{center}
\textit{Distinction clé :} La stérilisation (Lot C) détruit les spores (grâce à la combinaison pH acide + \SI{100}{\degreeCelsius}), la pasteurisation (Lot B) ne les détruit pas.

\item \textbf{Stabilité Lots B et C vs Altération Lot A :}
\begin{itemize}
\item \textbf{Lot A :} Absence de traitement $\to$ flore initiale (levures, bactéries) se développe librement (pH acide favorable, $a_w$ élevée, nutriments, $\ce{O2}$ résiduel) $\to$ altération rapide.
\item \textbf{Lot B (Pasteurisation) :} Destruction des formes végétatives $\to$ charge résiduelle très faible ($\num{3e2}$ UFC/mL), inférieure au seuil de détection organoleptique. \textbf{Pourquoi quelques UFC ?} Survivance de \emph{spores} thermorésistantes (bactéries \emph{Bacillus}, \emph{Clostridium} ; moisissures \emph{Byssochlamys}) non détruites à \SI{70}{\degreeCelsius}. Dans un jus acide (pH $\approx$ 3,5--4,0), anaérobie (bouteille pleine), ces spores ne germent généralement pas (inhibées par pH + absence $\ce{O2}$ + compétition), d'où stabilité pour plusieurs semaines/mois.
\item \textbf{Lot C (Stérilisation) :} Destruction formes végétatives + spores $\to$ stérilité commerciale ($< 10$ UFC/mL = limite détection). Stabilité longue (années).
\end{itemize}

\item \textbf{Calcul Vitamine C et Compromis :}
Jus frais : \SI{30}{\mg}/\SI{100}{\mL}.
\begin{itemize}
\item Lot B (Pasteurisation) : Perte \SI{15}{\%} $\Rightarrow$ Teneur = $30 \times (1 - 0,15) = 30 \times 0,85 = \textbf{\SI{25,5}{\mg}/\SI{100}{\mL}}$.
\item Lot C (Stérilisation) : Perte \SI{40}{\%} $\Rightarrow$ Teneur = $30 \times (1 - 0,40) = 30 \times 0,60 = \textbf{\SI{18,0}{\mg}/\SI{100}{\mL}}$.
\end{itemize}
\textbf{Discussion du compromis :}
La stérilisation (Lot C) offre une \textbf{sécurité microbiologique maximale} (destruction spores, conservation années à T° ambiante) mais au prix d'une \textbf{perte nutritionnelle majeure} (\SI{40}{\%} vitamine C, goût « cuit » par réaction de Maillard/caramélisation). La pasteurisation (Lot B) offre une \textbf{sécurité suffisante pour une distribution courte} (quelques semaines à mois si pH $< 4,5$, hygiène parfaite, chaîne du froid non requise mais stockage frais préférable) avec une \textbf{meilleure qualité nutritionnelle} (\SI{25,5}{\mg} vs \SI{18}{\mg}, soit +\SI{42}{\%} de vitamine C conservée) et organoleptique (goût frais).

\item \textbf{Conclusion argumentée (Conseil pour petite unité artisanale, vente dans la semaine) :}
\textbf{Je conseille la pasteurisation (Lot B, \SI{70}{\degreeCelsius} / \SI{15}{\minute}).}
\textbf{Justification :}
\begin{enumerate}
\item \textbf{Adéquation durée de vie / circuit court :} La vente se fait dans la semaine. La pasteurisation assure une stabilité de plusieurs semaines à température ambiante (pH mangue $< 4,5$ inhibe germination spores résiduelles), largement suffisante.
\item \textbf{Qualité supérieure :} Rétention vitamine C significativement plus élevée (\SI{25,5}{\mg} vs \SI{18}{\mg}/\SI{100}{\mL}), goût frais préservé (pas de note « cuit »), couleur vive.
\item \textbf{Faisabilité technique/économique :} Équipement simple (marmite, thermomètre, source de chaleur bois/gaz), pas d'autoclave coûteux, consommation d'énergie et temps de traitement moindres (\SI{70}{\degreeCelsius} vs \SI{100}{\degreeCelsius} \SI{20}{\minute}), investissement accessible.
\item \textbf{Rentabilité :} Produit de meilleure qualité se vend mieux/prix supérieur ; coûts de production inférieurs.
\end{enumerate}
La stérilisation (Lot C) est surdimensionnée (coût, qualité) pour un circuit de distribution hebdomadaire.
\end{enumerate}

\vspace{0,5cm}
\noindent\textbf{Barème indicatif (sur 20) :}
\begin{itemize}
\item Q1 : Problème + Hypothèse clairs : 2 pts.
\item Q2 : Interprétation Lot A (micro-organismes, fermentation, équation) : 4 pts.
\item Q3 : Nommage procédés + Distinction (T°, durée, spores) : 3 pts.
\item Q4 : Explication stabilité B/C vs A + Pourquoi UFC résiduels B : 4 pts.
\item Q5 : Calculs Vit C justes + Discussion compromis (sécurité vs qualité) : 4 pts.
\item Q6 : Choix argumenté (pasteurisation) avec 3 arguments contextuels pertinents : 3 pts.
\end{itemize}
\textbf{Points de vigilance :} Équation chimique \ce{} équilibrée. Distinction pasteurisation/stérilisation sur les spores. Justification du choix Lot B liée au contexte (vente semaine, artisanat). Unités vitamine C (mg/100 mL).
\end{corrige}

\begin{corrige}[Exercice 11 — Type Bac : transformation de la tomate et filière de conservation]
\begin{enumerate}
\item \textbf{Gain potentiel par kg de pâte de tomate :}
\begin{itemize}
\item Coût matière première : \SI{6}{\kg} tomates fraiches $\times$ \SI{150}{\FCFA}/\si{\kg} = \textbf{\SI{900}{\FCFA}}.
\item Coûts de transformation (énergie, main d'œuvre, emballage, charges) : \textbf{\SI{500}{\FCFA}} / kg pâte.
\item Coût de revient total : 900 + 500 = \textbf{\SI{1400}{\FCFA}} / kg pâte.
\item Prix de vente pâte : \textbf{\SI{2000}{\FCFA}} / kg pâte.
\item \textbf{Gain (marge brute) = 2000 - 1400 = \textbf{\SI{600}{\FCFA} par kg de pâte produite}.}
\end{itemize}
(Soit un gain de \SI{100}{\FCFA} par kg de tomates fraîches transformées).

\item \textbf{Logigramme de fabrication (d'après tableau charges microbiennes) :}
\begin{popfigure}
\begin{tikzpicture}[node distance=1.1cm and 0.3cm, >=Stealth,
  etape/.style={rectangle, draw=popgreen, thick, fill=popgreenL, rounded corners, minimum width=3.2cm, minimum height=1cm, align=center, font=\small},
  fleche/.style={->, thick, popdark}]
\node[etape, align=center] (1){Tomates entières\\après récolte};
\node[etape, below of=1, align=center] (2){Tri et lavage\\eau propre};
\node[etape, below of=2, align=center] (3){Ébouillantage \SI{95}{\degreeCelsius}, \SI{3}{\minute}\\+ Broyage};
\node[etape, below of=3, align=center] (4){Concentration\\Cuisson \SI{30}{\minute} (\SI{28}{\%} MS)};
\node[etape, below of=4, align=center] (5){Remplissage bocaux\\+ Stérilisation \SI{115}{\degreeCelsius}, \SI{20}{\minute}};
\node[etape, below of=5, align=center] (6){Refroidissement\\Stockage};
\draw[fleche] (1) -- (2);
\draw[fleche] (2) -- (3);
\draw[fleche] (3) -- (4);
\draw[fleche] (4) -- (5);
\draw[fleche] (5) -- (6);
\end{tikzpicture}
\legende{Logigramme de fabrication de la pâte de tomate stérilisée (région de l'Ouest).}
\end{popfigure}

\item \textbf{Analyse facteurs de réduction (FR) :}
\begin{center}
\begin{tabularx}{\linewidth}{|l|X|X|}\hline
\rowcolor{popblueL} \textbf{Étape} & \textbf{Charge (UFC/g)} & \textbf{Facteur de réduction (FR)} \\ \hline
Tomates entières (témoin) & $\num{1,0e5}$ & 1 \\ \hline
Tri + Lavage & $\num{2,0e4}$ & $\num{1,0e5} / \num{2,0e4} = \mathbf{5}$ \\ \hline
Ébouillantage + Broyage & $\num{5,0e2}$ & $\num{2,0e4} / \num{5,0e2} = \mathbf{40}$ \\ \hline
Concentration (cuisson \SI{30}{\minute}) & $\num{5,0e1}$ & $\num{5,0e2} / \num{5,0e1} = \mathbf{10}$ \\ \hline
Stérilisation bocaux (\SI{115}{\degreeCelsius}, \SI{20}{\minute}) & $< 1$ & $\num{5,0e1} / 1 \approx \mathbf{50}$ (minimum) \\ \hline
\end{tabularx}
\end{center}
\textbf{Deux étapes les plus efficaces (FR max) :}
1. \textbf{Stérilisation des bocaux (FR $\ge$ 50).}
2. \textbf{Ébouillantage + Broyage (FR = 40).}

\item \textbf{Facteur d'altération visé et mécanisme pour ces deux étapes :}
\begin{enumerate}
\item \textbf{Ébouillantage (\SI{95}{\degreeCelsius}, \SI{3}{\minute}) :}
\begin{itemize}
\item \textbf{Facteur visé :} \textbf{Enzymes} (PPO, POD, pectinestérases) \textbf{et Micro-organismes végétatifs} (flore de surface).
\item \textbf{Mécanisme :} \textbf{Dénaturation thermique des protéines} (enzymes et protéines structurales microbiennes). La chaleur humide à \SI{95}{\degreeCelsius} casse les liaisons hydrogène et hydrophobes, entraînant la perte de la structure tridimensionnelle native $\to$ inactivation enzymatique irréversible et mort cellulaire des formes végétatives. Chasse aussi l'$\ce{O2}$ intercellulaire (limite oxydation).
\end{itemize}
\item \textbf{Stérilisation bocaux (\SI{115}{\degreeCelsius}, \SI{20}{\minute}) :}
\begin{itemize}
\item \textbf{Facteur visé :} \textbf{Micro-organismes (végétatives + Spores)}.
\item \textbf{Mécanisme :} \textbf{Destruction thermique des spores} par chaleur humide à \SI{115}{\degreeCelsius} (pression $\approx$ \SI{0,7}{\bar} en autoclave). La température $> \SI{100}{\degreeCelsius}$ permet la coagulation/dénaturation irréversible des protéines du cortex de la spore et la destruction de son matériel génétique (ADN), assurant la stérilité commerciale. (Note : pour tomate pH $\approx$ 4,2--4,5, limite botulisme, \SI{115}{\degreeCelsius} est standard sécurité).
\end{itemize}
\end{enumerate}

\item \textbf{Rôle de la concentration (cuisson \SI{30}{\minute}) sur la conservation (hors chaleur) :}
La concentration par évaporation de l'eau \textbf{abaisse l'activité de l'eau ($a_w$)} de la pâte. En passant de $\approx$ \SI{94}{\%} eau (tomate fraîche) à \SI{72}{\%} eau (\SI{28}{\%} MS), l'$a_w$ chute de $\approx 0,99$ à $\approx 0,92\text{--}0,94$. Bien que cette $a_w$ ne soit pas inhibatrice seule pour tous les microbes, elle constitue une \textbf{barrière supplémentaire (hurdle)} qui :
\begin{itemize}
\item Ralentit/empêche la croissance des bactéries résiduelles (exigences $a_w > 0,90$).
\item Diminue la vitesse des réactions chimiques d'altération (oxydation, Maillard, dégradation vitamines) qui nécessitent l'eau comme solvant/réactif.
\item Augmente la concentration en solutés (sucres, acides, sels) $\to$ effet osmotique additionnel.
\end{itemize}
C'est l'association \emph{Chaleur (cuisson) + Baisse $a_w$ (concentration) + Acidité (pH tomate) + Stérilisation finale} qui garantit la stabilité.

\item \textbf{Intérêts de la transformation locale (région de l'Ouest) :}
\begin{enumerate}
\item \textbf{Économique 1 : Valorisation du surplus saisonnier et réduction des pertes post-récolte.} En saison de pointe (déc--mars), le prix s'effondre (\SI{150}{\FCFA}/\si{\kg}) par surproduction et pertes (pourriture, transport). La transformation en pâte (produit stable, volume réduit 6:1) permet d'écouler la production excédentaire, générant un revenu (gain \SI{600}{\FCFA}/\si{\kg} pâte) au lieu de pertes totales.
\item \textbf{Économique 2 : Lissage de l'offre et des prix / Création de valeur ajoutée locale.} La pâte stockée est vendue en saison sèche (prix tomate fraîche \SI{1200}{\FCFA}/\si{\kg}). Cela stabilise l'approvisionnement du marché local, réduit la dépendance aux importations (concentré de tomate importé), crée des emplois (récolte, transformation, conditionnement, vente) et fixe la richesse dans la région.
\item \textbf{Nutritionnel : Disponibilité toute l'année en \cle{lycopène} et micronutriments.} La tomate est la principale source alimentaire de \textbf{lycopène} (caroténoïde antioxydant puissant, prévention cancers/cardiovasculaires). La cuisson/concentration \emph{augmente} la biodisponibilité du lycopène (rupture matrice cellulaire, isomérisation cis-trans). La pâte de tomate locale permet aux populations de bénéficier de cet apport protecteur hors saison de production fraîche, améliorant la sécurité nutritionnelle.
\end{enumerate}
\end{enumerate}

\vspace{0,5cm}
\noindent\textbf{Barème indicatif (sur 20) :}
\begin{itemize}
\item Q1 : Calcul gain correct (\SI{600}{\FCFA}/\si{\kg} pâte) avec étapes : 3 pts.
\item Q2 : Logigramme complet, ordonné, flèches : 3 pts.
\item Q3 : Calcul FR toutes étapes + Identification 2 meilleures (Stérilisation, Ébouillantage) : 3 pts.
\item Q4 : Facteurs + Mécanismes précis (Dénaturation protéines enzymes/végétatives ; Destruction spores \SI{115}{\degreeCelsius}) : 4 pts.
\item Q5 : Explication rôle concentration sur $a_w$ (effet hurdle) : 3 pts.
\item Q6 : 2 intérêts économiques contextuels + 1 intérêt nutritionnel (lycopène/biodisponibilité) : 4 pts.
\end{itemize}
\textbf{Points de vigilance :} Unités FCFA/kg. Logigramme : étapes dans l'ordre du tableau. FR = ratio témoin/traité (pas différence). Mécanisme ébouillantage = enzymes (principal) + flore végétative. Mécanisme stérilisation = spores (température $> \SI{100}{\degreeCelsius}$). Concentration = baisse $a_w$ (pas seulement « évaporation eau »). Lycopène = mot-clé nutritionnel attendu pour tomate transformée.
\end{corrige}

\newpage

\coursec{Fiche bilan}

\begin{popfigure}
\begin{center}
\begin{center}\tcbox[colback=popblue,colframe=popblue,coltext=white,arc=9pt,boxsep=4pt,left=6pt,right=6pt]{\bfseries\small Conservation des fruits de saison}\end{center}
\vspace{-2pt}
\begin{tcbraster}[raster columns=3,raster equal height=rows,raster column skip=6pt,raster row skip=6pt,enhanced,blankest]
\begin{tcolorbox}[enhanced,colback=white,colframe=poporange,boxrule=0.9pt,arc=3pt,title={Enjeux},fonttitle=\bfseries\scriptsize,coltitle=white,colbacktitle=poporange,left=4pt,right=4pt,top=2pt,bottom=2pt,boxsep=1pt,toptitle=1.5pt,bottomtitle=1.5pt,halign=left]
\begin{itemize}[nosep,leftmargin=*,itemsep=1pt,font=\scriptsize]
\item Abondance saisonnière
\item Pertes après récolte
\end{itemize}
\end{tcolorbox}
\begin{tcolorbox}[enhanced,colback=white,colframe=popgreen,boxrule=0.9pt,arc=3pt,title={Facteurs d'altération},fonttitle=\bfseries\scriptsize,coltitle=white,colbacktitle=popgreen,left=4pt,right=4pt,top=2pt,bottom=2pt,boxsep=1pt,toptitle=1.5pt,bottomtitle=1.5pt,halign=left]
\begin{itemize}[nosep,leftmargin=*,itemsep=1pt,font=\scriptsize]
\item Micro-organismes
\item Enzymes et oxydation
\item Eau disponible $a_w$
\end{itemize}
\end{tcolorbox}
\begin{tcolorbox}[enhanced,colback=white,colframe=poppurple,boxrule=0.9pt,arc=3pt,title={Principes},fonttitle=\bfseries\scriptsize,coltitle=white,colbacktitle=poppurple,left=4pt,right=4pt,top=2pt,bottom=2pt,boxsep=1pt,toptitle=1.5pt,bottomtitle=1.5pt,halign=left]
\begin{itemize}[nosep,leftmargin=*,itemsep=1pt,font=\scriptsize]
\item Chaleur : stérilisation
\item Sucre, acide : osmose
\item Séchage : baisse de $a_w$
\end{itemize}
\end{tcolorbox}
\begin{tcolorbox}[enhanced,colback=white,colframe=poppink,boxrule=0.9pt,arc=3pt,title={Techniques},fonttitle=\bfseries\scriptsize,coltitle=white,colbacktitle=poppink,left=4pt,right=4pt,top=2pt,bottom=2pt,boxsep=1pt,toptitle=1.5pt,bottomtitle=1.5pt,halign=left]
\begin{itemize}[nosep,leftmargin=*,itemsep=1pt,font=\scriptsize]
\item Séchage (mangue)
\item Conserve, jus
\item Confiture, pâte (tomate)
\end{itemize}
\end{tcolorbox}
\begin{tcolorbox}[enhanced,colback=white,colframe=popgold,boxrule=0.9pt,arc=3pt,title={Comparaison},fonttitle=\bfseries\scriptsize,coltitle=white,colbacktitle=popgold,left=4pt,right=4pt,top=2pt,bottom=2pt,boxsep=1pt,toptitle=1.5pt,bottomtitle=1.5pt,halign=left]
\begin{itemize}[nosep,leftmargin=*,itemsep=1pt,font=\scriptsize]
\item Durée de conservation
\item Qualité nutritionnelle
\item Coût
\end{itemize}
\end{tcolorbox}
\begin{tcolorbox}[enhanced,colback=white,colframe=popmuted,boxrule=0.9pt,arc=3pt,title={Intérêts},fonttitle=\bfseries\scriptsize,coltitle=white,colbacktitle=popmuted,left=4pt,right=4pt,top=2pt,bottom=2pt,boxsep=1pt,toptitle=1.5pt,bottomtitle=1.5pt,halign=left]
\begin{itemize}[nosep,leftmargin=*,itemsep=1pt,font=\scriptsize]
\item Économique
\item Nutritionnel
\item Environnemental
\end{itemize}
\end{tcolorbox}
\end{tcbraster}

\end{center}
\legende{Carte mentale de la leçon : les six grandes idées autour de la conservation des fruits de saison (mangue, tomate).}
\end{popfigure}

\begin{bilan}
\textbf{Définitions clés}
\begin{itemize}
  \item \cle{Fruit de saison} : fruit produit en abondance sur une courte période de l'année (mangue : mars--juillet ; tomate : saison sèche), souvent perdu faute de conservation.
  \item \cle{Conservation} : ensemble des techniques visant à empêcher ou ralentir l'altération d'un aliment afin de prolonger sa durée de consommation.
  \item \cle{Transformation} : modification de l'état du fruit (jus, confiture, pâte, séché) qui le stabilise et lui donne une valeur ajoutée.
  \item \cle{Stérilisation} : destruction des micro-organismes et de leurs spores par la chaleur (environ \SI{100}{\celsius} à \SI{120}{\celsius}) en récipient hermétique.
  \item \cle{Pasteurisation} : chauffage modéré (environ \SI{65}{\celsius} à \SI{85}{\celsius}) qui détruit les germes pathogènes sans tout stériliser.
  \item \cle{Activité de l'eau} ($a_w$) : fraction d'eau disponible pour les micro-organismes ; en dessous de $a_w \approx 0{,}6$, aucun micro-organisme ne se développe.
\end{itemize}

\textbf{Facteurs d'altération et principes de lutte}
\begin{center}
\begin{tabularx}{\linewidth}{|l|X|X|}
\hline
\rowcolor{popblueL} \textbf{Facteur} & \textbf{Effet sur le fruit} & \textbf{Principe de conservation} \\
\hline
Micro-organismes (bactéries, moisissures, levures) & Pourriture, fermentation, toxines & Chaleur (stérilisation), acidité, sucre concentré, séchage, hygiène \\
\hline
Enzymes du fruit & Brunissement, ramollissement, perte d'arômes & Blanchiment (chauffage bref) qui dénature les enzymes \\
\hline
Oxydation (dioxygène) & Brunissement, destruction de la vitamine C & Récipient hermétique, jus de citron (acide ascorbique) \\
\hline
Eau disponible ($a_w$ élevée) & Milieu favorable aux germes & Séchage, concentration par le sucre ou l'évaporation \\
\hline
\end{tabularx}
\end{center}

\textbf{Techniques à connaître (mangue et tomate)}
\begin{itemize}
  \item \cle{Séchage} : déshydratation au soleil ou en séchoir ($a_w$ abaissée) ; mangues séchées, tomates séchées. Simple, peu coûteux, mais sensible à l'humidité et aux poussières.
  \item \cle{Mise en conserve / stérilisation} : conditionnement hermétique + chauffage ; conservation longue (1 à 2 ans) ; exige du matériel (bocaux, autoclave ou stérilisateur).
  \item \cle{Confiture} : cuisson du fruit avec environ \SI{1}{kg} de sucre pour \SI{1}{kg} de fruit ; le sucre concentre le milieu (osmose) et bloque les germes.
  \item \cle{Jus} : extraction, filtration, pasteurisation, embouteillage à chaud ; conservation moyenne, à consommer rapidement après ouverture.
  \item \cle{Pâte de tomate} : concentration par évaporation (double concentré : environ 28 à 30 \% d'extrait sec), puis stérilisation.
\end{itemize}

\textbf{Méthode express : comparer deux techniques}
\begin{enumerate}
  \item Choisir des critères : durée de conservation, coût, qualité nutritionnelle, simplicité.
  \item Appliquer chaque technique au même fruit, dans les mêmes conditions.
  \item Observer à intervalles réguliers (aspect, odeur, moisissures).
  \item Conclure par un tableau comparatif argumenté.
\end{enumerate}
\end{bilan}

\begin{aretenir}[Checklist avant le Bac]
\begin{itemize}
  \item $\square$ Je sais définir conservation, transformation, stérilisation, pasteurisation et activité de l'eau ($a_w$).
  \item $\square$ Je cite les trois grands facteurs d'altération : micro-organismes, enzymes, oxydation.
  \item $\square$ J'explique pourquoi la chaleur conserve : destruction des germes et dénaturation des enzymes.
  \item $\square$ J'explique le rôle du sucre dans la confiture : milieu concentré, eau indisponible par osmose.
  \item $\square$ J'explique le principe du séchage : abaissement de $a_w$ sous le seuil de développement microbien ($a_w \approx 0{,}6$).
  \item $\square$ Je décris les étapes d'une mise en conserve : lavage, préparation, conditionnement hermétique, stérilisation, refroidissement, étiquetage.
  \item $\square$ Je connais les proportions d'une confiture : environ \SI{1}{kg} de sucre pour \SI{1}{kg} de fruit.
  \item $\square$ Je sais que le blanchiment (chauffage bref) inactive les enzymes avant séchage ou congélation.
  \item $\square$ Je compare les techniques selon durée de conservation, coût et qualité nutritionnelle.
  \item $\square$ Je cite les intérêts : réduction des pertes, revenus hors saison, vitamines toute l'année, moins de déchets.
  \item $\square$ Je sais donner des exemples camerounais : mangues séchées du Nord, pâte de tomate artisanale, jus de mangue pasteurisé.
  \item $\square$ Je rédige une démarche d'investigation complète : observation, hypothèse, expérience, interprétation, conclusion.
\end{itemize}
\end{aretenir}

\findecours

\end{document}
